% Template:     Informe/Reporte LaTeX
% Documento:    Archivo principal
% Versión:      6.5.5 (02/09/2019)
% Codificación: UTF-8
%
% Autor: Pablo Pizarro R. @ppizarror
%        Facultad de Ciencias Físicas y Matemáticas
%        Universidad de Chile
%        pablo.pizarro@ing.uchile.cl, ppizarror.com
%
% Manual template: [https://latex.ppizarror.com/Template-Informe/]
% Licencia MIT:    [https://opensource.org/licenses/MIT/]

% CREACIÓN DEL DOCUMENTO
\documentclass[letterpaper,11pt]{article} % Articulo tamaño carta, 11pt
\usepackage[utf8]{inputenc} % Codificación UTF-8

% INFORMACIÓN DEL DOCUMENTO
\def\titulodelinforme {Título del informe}
\def\temaatratar {Tema a tratar}

\def\autordeldocumento {Nombre del autor}
\def\nombredelcurso {Curso}
\def\codigodelcurso {CO-1234}

\def\nombreuniversidad {Universidad de Chile}
\def\nombrefacultad {Facultad de Ciencias Físicas y Matemáticas}
\def\departamentouniversidad {Departamento de la Universidad}
\def\imagendepartamento {departamentos/fcfm}
\def\imagendepartamentoescala {0.2}
\def\localizacionuniversidad {Santiago, Chile}

% INTEGRANTES, PROFESORES Y FECHAS
\def\tablaintegrantes {
\begin{tabular}{ll}
	Integrantes:
	& \begin{tabular}[t]{l}
		Integrante 1 \\
		Integrante 2
	\end{tabular} \\
	Profesor:
	& \begin{tabular}[t]{l}
		Profesor 1
	\end{tabular} \\
	Auxiliar:
	& \begin{tabular}[t]{l}
		Auxiliar 1
	\end{tabular} \\
	Ayudantes:
	& \begin{tabular}[t]{l}
		Ayudante 1 \\
		Ayudante 2
	\end{tabular} \\
	\multicolumn{2}{l}{Ayudante de laboratorio: Ayudante 1} \\
	& \\
	\multicolumn{2}{l}{Fecha de realización: \today} \\
	\multicolumn{2}{l}{Fecha de entrega: \today} \\
	\multicolumn{2}{l}{\localizacionuniversidad}
\end{tabular}}{
}

% CONFIGURACIONES
% Template:     Informe/Reporte LaTeX
% Documento:    Configuraciones del template
% Versión:      4.2.7 (09/07/2017)
% Codificación: UTF-8
%
% Autor: Pablo Pizarro R.
%        Facultad de Ciencias Físicas y Matemáticas
%        Universidad de Chile
%        pablo.pizarro@ing.uchile.cl, ppizarror.com
%
% Manual template: [http://ppizarror.com/Template-Informe/]
% Licencia MIT:    [https://opensource.org/licenses/MIT/]

% CONFIGURACIONES GENERALES
\def\addemptypagetwosides {false}  % Añade pags. en blanco al imprimir a 2 caras
\def\anumsecaddtocounter {false}   % Función para insertar títulos anum aumenta contador
\def\defaultimagefolder {images/}  % Carpeta raíz de las imágenes
\def\defaultinterline {1.0}        % Interlineado por defecto [pt]
\def\defaultnewlinesize {11.0}     % Tamaño del salto de línea [pt]
\def\fontdocument {lmodern}        % Tipografía (lmodern,arial,helvet)
\def\numberedequation {true}       % Ecuaciones con \insert... numeradas
\def\pointdecimal {false}          % Decimales con punto en vez de coma
\def\romanpageuppercase {true}     % Páginas en número romano en mayúsculas
\def\showdotontitles {true}        % Punto al final de cada número de título/subtítulo
\def\tablepadding {1.0}            % Ancho de celda de las tablas

% CONFIGURACIÓN DE LAS LEYENDAS - CAPTION
\def\captionlessmarginimage {0.1}  % Margen sup/inf de fig si no hay leyenda [cm]
\def\captionlrmargin {2.0}         % Márgenes izq/der de la leyenda [cm]
\def\captionalignment {justified}  % Alineación leyenda: justified,centered,left,right
\def\captiontbmarginfigure {9.35}  % Margen sup/inf de la leyenda en figuras [pt]
\def\captiontbmargintable {7.0}    % Margen sup/inf de la leyenda en tablas [pt]
\def\captiontextbold {false}       % Etiquetas (Código,Figura,Tabla) en negrita
\def\codecaptiontop {true}         % Leyenda arriba del código fuente
\def\figurecaptiontop {false}      % Leyenda arriba de las imágenes
\def\showsectiononcaption {false}  % Muestra el número de sección en las leyendas
\def\tablecaptiontop {true}        % Leyenda arriba de las tablas

% CONFIGURACIÓN DEL ÍNDICE
\def\indexdepth {3}                % Profundidad máxima del índice
\def\indextitlemargin {7.0}        % Margen título en índice \insertindextitle [pt]
\def\showindex {true}              % Muestra el índice
\def\showindexofcode {true}        % Muestra la lista de códigos fuente
\def\showindexofcontents {true}    % Muestra la lista de contenidos
\def\showindexoffigures {true}     % Muestra la lista de figuras
\def\showindexoftables {true}      % Muestra la lista de tablas

% CONFIGURACIÓN DE LOS COLORES DEL DOCUMENTO
\def\captioncolor {black}          % Color de la etiqueta (Código,Figura,Tabla)
\def\captiontextcolor {black}      % Color de la leyenda
\def\citecolor {black}             % Color del número de las referencias o citas
\def\highlightcolor {yellow}       % Color del subrayado con \hl
\def\indextitlecolor {black}       % Color de los títulos del índice
\def\linkcolor {black}             % Color de los links del doc.
\def\maintextcolor {black}         % Color principal del texto
\def\portraittitlecolor {black}    % Color de los títulos de la portada
\def\showborderonlinks {false}     % Color de un links por un recuadro de color
\def\subsubtitlecolor {black}      % Color de los sub-subtítulos
\def\subtitlecolor {black}         % Color de los subtítulos
\def\tablelinecolor {black}        % Color de las líneas de las tablas
\def\titlecolor {black}            % Color de los títulos
\def\urlcolor {magenta}            % Color de los enlaces web (\url,\href)

% MÁRGENES DE FIGURAS
\def\marginbottomimages {-0.2}     % Margen inferior figura [cm]
\def\marginfloatimages {-13.0}     % Margen sup. fig. float \insertimageleft/right [pt]
\def\margintopimages {0.0}         % Margen superior figura [cm]

% ANEXO, CITAS, REFERENCIAS
\def\apaciterefsep {9}             % Separación entre referencias [pt] {apacite}
\def\bibtexrefsep {9}              % Separación entre referencias [pt] {bibtex}
\def\natbibrefsep {5}              % Separación entre referencias [pt] {natbib}
\def\referencenumsection {false}   % Sección de referencias numerada
\def\sectionappendixlastchar {.}   % Caracter entre el n° de la sec. de anexo y el título
\def\showappendixsecindex {false}  % Muestra el título de la sec. de anexos en el índice
\def\showappendixsectitle {false}  % Muestra el título de la sec. de anexo en el informe
\def\stylecitereferences {bibtex}  % Estilo citas y referencias (bibtex,apacite,natbib)
\def\twocolumnreferences {false}   % Referencias en dos columnas

% CONFIGURACIÓN PORTADA Y HEADERS
\def\gradecodeonportrait {false}   % Muestra el código del curso en la portada
\def\showfooter {true}             % Muestra el footer
\def\showheadertitle {true}        % Muestra título de la sección en el header

% MÁRGENES DE PÁGINA
\def\firstpagemargintop {3.8}      % Margen superior página portada [cm]
\def\pagemarginbottom {2.7}        % Margen inferior página [cm]
\def\pagemarginleft {2.54}         % Margen izquierdo página [cm]
\def\pagemarginright {2.54}        % Margen derecho página [cm]
\def\pagemargintop {3.0}           % Margen superior página [cm]

% ESTILO Y TAMAÑO DE TÍTULOS
\def\fontsizesubsubtitle {\large}  % Tamaño sub-subtítulos
\def\fontsizesubtitle {\Large}     % Tamaño subtítulos
\def\fontsizetitle {\huge}         % Tamaño títulos
\def\fontsizetitlei {\huge}        % Tamaño títulos en el índice
\def\stylesubsubtitle {\bfseries}  % Estilo sub-subtítulos
\def\stylesubtitle {\bfseries}     % Estilo subtítulos
\def\styletitle {\bfseries}        % Estilo títulos
\def\styletitlei {\bfseries}       % Estilo títulos en el índice

% OPCIONES DEL PDF COMPILADO
\def\cfgbookmarksopenlevel {1}     % Nivel de los marcadores a mostrar (1: secciones)
\def\cfgpdfbookmarkopen {true}     % Abre el panel de marcadores al abrir el pdf
\def\cfgpdfcenterwindow {true}     % Centra la ventana del lector al abrir el pdf
\def\cfgpdfdisplaydoctitle {true}  % Muestra el título del informe como título del pdf
\def\cfgpdffitwindow {false}       % Ajusta la ventana del lector al tamaño del pdf
\def\cfgpdftoolbar {true}          % Muestra la barra de herramientas del lector pdf

% NOMBRE DE OBJETOS
\def\nameabstract {Resumen}           % Nombre del resumen-abstract
\def\nameappendixsection{Anexos}      % Nombre de la sección de anexos/apéndices
\def\nameportraitpage {Portada}       % Etiqueta página de la portada
\def\namereferences {Referencias}     % Nombre de la sección de referencias
\def\nomltappendixsection {Anexo}     % Etiqueta sección en anexo/apéndices
\def\nomltcont {Índice de Contenidos} % Nombre del índice de contenidos
\def\nomltfigure {Lista de Figuras}   % Nombre del índice de la lista de figuras
\def\nomltsrc {Lista de Códigos}      % Nombre del índice de la lista de código
\def\nomlttable {Lista de Tablas}     % Nombre del índice de la lista de tablas
\def\nomltwfigure {Figura}            % Etiqueta leyenda de las figuras
\def\nomltwsrc {Código}               % Etiqueta leyenda del código fuente
\def\nomltwtable {Tabla}              % Etiqueta leyenda de las tablas


% IMPORTACIÓN DE LIBRERÍAS
% Template:     Informe/Reporte LaTeX
% Documento:    Importación de librerías
% Versión:      DEV
% Codificación: UTF-8
%
% Autor: Pablo Pizarro R.
%        Facultad de Ciencias Físicas y Matemáticas
%        Universidad de Chile
%        pablo@ppizarror.com
%
% Manual template: [https://latex.ppizarror.com/informe]
% Licencia MIT:    [https://opensource.org/licenses/MIT]

% -----------------------------------------------------------------------------
% LIBRERÍAS IMPORTANTES
% -----------------------------------------------------------------------------

% Define el idioma
\usepackage[spanish,es-nosectiondot,es-lcroman,es-noquoting]{babel}

% Manejo de condicionales
\usepackage{ifthen}

% Lanza un mensaje de error indicando mala configuración
%	#1	Parámetros opcionales (nostop,noheader)
%	#2	Mensaje de error
% 	#3	Configuración usada
%	#4	Valores esperados
\newcommand{\throwbadconfig}[4][]{
	\ifthenelse{\equal{#1}{noheader}}{
		\errmessage{LaTeX Warning: #4}
	}{
		\ifthenelse{\equal{#1}{noheader-nostop}}{
			\errmessage{LaTeX Warning: #4}
		}{
			\errmessage{LaTeX Warning: #2 \noexpand #3=#3. Valores esperados: #4}
		}
	}
	\ifthenelse{\equal{#1}{nostop}}{}{
		\ifthenelse{\equal{#1}{noheader-nostop}}{}{
			\stop
		}
	}
}

% Librerías matemáticas
\ifthenelse{\equal{\equationleftalign}{true}}{
	\usepackage[fleqn]{amsmath}
}{
	\usepackage{amsmath}
}

% -----------------------------------------------------------------------------
% PARCHES DE LIBRERÍAS
% -----------------------------------------------------------------------------
\let\counterwithout\relax
\let\counterwithin\relax

% -----------------------------------------------------------------------------
% SE GUARDAN VARIABLES ANTES DE CARGAR LIBRERÍAS
% -----------------------------------------------------------------------------
\let\RE\Re
\let\IM\Im

% -----------------------------------------------------------------------------
% LIBRERÍAS INDEPENDIENTES
% -----------------------------------------------------------------------------
\usepackage{amssymb}       % Librerías matemáticas
\usepackage{amsthm}        % Definición de teoremas
\usepackage{array}         % Nuevas características a las tablas
\usepackage{bigstrut}      % Líneas horizontales en tablas
\usepackage{bm}            % Caracteres en negrita en ecuaciones
\usepackage{booktabs}      % Permite manejar elementos visuales en tablas
\usepackage{caption}       % Leyendas
\usepackage{changepage}    % Condicionales para administrar páginas
\usepackage{chngcntr}      % Añade números a las leyendas
\usepackage{color}         % Colores
\usepackage{datetime}      % Fechas
\usepackage{floatpag}      % Maneja números de páginas
\usepackage{floatrow}      % Permite administrar posiciones en los caption
\usepackage{framed}        % Permite creación de recuadros
\usepackage{gensymb}       % Simbología común
\usepackage{geometry}      % Dimensiones y geometría del documento
\usepackage{graphicx}      % Propiedades extra para los gráficos
\usepackage{lipsum}        % Permite crear párrafos de prueba
\usepackage{listings}      % Permite añadir código fuente
\usepackage{listingsutf8}  % Acepta codificación UTF-8 en código fuente
\usepackage{longtable}     % Permite utilizar tablas en varias hojas
\usepackage{mathtools}     % Permite utilizar notaciones matemáticas
\usepackage{multicol}      % Múltiples columnas
\usepackage{needspace}     % Maneja los espacios en página
\usepackage{pdflscape}     % Modo página horizontal de página
\usepackage{pdfpages}      % Permite administrar páginas en pdf
\usepackage{physics}       % Paquete de matemáticas
\usepackage{rotating}      % Permite rotación de objetos
\usepackage{sectsty}       % Cambia el estilo de los títulos
\usepackage{selinput}      % Compatibilidad con acentos
\usepackage{setspace}      % Cambia el espacio entre líneas
\usepackage{soul}          % Permite subrayar texto
\usepackage{subfig}        % Permite agrupar imágenes
\usepackage{textcomp}      % Simbología común
\usepackage{url}           % Permite añadir enlaces
\usepackage{wasysym}       % Contiene caracteres misceláneos
\usepackage{wrapfig}       % Posición de imágenes
\usepackage{xspace}        % Adminsitra espacios en párrafos y líneas

% En 6.3.7 se desactiva cellspace
% \usepackage{cellspace}     % Tamaños en celdas de tablas

% \usepackage{colortbl}, xcolor al cargar table ya importa colortbl (v6.3.3)
% El paquete colortbl queda obsoleto con array+xcolor si la versión es inferior
% a 1.0c, se recomienda actualizar o ocurren errores asociados a @parbox

% Se prefiere no utilizar matlab-prettifier ya que es muy pesada para lo que
% la librería ofrece
% \usepackage{matlab-prettifier}

% -----------------------------------------------------------------------------
% LIBRERÍAS CON PARÁMETROS
% -----------------------------------------------------------------------------
\usepackage[makeroom]{cancel} % Cancelar términos en fórmulas
\usepackage[inline]{enumitem} % Permite enumerar ítems
\usepackage[subfigure,titles]{tocloft} % Maneja entradas en el índice
\usepackage[figure,table,lstlisting]{totalcount} % Contador de objetos
\usepackage[normalem]{ulem} % Permite tachar y subrayar
\usepackage[dvipsnames,table,usenames]{xcolor} % Paquete de colores avanzado

% -----------------------------------------------------------------------------
% LIBRERÍAS CONDICIONALES
% -----------------------------------------------------------------------------
\ifthenelse{\equal{\footnotepagetoprule}{true}}{ % Regla superior
	\usepackage[bottom,hang]{footmisc} % Estilo pie de página
}{
	\usepackage[bottom,norule,hang]{footmisc}
}
\ifthenelse{\equal{\showdotaftersnum}{true}}{ % Agrega punto a títulos/subtítulos
	\usepackage{secdot}
	\sectiondot{subsection}
	\sectiondot{subsubsection}}{
}
% Desde 6.2.8 se debe cargar al final para evitar errores:
%	- Option clash for package hyperref
% 	- name has been referenced but does not exist, replaced by a fixed one
% Desde 6.5.6 se carga después de las referencias (apacite, natbib, bibtex)
\usepackage[pdfencoding=auto,psdextra]{hyperref} % Enlaces, referencias
\ifthenelse{\equal{\stylecitereferences}{natbib}}{ % Formato citas natbib
	\ifthenelse{\equal{\natbibrefstyle}{apa}}{
		\ifthenelse{\equal{\natbibsquare}{true}}{
			\usepackage[square]{natbib}
		}{
			\usepackage[round]{natbib}
		}
	}{
	\ifthenelse{\equal{\natbibrefstyle}{ieeetr}}{
		\ifthenelse{\equal{\natbibsquare}{true}}{
			\usepackage[square,numbers]{natbib}
		}{
			\usepackage[round,numbers]{natbib}
		}
	}{
	\ifthenelse{\equal{\natbibrefstyle}{unsrt}}{
		\ifthenelse{\equal{\natbibsquare}{true}}{
			\usepackage[square,numbers]{natbib}
		}{
			\usepackage[round,numbers]{natbib}
		}
	}{
	\ifthenelse{\equal{\natbibrefstyle}{abbrvnat}}{
		\ifthenelse{\equal{\natbibsquare}{true}}{
			\usepackage[square,numbers]{natbib}
		}{
			\usepackage[round,numbers]{natbib}
		}
	}{
		\ifthenelse{\equal{\natbibnumbers}{true}}{
			\ifthenelse{\equal{\natbibsquare}{true}}{
				\usepackage[square,numbers]{natbib}
			}{
				\usepackage[round,numbers]{natbib}
			}
		}{
			\ifthenelse{\equal{\natbibsquare}{true}}{
				\usepackage[square]{natbib}
			}{
				\usepackage[round]{natbib}
			}
		}
	}}}}
	\usepackage[nottoc,notlof,notlot]{tocbibind}
}{
	\ifthenelse{\equal{\stylecitereferences}{apacite}}{ % Formato citas apacite
		\usepackage{apacite}
		\usepackage[nottoc,notlof,notlot]{tocbibind}
	}{
		\ifthenelse{\equal{\stylecitereferences}{bibtex}}{ % Formato citas bibtex
		}{}
	}
}
\ifthenelse{\equal{\showappendixsecindex}{true}}{ % Anexos/Apéndices
	\usepackage[toc]{appendix}
}{
	\usepackage{appendix}
}
\ifthenelse{\equal{\importtikz}{true}}{ % Importa la librería tikz
	\usepackage{tikz}}{
}

% -----------------------------------------------------------------------------
% ESTILO PORTADA
% -----------------------------------------------------------------------------
\ifthenelse{\equal{\hfstyle}{style11}}{
	\usepackage{lastpage}}{}
\ifthenelse{\equal{\hfstyle}{style12}}{
	\usepackage{lastpage}}{}
\ifthenelse{\equal{\hfstyle}{style13}}{
	\usepackage{lastpage}}{}
\ifthenelse{\equal{\hfstyle}{style14}}{
	\usepackage{lastpage}}{
}

% -----------------------------------------------------------------------------
% LIBRERÍAS DEPENDIENTES
% -----------------------------------------------------------------------------
\usepackage{bookmark}      % Administración de marcadores en pdf
\usepackage{fancyhdr}      % Encabezados y pie de páginas
\usepackage{float}         % Administrador de posiciones de objetos
\usepackage{hyperxmp}      % Etiquetas opcionales para el pdf compilado
\usepackage{multirow}      % Agrega nuevas opciones a las tablas
\usepackage{notoccite}     % Desactiva las citas en el índice
\usepackage{titlesec}      % Administración de títulos

% -----------------------------------------------------------------------------
% TIPOGRAFÍA DEL DOCUMENTO
% -----------------------------------------------------------------------------

% Tipografías clásicas
\ifthenelse{\equal{\fontdocument}{lmodern}}{ % Default
	\usepackage{lmodern}
}{
\ifthenelse{\equal{\fontdocument}{arial}}{
	\usepackage{helvet}
	\renewcommand{\familydefault}{\sfdefault}
}{
\ifthenelse{\equal{\fontdocument}{arial2}}{
	\usepackage{arial}
}{
\ifthenelse{\equal{\fontdocument}{times}}{
	\usepackage{mathptmx}
}{
\ifthenelse{\equal{\fontdocument}{helvet}}{
	\usepackage{helvet}
}{
\ifthenelse{\equal{\fontdocument}{alegreyasans}}{
	\usepackage[sfdefault]{AlegreyaSans}
	\renewcommand*\oldstylenums[1]{{\AlegreyaSansOsF #1}}
}{
\ifthenelse{\equal{\fontdocument}{mathpazo}}{
	\usepackage{mathpazo}
}{

% Otros
\ifthenelse{\equal{\fontdocument}{accantis}}{
	\usepackage{accanthis}
}{
\ifthenelse{\equal{\fontdocument}{alegreya}}{
	\usepackage{Alegreya}
	\renewcommand*\oldstylenums[1]{{\AlegreyaOsF #1}}
}{
\ifthenelse{\equal{\fontdocument}{algolrevived}}{
	\usepackage{algolrevived}
}{
\ifthenelse{\equal{\fontdocument}{antiqua}}{
	\usepackage{antiqua}
}{
\ifthenelse{\equal{\fontdocument}{antpolt}}{
	\usepackage{antpolt}
}{
\ifthenelse{\equal{\fontdocument}{antpoltlight}}{
	\usepackage[light]{antpolt}
}{
\ifthenelse{\equal{\fontdocument}{anttor}}{
	\usepackage[math]{anttor}
}{
\ifthenelse{\equal{\fontdocument}{anttorcondensed}}{
	\usepackage[condensed,math]{anttor}
}{
\ifthenelse{\equal{\fontdocument}{anttorlight}}{
	\usepackage[light,math]{anttor}
}{
\ifthenelse{\equal{\fontdocument}{anttorlightcondensed}}{
	\usepackage[light,condensed,math]{anttor}
}{
\ifthenelse{\equal{\fontdocument}{arev}}{
	\usepackage{arev}
}{
\ifthenelse{\equal{\fontdocument}{arimo}}{
	\usepackage[sfdefault]{arimo}
	\renewcommand*\familydefault{\sfdefault}
}{
\ifthenelse{\equal{\fontdocument}{aurical}}{
	\usepackage{aurical}
}{
\ifthenelse{\equal{\fontdocument}{avant}}{
	\usepackage{avant}
}{
\ifthenelse{\equal{\fontdocument}{baskervald}}{
	\usepackage{baskervald}
}{
\ifthenelse{\equal{\fontdocument}{berasans}}{
	\usepackage[scaled]{berasans}
	\renewcommand*\familydefault{\sfdefault}
}{
\ifthenelse{\equal{\fontdocument}{beraserif}}{
	\usepackage{bera}
}{
\ifthenelse{\equal{\fontdocument}{biolinum}}{
	\usepackage{libertine}
	\renewcommand*\familydefault{\sfdefault}
}{
\ifthenelse{\equal{\fontdocument}{cabin}}{
	\usepackage[sfdefault]{cabin}
	\renewcommand*\familydefault{\sfdefault}
}{
\ifthenelse{\equal{\fontdocument}{cabincondensed}}{
	\usepackage[sfdefault,condensed]{cabin}
	\renewcommand*\familydefault{\sfdefault}
}{
\ifthenelse{\equal{\fontdocument}{cantarell}}{
	\usepackage[default]{cantarell}
}{
\ifthenelse{\equal{\fontdocument}{caladea}}{
	\usepackage{caladea}
}{
\ifthenelse{\equal{\fontdocument}{carlito}}{
	\usepackage[sfdefault]{carlito}
	\renewcommand*\familydefault{\sfdefault}
}{
\ifthenelse{\equal{\fontdocument}{chivolight}}{
	\usepackage[familydefault,light]{Chivo}
}{
\ifthenelse{\equal{\fontdocument}{chivoregular}}{
	\usepackage[familydefault,regular]{Chivo}
}{
\ifthenelse{\equal{\fontdocument}{clearsans}}{
	\usepackage[sfdefault]{ClearSans}
	\renewcommand*\familydefault{\sfdefault}
}{
\ifthenelse{\equal{\fontdocument}{comfortaa}}{
	\usepackage[default]{comfortaa}
}{
\ifthenelse{\equal{\fontdocument}{comicneue}}{
	\usepackage[default]{comicneue}
}{
\ifthenelse{\equal{\fontdocument}{comicneueangular}}{
	\usepackage[default,angular]{comicneue}
}{
\ifthenelse{\equal{\fontdocument}{crimson}}{
	\usepackage{crimson}
}{
\ifthenelse{\equal{\fontdocument}{cyklop}}{
	\usepackage{cyklop}
}{
\ifthenelse{\equal{\fontdocument}{dejavusans}}{
	\usepackage{DejaVuSans}
	\renewcommand*\familydefault{\sfdefault}
}{
\ifthenelse{\equal{\fontdocument}{dejavusanscondensed}}{
	\usepackage{DejaVuSansCondensed}
	\renewcommand*\familydefault{\sfdefault}
}{
\ifthenelse{\equal{\fontdocument}{droidsans}}{
	\usepackage[defaultsans]{droidsans}
	\renewcommand*\familydefault{\sfdefault}
}{
\ifthenelse{\equal{\fontdocument}{fetamont}}{ % Falta implementar aún
	\usepackage{fetamont}
	\renewcommand*\familydefault{\sfdefault}
}{
\ifthenelse{\equal{\fontdocument}{firasans}}{
	\usepackage[sfdefault]{FiraSans}
	\renewcommand*\familydefault{\sfdefault}
}{
\ifthenelse{\equal{\fontdocument}{iwona}}{
	\usepackage[math]{iwona}
}{
\ifthenelse{\equal{\fontdocument}{iwonacondensed}}{
	\usepackage[math]{iwona}
}{
\ifthenelse{\equal{\fontdocument}{iwonalight}}{
	\usepackage[light,math]{iwona}
}{
\ifthenelse{\equal{\fontdocument}{iwonalightcondensed}}{
	\usepackage[light,condensed,math]{iwona}
}{
\ifthenelse{\equal{\fontdocument}{kurier}}{
	\usepackage[math]{kurier}
}{
\ifthenelse{\equal{\fontdocument}{kuriercondensed}}{
	\usepackage[condensed,math]{kurier}
}{
\ifthenelse{\equal{\fontdocument}{kurierlight}}{
	\usepackage[light,math]{kurier}
}{
\ifthenelse{\equal{\fontdocument}{kurierlightcondensed}}{
	\usepackage[light,condensed,math]{kurier}
}{
\ifthenelse{\equal{\fontdocument}{lato}}{
	\usepackage[default]{lato}
}{
\ifthenelse{\equal{\fontdocument}{libris}}{
	\usepackage{libris}
	\renewcommand*\familydefault{\sfdefault}
}{
\ifthenelse{\equal{\fontdocument}{lxfonts}}{
	\usepackage{lxfonts}
}{
\ifthenelse{\equal{\fontdocument}{merriweather}}{
	\usepackage[sfdefault]{merriweather}
}{
\ifthenelse{\equal{\fontdocument}{merriweatherlight}}{
	\usepackage[sfdefault,light]{merriweather}
}{
\ifthenelse{\equal{\fontdocument}{mintspirit}}{
	\usepackage[default]{mintspirit}
}{
\ifthenelse{\equal{\fontdocument}{montserratalternatesextralight}}{
	\usepackage[defaultfam,extralight,tabular,lining,alternates]{montserrat}
	\renewcommand*\oldstylenums[1]{{\fontfamily{Montserrat-TOsF}\selectfont #1}}
}{
\ifthenelse{\equal{\fontdocument}{montserratalternatesregular}}{
	\usepackage[defaultfam,tabular,lining,alternates]{montserrat}
	\renewcommand*\oldstylenums[1]{{\fontfamily{Montserrat-TOsF}\selectfont #1}}
}{
\ifthenelse{\equal{\fontdocument}{montserratalternatesthin}}{
	\usepackage[defaultfam,thin,tabular,lining,alternates]{montserrat}
	\renewcommand*\oldstylenums[1]{{\fontfamily{Montserrat-TOsF}\selectfont #1}}
}{
\ifthenelse{\equal{\fontdocument}{montserratextralight}}{
	\usepackage[defaultfam,extralight,tabular,lining]{montserrat}
	\renewcommand*\oldstylenums[1]{{\fontfamily{Montserrat-TOsF}\selectfont #1}}
}{
\ifthenelse{\equal{\fontdocument}{montserratlight}}{
	\usepackage[defaultfam,light,tabular,lining]{montserrat}
	\renewcommand*\oldstylenums[1]{{\fontfamily{Montserrat-TOsF}\selectfont #1}}
}{
\ifthenelse{\equal{\fontdocument}{montserratregular}}{
	\usepackage[defaultfam,tabular,lining]{montserrat}
	\renewcommand*\oldstylenums[1]{{\fontfamily{Montserrat-TOsF}\selectfont #1}}
}{
\ifthenelse{\equal{\fontdocument}{montserratthin}}{
	\usepackage[defaultfam,thin,tabular,lining]{montserrat}
	\renewcommand*\oldstylenums[1]{{\fontfamily{Montserrat-TOsF}\selectfont #1}}
}{
\ifthenelse{\equal{\fontdocument}{nimbussans}}{
	\usepackage{nimbussans}
	\renewcommand*\familydefault{\sfdefault}
}{
\ifthenelse{\equal{\fontdocument}{noto}}{
	\usepackage[sfdefault]{noto}
	\renewcommand*\familydefault{\sfdefault}
}{
\ifthenelse{\equal{\fontdocument}{opensans}}{
	\usepackage[default,osfigures,scale=0.95]{opensans}
}{
\ifthenelse{\equal{\fontdocument}{overlock}}{
	\usepackage[sfdefault]{overlock}
	\renewcommand*\familydefault{\sfdefault}
}{
\ifthenelse{\equal{\fontdocument}{paratype}}{
	\usepackage{paratype}
	\renewcommand*\familydefault{\sfdefault}
}{
\ifthenelse{\equal{\fontdocument}{paratypesanscaption}}{
	\usepackage{PTSansCaption}
	\renewcommand*\familydefault{\sfdefault}
}{
\ifthenelse{\equal{\fontdocument}{paratypesansnarrow}}{
	\usepackage{PTSansNarrow}
	\renewcommand*\familydefault{\sfdefault}
}{
\ifthenelse{\equal{\fontdocument}{quattrocento}}{
	\usepackage[sfdefault]{quattrocento}
}{
\ifthenelse{\equal{\fontdocument}{raleway}}{
	\usepackage[default]{raleway}
}{
\ifthenelse{\equal{\fontdocument}{roboto}}{
	\usepackage[sfdefault]{roboto}
}{
\ifthenelse{\equal{\fontdocument}{robotocondensed}}{
	\usepackage[sfdefault,condensed]{roboto}
}{
\ifthenelse{\equal{\fontdocument}{robotolight}}{
	\usepackage[sfdefault,light]{roboto}
}{
\ifthenelse{\equal{\fontdocument}{robotolightcondensed}}{
	\usepackage[sfdefault,light,condensed]{roboto}
}{
\ifthenelse{\equal{\fontdocument}{robotothin}}{
	\usepackage[sfdefault,thin]{roboto}
}{
\ifthenelse{\equal{\fontdocument}{rosario}}{
	\usepackage[familydefault]{Rosario}
}{
\ifthenelse{\equal{\fontdocument}{sourcesanspro}}{
	\usepackage[default]{sourcesanspro}
}{
\ifthenelse{\equal{\fontdocument}{uarial}}{
	\usepackage{uarial}
	\renewcommand*\familydefault{\sfdefault}
}{
\ifthenelse{\equal{\fontdocument}{ugq}}{
	\renewcommand*\sfdefault{ugq}
	\renewcommand*\familydefault{\sfdefault}
}{
\ifthenelse{\equal{\fontdocument}{universalis}}{
	\usepackage[sfdefault]{universalis}
}{
\ifthenelse{\equal{\fontdocument}{universaliscondensed}}{
	\usepackage[condensed,sfdefault]{universalis}
}{
\ifthenelse{\equal{\fontdocument}{venturis}}{
	\usepackage[lf]{venturis}
	\renewcommand*\familydefault{\sfdefault}
}{
	\throwbadconfig[nostop]{Fuente desconocida}{\fontdocument}{(Fuentes clasicas) lmodern,arial,arial2,helvet,times,alegreyasans,mathpazo}
	\throwbadconfig[noheader-nostop]{Fuente desconocida}{\fontdocument}{(Fuentes adicionales) accantis,alegreya,algolrevived,antiqua,antpolt,antpoltlight,anttor,anttorcondensed,anttorlight,anttorlightcondensed,arev,arimo,aurical,avant,baskervald,berasans,beraserif,biolinum,cabin,cabincondensed,cantarell,caladea,carlito,chivolight,chivoregular,clearsans,comfortaa,comicneue,comicneueangular,crimson,cyklop,dejavusans,dejavusanscondensed,droidsans,firasans,iwona,iwonacondensed,iwonalight,iwonalightcondensed,kurier}
	\throwbadconfig[noheader-nostop]{Fuente desconocida}{\fontdocument}{kuriercondensed,kurierlight,kurierlightcondensed,lato,libris,lxfonts,merriweather,merriweatherlight,mintspirit,montserratalternatesextralight,montserratalternatesregular,montserratalternatesthin,montserratextralight,montserratlight,montserratregular,montserratthin,nimbussans,noto,opensans,overlock,paratype,paratypesanscaption,paratypesansnarrow,quattrocento,raleway,roboto,robotolight,robotolightcondensed,robotothin,rosario,sourcesanspro,uarial,ugq}
	\throwbadconfig[noheader]{Fuente desconocida}{\fontdocument}{universalis,universaliscondensed,venturis}
	}}}}}}}}}}}}}}}}}}}}}}}}}}}}}}}}}}}}}}}}}}}}}}}}}}}}}}}}}}}}}}}}}}}}}}}}}}}}}}}}}}}}
}

% -----------------------------------------------------------------------------
% TIPOGRAFÍA TYPEWRITER
% -----------------------------------------------------------------------------
\ifthenelse{\equal{\fonttypewriter}{tmodern}}{ % Default
	\renewcommand*\ttdefault{lmvtt}
}{
\ifthenelse{\equal{\fonttypewriter}{anonymouspro}}{
	\usepackage[ttdefault=true]{AnonymousPro}
}{
\ifthenelse{\equal{\fonttypewriter}{ascii}}{
	\usepackage{ascii}
	\let\SI\relax
}{
\ifthenelse{\equal{\fonttypewriter}{beramono}}{
	\usepackage[scaled]{beramono}
}{
\ifthenelse{\equal{\fonttypewriter}{cmpica}}{
	\usepackage{addfont}
	\addfont{OT1}{cmpica}{\pica}
	\addfont{OT1}{cmpicab}{\picab}
	\addfont{OT1}{cmpicati}{\picati}
	\renewcommand*\ttdefault{pica}
}{
\ifthenelse{\equal{\fonttypewriter}{courier}}{
	\usepackage{courier}
}{
\ifthenelse{\equal{\fonttypewriter}{dejavusansmono}}{
	\usepackage[scaled]{DejaVuSansMono}
}{
\ifthenelse{\equal{\fonttypewriter}{firamono}}{
	\usepackage[scale=0.85]{FiraMono}
}{
\ifthenelse{\equal{\fonttypewriter}{gomono}}{
	\usepackage[scale=0.85]{GoMono}
}{
\ifthenelse{\equal{\fonttypewriter}{inconsolata}}{
	\usepackage{inconsolata}
}{
\ifthenelse{\equal{\fonttypewriter}{nimbusmono}}{
	\usepackage{nimbusmono}
}{
\ifthenelse{\equal{\fonttypewriter}{newtxtt}}{
	\usepackage[zerostyle=d]{newtxtt}
}{
\ifthenelse{\equal{\fonttypewriter}{nimbusmono}}{
	\usepackage{nimbusmono}
}{
\ifthenelse{\equal{\fonttypewriter}{nimbusmononarrow}}{
	\usepackage{nimbusmononarrow}
}{
\ifthenelse{\equal{\fonttypewriter}{lcmtt}}{
	\renewcommand*\ttdefault{lcmtt}
}{
\ifthenelse{\equal{\fonttypewriter}{sourcecodepro}}{
	\usepackage[ttdefault=true,scale=0.85]{sourcecodepro}
}{
\ifthenelse{\equal{\fonttypewriter}{texgyrecursor}}{
	\usepackage{tgcursor}
}{
	\throwbadconfig{Fuente desconocida}{\fonttypewriter}{anonymouspro,ascii,beramono,
		cmpica,courier,dejavusansmono,firamono,gomono,inconsolata,kpmonospaced,lcmtt,
		newtxtt,nimbusmono,nimbusmononarrow,texgyrecursor,tmodern}
	}}}}}}}}}}}}}}}}
}

% -----------------------------------------------------------------------------
% FINALES
% -----------------------------------------------------------------------------
\usepackage[T1]{fontenc} % Caracteres acentuados
\ifthenelse{\equal{\showlinenumbers}{true}}{ % Muestra los números de línea
	\usepackage[switch,columnwise,running]{lineno}}{
}
\usepackage{csquotes} % Citas y comillas, se debe usar después de lineno [6.4.2]
\inputencoding{utf8}

% END

% IMPORTACIÓN DE FUNCIONES Y ENTORNOS
% Template:     Informe/Reporte LaTeX
% Documento:    Carga las funciones del template
% Versión:      5.2.1 (21/05/2018)
% Codificación: UTF-8
%
% Autor: Pablo Pizarro R. @ppizarror
%        Facultad de Ciencias Físicas y Matemáticas
%        Universidad de Chile
%        pablo.pizarro@ing.uchile.cl, ppizarror.com
%
% Manual template: [http://latex.ppizarror.com/Template-Informe/]
% Licencia MIT:    [https://opensource.org/licenses/MIT/]

% Template:     Informe/Reporte LaTeX
% Documento:    Funciones del núcleo del template
% Versión:      5.5.1 (04/09/2018)
% Codificación: UTF-8
%
% Autor: Pablo Pizarro R. @ppizarror
%        Facultad de Ciencias Físicas y Matemáticas
%        Universidad de Chile
%        pablo.pizarro@ing.uchile.cl, ppizarror.com
%
% Manual template: [http://latex.ppizarror.com/Template-Informe/]
% Licencia MIT:    [https://opensource.org/licenses/MIT/]

% Lanza un mensaje de error
% 	#1	Función del error
%	#2	Mensaje
\newcommand{\throwerror}[2]{
	\errmessage{LaTeX Error: \noexpand#1 #2 (linea \the\inputlineno)}
	\stop
}

% Lanza un mensaje de advertencia
%	#1	Mensaje
\newcommand{\throwwarning}[1]{
	\errmessage{LaTeX Warning: #1 (linea \the\inputlineno)}
}

% Lanza un mensaje de error indicando mala configuración dentro de begin{document}
%	#1	Mensaje de error
% 	#2	Configuración usada
%	#3	Valores esperados
\newcommand{\throwbadconfigondoc}[3]{
	\errmessage{#1 \noexpand #2=#2. Valores esperados: #3}
	\stop
}

% Comprueba si una variable está definida
%	#1	Variable
\newcommand{\checkvardefined}[1]{
	\ifthenelse{\isundefined{#1}}{
		\errmessage{LaTeX Warning: Variable \noexpand#1 no definida}
		\stop}{
	}
}

% Comprueba si una variable está definida
%	#1	Variable
%	#2	Mensaje
\newcommand{\checkextravarexist}[2]{
	\ifthenelse{\isundefined{#1}}{
		\errmessage{LaTeX Warning: Variable \noexpand#1 no definida}
		\ifx\hfuzz#2\hfuzz
			\errmessage{LaTeX Warning: Defina la variable en el bloque de INFORMACION DEL DOCUMENTO al comienzo del archivo principal del Template}
		\else
			\errmessage{LaTeX Warning: #2}
		\fi}{
	}
}

% Lanza un mensaje de error si una variable no ha sido definida
% 	#1	Función del error
%	#2	Variable
%	#3	Mensaje
\newcommand{\emptyvarerr}[3]{
	\ifx\hfuzz#2\hfuzz
		\errmessage{LaTeX Warning: \noexpand#1 #3 (linea \the\inputlineno)}
	\fi
}

% Cambiar el margen de los caption
% 	#1	Margen en centímetros
\newcommand{\setcaptionmargincm}[1]{
	\captionsetup{margin=#1cm}
}

% Cambia márgenes de las páginas [cm]
% 	#1	Margen izquierdo
%	#2	Margen superior
%	#3	Margen derecho
%	#4	Margen inferior
\newcommand{\setpagemargincm}[4]{
	\newgeometry{left=#1cm, top=#2cm, right=#3cm, bottom=#4cm}
}

% Cambia los márgenes del documento
%	#1 Margen izquierdo
%	#2 Margen derecho
\newcommand{\changemargin}[2]{
	\emptyvarerr{\changemargin}{#1}{Margen izquierdo no definido}
	\emptyvarerr{\changemargin}{#2}{Margen derecho no definido}
	\list{}{\rightmargin#2\leftmargin#1}\item[]
}
\let\endchangemargin=\endlist

% Imagen de prueba tikz
\newcommand{\bgtemplatetestimg}{
	\definecolor{c040302}{RGB}{4, 3, 2}
	\definecolor{cb26929}{RGB}{178,105, 41}
	\definecolor{cfa943c}{RGB}{250, 148, 60}
	\definecolor{cfdc594}{RGB}{253, 197, 148}
	\begin{tikzpicture}[y=1.0pt, x=1.0pt, yscale=-1, xscale=1, inner sep=0pt, outer sep=0pt, opacity=0.1]
	\path[fill=cfdc594](46.3187,318.2550)..controls(46.0737,317.5650)and
	(45.7487,314.0750)..(45.5987,310.5)--(45.3247,304)--(37.8247,303.5)--(30.3247,303)--(29.8247,295.5)-- (29.3247,288)--(21.8247,287.5)--(14.3247,287)--(13.8247,271.5)--(13.3247,256)--(5.8247,255.5)-- (-1.6753,255)--(-1.6753,206.5)--(-1.6753,158)--(5.8247,157.5)--(13.3247,157)--(13.8247,141.5)-- (14.3247,126)--(21.8247,125.5)--(29.3247,125)--(29.8247,109.5)--(30.3247,94)--(37.8247,93.5)-- (45.3247,93)--(45.8247,69.5)--(46.3247,46)--(53.8247,45.5)--(61.3247,45)--(61.8247,29.5)-- (62.3247,14)--(69.8247,13.5)--(77.3247,13)--(77.8247,5.5)--(78.3247,-2)--(94.8247,-2)-- (111.3247,-2)--(111.8247,5.5)--(112.3247,13)--(119.8247,13.5)--(127.3247,14)--(127.8247,21.5)-- (128.3247,29)--(166.8247,29)--(205.3247,29)--(205.8247,21.5)--(206.3247,14)--(213.8247,13.5)-- (221.3247,13)--(221.8247,5.5)--(222.3247,-2)--(230.8247,-2)--(239.3247,-2)--(239.8247,5.5)-- (240.3247,13)--(247.8247,13.5)--(255.3247,14)--(255.8247,21.5)--(256.3247,29)--(263.8247,29.5)-- (271.3247,30)--(271.8247,45.5)--(272.3247,61)--(279.8247,61.5)--(287.3247,62)--(287.8247,85.5)-- (288.3247,109)--(295.8247,109.5)--(303.3247,110)--(303.8247,133.5)--(304.3247,157)--(311.8247,157.5)-- (319.3247,158)--(319.3247,198.5)--(319.3247,239)--(311.8247,239.5)--(304.3247,240)--(303.8247,255.5)-- (303.3247,271)--(295.8247,271.5)--(288.3247,272)--(287.8247,279.5)--(287.3247,287)--(279.8247,287.5)-- (272.3247,288)--(271.8247,295.5)--(271.3247,303)--(263.8247,303.5)--(256.3247,304)--(255.8247,311.5)-- (255.3247,319)--(151.0447,319.2550)..controls(68.1357,319.4570)and (46.6737,319.2530)..(46.3187,318.2550)--cycle(173.8247,142.5)-- (173.8247,127.5)--(166.8247,127.5)--(159.8247,127.5)-- (159.8247,142.5)--(159.8247,157.5)--(166.8247,157.5)-- (173.8247,157.5)--cycle(269.8247,126.5)--(269.8247,111.5)-- (262.8247,111.5)--(255.8247,111.5)--(255.8247,126.5)-- (255.8247,141.5)--(262.8247,141.5)--(269.8247,141.5)--cycle; \path[fill=cfa943c](47.3187,317.2550)..controls(47.0737,316.5650)and (46.7487,313.0750)..(46.5987,309.5)--(46.3247,303)--(38.8247,302.5)--(31.3247,302)--(30.8247,294.5)-- (30.3247,287)--(22.8247,286.5)--(15.3247,286)--(14.8247,270.5)--(14.3247,255)--(6.8247,254.5)-- (-0.6753,254)--(-0.6753,206.5)--(-0.6753,159)--(6.8247,158.5)--(14.3247,158)--(14.8247,142.5)-- (15.3247,127)--(22.8247,126.5)--(30.3247,126)--(30.8247,110.5)--(31.3247,95)--(38.8247,94.5)-- (46.3247,94)--(46.8247,70.5)--(47.3247,47)--(54.8247,46.5)--(62.3247,46)--(62.8247,30.5)-- (63.3247,15)--(70.8247,14.5)--(78.3247,14)--(78.8247,6.5)--(79.3247,-1)--(94.8247,-1)-- (110.3247,-1)--(110.8247,6.5)--(111.3247,14)--(118.8247,14.5)--(126.3247,15)--(126.8247,22.5)-- (127.3247,30)--(166.8247,30)--(206.3247,30)--(206.8247,22.5)--(207.3247,15)--(214.8247,14.5)-- (222.3247,14)--(222.8247,6.5)--(223.3247,-1)--(230.8247,-1)--(238.3247,-1)--(238.8247,6.5)-- (239.3247,14)--(246.8247,14.5)--(254.3247,15)--(254.8247,22.5)--(255.3247,30)--(262.8247,30.5)-- (270.3247,31)--(270.8247,46.5)--(271.3247,62)--(278.8247,62.5)--(286.3247,63)--(286.8247,86.5)-- (287.3247,110)--(294.8247,110.5)--(302.3247,111)--(302.8247,134.5)--(303.3247,158)--(310.8247,158.5)-- (318.3247,159)--(318.3247,198.5)--(318.3247,238)--(310.8247,238.5)--(303.3247,239)--(302.8247,254.5)-- (302.3247,270)--(294.8247,270.5)--(287.3247,271)--(286.8247,278.5)--(286.3247,286)--(278.8247,286.5)-- (271.3247,287)--(270.8247,294.5)--(270.3247,302)--(262.8247,302.5)--(255.3247,303)--(254.8247,310.5)-- (254.3247,318)--(151.0447,318.2550)..controls(68.9357,318.4570)and (47.6737,318.2520)..(47.3187,317.2550)--cycle(253.8247,294.70).. controls(253.8247,290.7440)and(254.3267,287.3980)..(255.0247,286.70).. controls(255.7227,286.0020)and(259.0687,285.5)..(263.0247,285.5)-- (269.8247,285.5)--(269.8247,278.70)..controls(269.8247,274.7440)and (270.3267,271.3980)..(271.0247,270.70)..controls(271.7227,270.0020)and (275.0687,269.5)..(279.0247,269.5)--(285.8247,269.5)-- (285.8247,254.70)..controls(285.8247,244.5220)and(286.1997,239.5250).. (287.0247,238.70)..controls(287.7227,238.0020)and(291.0687,237.5).. (295.0247,237.5)--(301.8247,237.5)--(301.8247,206)-- (301.8247,174.5)--(294.8247,174.5)--(287.8247,174.5)-- (287.8247,189.80)..controls(287.8247,200.3670)and(287.4537,205.4710).. (286.6247,206.30)..controls(285.9267,206.9980)and(282.5807,207.5).. (278.6247,207.5)--(271.8247,207.5)--(271.8247,214.30)..controls (271.8247,218.2560)and(271.3227,221.6020)..(270.6247,222.30)..controls (269.9267,222.9980)and(266.5807,223.5)..(262.6247,223.5)-- (255.8247,223.5)--(255.8247,230.3780)..controls(255.8247,234.9380)and (255.3687,237.6340)..(254.4727,238.3780)..controls(252.4847,240.0280)and (176.6787,239.9540)..(175.0247,238.30)..controls(174.3337,237.6090)and (173.8247,234.30)..(173.8247,230.5)..controls(173.8247,226.70)and (174.3337,223.3910)..(175.0247,222.70)..controls(175.9167,221.8080)and (186.1807,221.5)..(215.0247,221.5)--(253.8247,221.5)-- (253.8247,214.5)--(253.8247,207.5)--(249.0747,207.3720)..controls (246.4627,207.3020)and(235.9997,206.9640)..(225.8247,206.6220)-- (207.3247,206)--(206.8247,198.5)--(206.3247,191)--(190.8247,190.5)--(175.3247,190)--(174.8247,182.5)-- (174.3247,175)--(134.8247,175)--(95.3247,175)--(94.8247,182.5)--(94.3247,190)--(86.8247,190.5)-- (79.3247,191)--(79.3247,222.5)--(79.3247,254)--(86.8247,254.5)--(94.3247,255)--(94.8247,270.5)-- (95.3247,286)--(102.8247,286.5)--(110.3247,287)-- (110.6187,294.2500)--(110.9127,301.5)--(182.3687,301.5)-- (253.8247,301.5)--cycle(174.6247,158.30)..controls (176.4747,156.4500)and(176.2867,128.1290)..(174.4147,126.5750)..controls (173.5527,125.8600)and(170.3467,125.5210)..(166.1647,125.7030)-- (159.3247,126)--(159.0477,141.4580)..controls(158.8957,149.9600)and (158.9937,157.4970)..(159.2667,158.2080)..controls(159.9187,159.9070)and (172.9397,159.9850)..(174.6257,158.30)--cycle(270.6247,142.30).. controls(272.4747,140.4500)and(272.2867,112.1290)..(270.4147,110.5750).. controls(269.5527,109.8600)and(266.3467,109.5210)..(262.1647,109.7030)-- (255.3247,110)--(255.0477,125.4580)..controls(254.8957,133.9600)and (254.9937,141.4970)..(255.2667,142.2080)..controls(255.9187,143.9070)and (268.9397,143.9850)..(270.6257,142.30)--cycle(78.8247,86.5)-- (79.3247,79)--(86.8247,78.5)--(94.3247,78)--(94.8247,70.5)--(95.3247,63)--(102.8247,62.5)-- (110.3247,62)--(110.8247,54.5)--(111.3247,47)-- (118.5747,46.7060)--(125.8247,46.4120)--(125.8247,38.9560)-- (125.8247,31.5)--(119.0247,31.5)..controls(115.0687,31.5)and (111.7227,30.9980)..(111.0247,30.30)..controls(110.3267,29.6020)and (109.8247,26.2560)..(109.8247,22.30)--(109.8247,15.5)-- (94.8247,15.5)--(79.8247,15.5)--(79.8247,30.30)..controls (79.8247,40.4780)and(79.4497,45.4750)..(78.6247,46.30)..controls (77.9267,46.9980)and(74.5807,47.5)..(70.6247,47.5)-- (63.8247,47.5)--(63.8247,71.0440)--(63.8247,94.5880)-- (71.0747,94.2940)--(78.3247,94)--cycle(253.8247,39)-- (253.8247,31.5)--(247.0247,31.5)..controls(243.0687,31.5)and (239.7227,30.9980)..(239.0247,30.30)..controls(238.3267,29.6020)and (237.8247,26.2560)..(237.8247,22.30)--(237.8247,15.5)-- (230.8247,15.5)--(223.8247,15.5)--(223.8247,22.30)..controls (223.8247,26.2560)and(223.3227,29.6020)..(222.6247,30.30)..controls (221.9267,30.9980)and(218.5807,31.5)..(214.6247,31.5)-- (207.8247,31.5)--(207.8247,39)--(207.8247,46.5)-- (230.8247,46.5)--(253.8247,46.5)--cycle; \path[fill=cb26929](47.3187,317.2550)..controls(47.0737,316.5650)and (46.7487,313.0750)..(46.5987,309.5)--(46.3247,303)--(38.8247,302.5)--(31.3247,302)--(30.8247,294.5)-- (30.3247,287)--(22.8247,286.5)--(15.3247,286)--(14.8247,270.5)--(14.3247,255)--(6.8247,254.5)-- (-0.6753,254)--(-0.6753,206.5)--(-0.6753,159)--(6.8247,158.5)--(14.3247,158)--(14.8247,142.5)-- (15.3247,127)--(22.8247,126.5)--(30.3247,126)--(30.8247,110.5)--(31.3247,95)--(38.8247,94.5)-- (46.3247,94)--(46.8247,70.5)--(47.3247,47)--(54.8247,46.5)--(62.3247,46)--(62.8247,30.5)-- (63.3247,15)--(70.8247,14.5)--(78.3247,14)--(78.8247,6.5)--(79.3247,-1)--(94.8247,-1)-- (110.3247,-1)--(110.8247,6.5)--(111.3247,14)--(118.8247,14.5)--(126.3247,15)--(126.8247,22.5)-- (127.3247,30)--(166.8247,30)--(206.3247,30)--(206.8247,22.5)--(207.3247,15)--(214.8247,14.5)-- (222.3247,14)--(222.8247,6.5)--(223.3247,-1)--(230.8247,-1)--(238.3247,-1)--(238.8247,6.5)-- (239.3247,14)--(246.8247,14.5)--(254.3247,15)--(254.8247,22.5)--(255.3247,30)--(262.8247,30.5)-- (270.3247,31)--(270.8247,46.5)--(271.3247,62)--(278.8247,62.5)--(286.3247,63)--(286.8247,86.5)-- (287.3247,110)--(294.8247,110.5)--(302.3247,111)--(302.8247,134.5)--(303.3247,158)--(310.8247,158.5)-- (318.3247,159)--(318.3247,198.5)--(318.3247,238)--(310.8247,238.5)--(303.3247,239)--(302.8247,254.5)-- (302.3247,270)--(294.8247,270.5)--(287.3247,271)--(286.8247,278.5)--(286.3247,286)--(278.8247,286.5)-- (271.3247,287)--(270.8247,294.5)--(270.3247,302)--(262.8247,302.5)--(255.3247,303)--(254.8247,310.5)-- (254.3247,318)--(151.0447,318.2550)..controls(68.9357,318.4570)and (47.6737,318.2520)..(47.3187,317.2550)--cycle(253.8247,294.70).. controls(253.8247,290.7440)and(254.3267,287.3980)..(255.0247,286.70).. controls(255.7227,286.0020)and(259.0687,285.5)..(263.0247,285.5)-- (269.8247,285.5)--(269.8247,278.70)..controls(269.8247,274.7440)and (270.3267,271.3980)..(271.0247,270.70)..controls(271.7227,270.0020)and (275.0687,269.5)..(279.0247,269.5)--(285.8247,269.5)-- (285.8247,254.70)..controls(285.8247,244.5220)and(286.1997,239.5250).. (287.0247,238.70)..controls(287.7227,238.0020)and(291.0717,237.5).. (295.0367,237.5)--(301.8487,237.5)--(301.5867,198.7500)-- (301.3247,160)--(294.8247,160)--(288.3247,160)-- (288.0547,182.4210)..controls(287.8807,196.9060)and(287.3947,205.3120).. (286.6827,206.1710)..controls(285.9607,207.0410)and(283.2017,207.5).. (278.7027,207.5)--(271.8247,207.5)--(271.8247,214.30)..controls (271.8247,218.2560)and(271.3227,221.6020)..(270.6247,222.30)..controls (269.9267,222.9980)and(266.5807,223.5)..(262.6247,223.5)-- (255.8247,223.5)--(255.8247,230.3780)..controls(255.8247,234.9380)and (255.3687,237.6340)..(254.4727,238.3780)..controls(252.4847,240.0280)and (176.6787,239.9540)..(175.0247,238.30)..controls(174.3337,237.6090)and (173.8247,234.30)..(173.8247,230.5)..controls(173.8247,226.70)and (174.3337,223.3910)..(175.0247,222.70)..controls(175.9167,221.8080)and (186.1807,221.5)..(215.0247,221.5)--(253.8247,221.5)-- (253.8247,214.5)--(253.8247,207.5)--(247.0747,207.4850)..controls (243.3627,207.4770)and(239.7747,207.1200)..(239.1027,206.6930)..controls (238.2507,206.1520)and(237.7967,201.3880)..(237.6027,190.9580)-- (237.3247,176)--(229.8247,175.5)--(222.3247,175)--(222.3247,166.5)--(222.3247,158)--(253.8247,157.5)-- (285.3247,157)--(285.3247,134.5)--(285.3247,112)--(277.8247,111.5)--(270.3247,111)--(269.8247,87.5)-- (269.3247,64)--(261.8247,63.5)--(254.3247,63)-- (254.0497,47.2500)--(253.7737,31.5)--(246.9987,31.5)..controls (243.0637,31.5)and(239.7217,30.9970)..(239.0247,30.30)..controls (238.3267,29.6020)and(237.8247,26.2560)..(237.8247,22.30)-- (237.8247,15.5)--(230.8247,15.5)--(223.8247,15.5)-- (223.8247,22.30)..controls(223.8247,26.2560)and(223.3227,29.6020).. (222.6247,30.30)..controls(221.9277,30.9970)and(218.5897,31.5).. (214.6667,31.5)--(207.9087,31.5)--(207.6167,39.2500)-- (207.3247,47)--(166.8247,47)--(126.3247,47)-- (126.0327,39.2500)--(125.7407,31.5)--(118.9827,31.5)..controls (115.0597,31.5)and(111.7217,30.9970)..(111.0247,30.30)..controls (110.3267,29.6020)and(109.8247,26.2560)..(109.8247,22.30)-- (109.8247,15.5)--(94.8247,15.5)--(79.8247,15.5)-- (79.8247,30.30)..controls(79.8247,40.4780)and(79.4497,45.4750).. (78.6247,46.30)..controls(77.9277,46.9970)and(74.5847,47.5).. (70.6427,47.5)--(63.8607,47.5)--(63.5927,71.2500)-- (63.3247,95)--(55.8247,95.5)--(48.3247,96)--(47.8247,111.5)--(47.3247,127)--(39.8247,127.5)-- (32.3247,128)--(31.8247,143.5)--(31.3247,159)--(23.8247,159.5)--(16.3247,160)--(16.3247,175)-- (16.3247,190)--(23.3247,190.5)--(30.3247,191)--(30.8247,198.5)--(31.3247,206)--(38.8247,206.5)-- (46.3247,207)--(46.8247,214.5)--(47.3247,222)--(62.8247,222.5)--(78.3247,223)--(78.8247,238.5)-- (79.3247,254)--(86.8247,254.5)--(94.3247,255)--(94.8247,270.5)--(95.3247,286)--(102.8247,286.5)-- (110.3247,287)--(110.6187,294.2500)--(110.9127,301.5)-- (182.3687,301.5)--(253.8247,301.5)--cycle(126.2657,158.2080).. controls(125.9927,157.4970)and(125.8947,149.9600)..(126.0467,141.4580)-- (126.3247,126)--(142.5747,125.7250)--(158.8247,125.4500)-- (158.8247,142.4750)--(158.8247,159.5)--(142.7937,159.5)..controls (130.6137,159.5)and(126.6437,159.1900)..(126.2667,158.2080)-- cycle(222.2657,142.2080)..controls(221.9927,141.4970)and (221.8947,133.9600)..(222.0467,125.4580)--(222.3247,110)-- (238.5747,109.7250)--(254.8247,109.4500)--(254.8247,126.4750)-- (254.8247,143.5)--(238.7937,143.5)..controls(226.6137,143.5)and (222.6437,143.1900)..(222.2667,142.2080)--cycle; \path[fill=c040302](46.3247,303)--(79.0747,303.0080)..controls (95.9877,303.0120)and(142.2457,303.1250)..(181.8707,303.2580)-- (253.9167,303.5)--(253.6207,310.2500)--(253.3247,317)-- (151.0447,317.2550)..controls(65.3137,317.4680)and(48.6807,317.2880).. (48.2407,316.1430)..controls(46.9629,307.2110)and(46.3247,303).. (46.3247,303)--cycle(32.2567,300.1840)..controls(31.9597,299.4100) and(31.8537,296.1260)..(32.0207,292.8880)--(32.3247,287)-- (39.3247,287)--(46.3247,287)--(46.3247,294)-- (46.3247,301)--(39.5607,301.2960)..controls(34.4617,301.5190)and (32.6637,301.2460)..(32.2567,300.1840)--cycle(255.8247,294.5460)-- (255.8247,287.5)--(262.8707,287.5)--(269.9167,287.5)-- (269.6207,294.2500)--(269.3247,301)--(262.5747,301.2960)-- (255.8247,301.5920)--cycle(16.2817,284.2480)..controls(15.9977,283.5090) and(15.8917,276.6260)..(16.0457,268.9520)--(16.3247,255)-- (23.3247,255)--(30.3247,255)--(30.3247,270)-- (30.3247,285)--(23.5607,285.2960)..controls(18.6047,285.5130)and (16.6587,285.2330)..(16.2807,284.2480)--cycle(271.8247,278.5460)-- (271.8247,271.5)--(278.8707,271.5)--(285.9167,271.5)-- (285.6207,278.2500)--(285.3247,285)--(278.5747,285.2960)-- (271.8247,285.5920)--cycle(287.8247,254.5460)--(287.8247,239.5)-- (294.8517,239.5)--(301.8787,239.5)--(301.6017,254.2500)-- (301.3247,269)--(294.5747,269.2960)--(287.8247,269.5920)-- cycle(0.3007,252.2960)..controls(0.0267,251.5830)and(-0.0793,230.5250).. (0.0637,205.5)--(0.3247,160)--(6.9587,159.7070)-- (13.5927,159.4140)--(14.2807,190.7070)..controls(14.6597,207.9180)and (14.8247,228.9750)..(14.6467,237.5)--(14.3247,253)-- (7.5607,253.2960)..controls(2.7217,253.5080)and(0.6557,253.2230).. (0.2997,252.2960)--cycle(175.8247,230.5)--(175.8247,223.5)-- (214.8247,223.5)--(253.8247,223.5)--(253.8247,230.5)-- (253.8247,237.5)--(214.8247,237.5)--(175.8247,237.5)-- cycle(303.8247,198.5)--(303.8247,159.4090)--(310.5747,159.7050)-- (317.3247,160.0010)--(317.3247,198.5010)--(317.3247,237.0010)-- (310.5747,237.2970)--(303.8247,237.5930)--cycle(255.8247,214.70).. controls(255.8247,210.7440)and(255.3227,207.3980)..(254.6247,206.70).. controls(253.9267,206.0020)and(250.5807,205.5)..(246.6247,205.5)-- (239.8247,205.5)--(239.8247,190.70)..controls(239.8247,180.5220)and (239.4497,175.5250)..(238.6247,174.70)..controls(237.9267,174.0020)and (234.5807,173.5)..(230.6247,173.5)--(223.8247,173.5)-- (223.8247,166.5)--(223.8247,159.5)--(254.8247,159.5)-- (285.8247,159.5)--(285.8247,182.5)--(285.8247,205.5)-- (279.0247,205.5)..controls(275.0687,205.5)and(271.7227,206.0020).. (271.0247,206.70)..controls(270.3267,207.3980)and(269.8247,210.7440).. (269.8247,214.70)--(269.8247,221.5)--(262.8247,221.5)-- (255.8247,221.5)--cycle(16.0477,142.7500)--(16.3247,128)-- (23.0747,127.7040)--(29.8247,127.4090)--(29.8247,142.4540)-- (29.8247,157.5)--(22.7977,157.5)--(15.7717,157.5)-- cycle(127.8247,142.5)--(127.8247,127.5)--(142.8247,127.5)-- (157.8247,127.5)--(157.8247,142.5)--(157.8247,157.5)-- (142.8247,157.5)--(127.8247,157.5)--cycle(287.8247,134.4540)-- (287.8247,111.4080)--(294.5747,111.7040)--(301.3247,112)-- (301.5937,134.7500)--(301.8627,157.5)--(294.8437,157.5)-- (287.8247,157.5)--cycle(223.8247,126.5)--(223.8247,111.5)-- (238.8247,111.5)--(253.8247,111.5)--(253.8247,126.5)-- (253.8247,141.5)--(238.8247,141.5)--(223.8247,141.5)-- cycle(32.0477,110.7500)--(32.3247,96)--(39.0747,95.7040)-- (45.8247,95.4090)--(45.8247,110.4540)--(45.8247,125.5)-- (38.7977,125.5)--(31.7717,125.5)--cycle(271.8247,86.4540)-- (271.8247,63.4090)--(278.5747,63.7050)--(285.3247,64.0010)-- (285.5937,86.7510)--(285.8627,109.5010)--(278.8437,109.5010)-- (271.8247,109.5010)--cycle(48.0557,70.7500)--(48.3247,48)-- (55.0747,47.7040)--(61.8247,47.4090)--(61.8247,70.4540)-- (61.8247,93.5)--(54.8057,93.5)--(47.7867,93.5)-- cycle(255.8247,46.4540)--(255.8247,31.4090)--(262.5747,31.7050)-- (269.3247,32.0010)--(269.6017,46.7510)--(269.8787,61.5010)-- (262.8517,61.5010)--(255.8247,61.5010)--cycle(64.0477,30.7500)-- (64.3247,16)--(71.0747,15.7040)--(77.8247,15.4090)--(77.8247,30.4540)--(77.8247,45.5)--(70.7977,45.5)-- (63.7707,45.5)--cycle(127.8247,38.5)--(127.8247,31.5)--(166.8247,31.5)--(205.8247,31.5)--(205.8247,38.5)-- (205.8247,45.5)--(166.8247,45.5)--(127.8247,45.5)-- cycle(111.8247,22.4540)--(111.8247,15.4080)--(118.5747,15.7040)-- (125.3247,16)--(125.6207,22.7500)--(125.9167,29.5)-- (118.8707,29.5)--(111.8247,29.5)--cycle(208.0287,22.7500)-- (208.3247,16)--(215.0747,15.7040)--(221.8247,15.4090)-- (221.8247,22.4540)--(221.8247,29.5)--(214.7787,29.5)-- (207.7327,29.5)--cycle(239.8247,22.4540)--(239.8247,15.4080)-- (246.5747,15.7040)--(253.3247,16)--(253.6207,22.7500)-- (253.9167,29.5)--(246.8707,29.5)--(239.8247,29.5)--cycle(80.0287,6.7500)--(80.3247,0)--(94.8247,0)-- (109.3247,0)--(109.6207,6.7500)--(109.9167,13.5)--(94.8247,13.5)--(79.7337,13.5)--cycle(224.0287,6.7500)-- (224.3247,0)--(230.8247,0)--(237.3247,0)--(237.6207,6.7500)--(237.9167,13.5)--(230.8247,13.5)-- (223.7337,13.5)--cycle;
	\end{tikzpicture}
}
\def\bgtemplatetestcode {d0g3}

% Inicia el modo de página horizontal
\def\beginplandscape {
	\newgeometry{landscape,left=\pagemarginleft cm,top=\pagemargintop cm,right=\pagemarginright cm,bottom=\pagemarginbottom cm}
	\headwidth=\textwidth
}

% Termina el modo de página horizontal
\def\endplandscape {
	\restoregeometry
	\setpagemargincm{\pagemarginleft}{\pagemargintop}{\pagemarginright}{\pagemarginbottom}
	\headwidth=\textwidth
}

% Chequea que las funciones sólo puedan usarse en el entorno images
\newcommand{\checkonlyonenvimage}{
	\ifthenelse{\equal{\envimagesinitialized}{true}}{}{
		\throwwarning{Funciones \noexpand\addimage o \noexpand\addimageboxed no pueden usarse fuera del entorno \noexpand\images}
		\stop
	}
}

% Chequea que las funciones sólo puedan usarse fuera del entorno images
\newcommand{\checkoutsideenvimage}{
	\ifthenelse{\equal{\envimagesinitialized}{true}}{
		\throwwarning{Esta funcion solo puede usarse fuera del entorno \noexpand\images}
		\stop}{
	}
}

% Alerta función de inserción de imagen deprecada, utilizar entorno images
\newcommand{\alertdeprecatedcmdimage}[1]{
	\throwwarning{Funcion \noexpand#1 deprecada, se recomienda utilizar el entorno \noexpand\images en su lugar}
}

% Se importa la librería tikz
\newcommand{\coreimporttikz}{
	\ifthenelse{\equal{\importtikz}{false}}{
		\usepackage{tikz}}{
	}
}

% Template:     Informe/Reporte LaTeX
% Documento:    Funciones matemáticas
% Versión:      6.0.6 (30/10/2018)
% Codificación: UTF-8
%
% Autor: Pablo Pizarro R. @ppizarror
%        Facultad de Ciencias Físicas y Matemáticas
%        Universidad de Chile
%        pablo.pizarro@ing.uchile.cl, ppizarror.com
%
% Manual template: [https://latex.ppizarror.com/Template-Informe/]
% Licencia MIT:    [https://opensource.org/licenses/MIT/]

% Insertar sub-índice, a_b
% 	#1	Elemento inferior (a)
%	#2	Elemento superior (b)
\newcommand{\lpow}[2]{
	\ensuremath{{#1}_{#2}}
}

% Insertar elevado, a^b
% 	#1	Elemento inferior (a)
%	#2	Elemento superior (b)
\newcommand{\pow}[2]{
	\ensuremath{{#1}^{#2}}
}

% Inserta inverso función seno, sin^-1
%	#1	Elemento
\newcommand{\aasin}[1][]{
	\ifx\hfuzz#1\hfuzz
		\ensuremath{\sin^{-1}#1}
	\else
		\ensuremath{{\sin}^{-1}}
	\fi
}

% Inserta inverso función coseno, cos^-1
%	#1	Elemento
\newcommand{\aacos}[1][]{
	\ifx\hfuzz#1\hfuzz
		\ensuremath{\cos^{-1}#1}
	\else
		\ensuremath{\cos^{-1}}
	\fi
}

% Inserta inverso función tangente, tan^-1
%	#1	Elemento
\newcommand{\aatan}[1][]{
	\ifx\hfuzz#1\hfuzz
		\ensuremath{\tan^{-1}#1}
	\else
		\ensuremath{\tan^{-1}}
	\fi
}

% Inserta inverso función cosecante, csc^-1
%	#1	Elemento
\newcommand{\aacsc}[1][]{
	\ifx\hfuzz#1\hfuzz
		\ensuremath{\csc^{-1}#1}
	\else
		\ensuremath{\csc^{-1}}
	\fi
}

% Inserta inverso función secante, sec^-1
%	#1	Elemento
\newcommand{\aasec}[1][]{
	\ifx\hfuzz#1\hfuzz
		\ensuremath{\sec^{-1}#1}
	\else
		\ensuremath{\sec^{-1}}
	\fi
}

% Inserta inverso función cotangente, cot^-1
%	#1	Elemento
\newcommand{\aacot}[1][]{
	\ifx\hfuzz#1\hfuzz
		\ensuremath{\cot^{-1}#1}
	\else
		\ensuremath{\cot^{-1}}
	\fi
}

% Fracción de derivadas parciales af/ax
% 	#1	Función a derivar (f)
%	#2	Variable a derivar (x)
\newcommand{\fracpartial}[2]{
	\ensuremath{\pdv{#1}{#2}}
}

% Fracción de derivadas parciales dobles a^2f/ax^2
% 	#1	Función a derivar (f)
%	#2	Variable a derivar (x)
\newcommand{\fracdpartial}[2]{
	\ensuremath{\pdv[2]{#1}{#2}}
}

% Fracción de derivadas parciales en n, a^nf/ax^n
% 	#1	Función a derivar (f)
%	#2	Variable a derivar (x)
%	#3	Orden (n)
\newcommand{\fracnpartial}[3]{
	\ensuremath{\pdv[#3]{#1}{#2}}
}

% Fracción de derivadas df/dx
% 	#1	Función a derivar (f)
%	#2	Variable a derivar (x)
\newcommand{\fracderivat}[2]{
	\ensuremath{\dv{#1}{#2}}
}

% Fracción de derivadas dobles d^2/dx^2
% 	#1	Función a derivar (f)
%	#2	Variable a derivar (x)
\newcommand{\fracdderivat}[2]{
	\ensuremath{\dv[2]{#1}{#2}}
}

% Fracción de derivadas en n d^nf/dx^n
% 	#1	Función a derivar (f)
%	#2	Variable a derivar (x)
%	#3	Orden de la derivada (n)
\newcommand{\fracnderivat}[3]{
	\ensuremath{\dv[#3]{#1}{#2}}
}

% Llave superior de equivalencia
% 	#1	Elemento a igualar
%	#2	Igualdad
\newcommand{\topequal}[2]{
	\ensuremath{\overbrace{#1}^{\mathclap{#2}}}
}

% Llave inferior de equivalencia
% 	#1	Elemento a igualar
%	#2	Igualdad
\newcommand{\underequal}[2]{
	\ensuremath{\underbrace{#1}_{\mathclap{#2}}}
}

% Rectángulo superior de equivalencia
% 	#1	Elemento a igualar
%	#2	Igualdad
\newcommand{\topsequal}[2]{
	\ensuremath{\overbracket{#1}^{\mathclap{#2}}}
}

% Rectángulo inferior de equivalencia
% 	#1	Elemento a igualar
%	#2	Igualdad
\newcommand{\undersequal}[2]{
	\ensuremath{\underbracket{#1}_{\mathclap{#2}}}
}

% END
% Template:     Informe/Reporte LaTeX
% Documento:    Funciones para insertar ecuaciones
% Versión:      5.5.6 (30/09/2018)
% Codificación: UTF-8
%
% Autor: Pablo Pizarro R. @ppizarror
%        Facultad de Ciencias Físicas y Matemáticas
%        Universidad de Chile
%        pablo.pizarro@ing.uchile.cl, ppizarror.com
%
% Manual template: [http://latex.ppizarror.com/Template-Informe/]
% Licencia MIT:    [https://opensource.org/licenses/MIT/]

% Redimensiona una ecuación en textwidth
% 	#1	Tamaño del nuevo objeto (En textwidth)
%	#2	Ecuación a redimensionar
\newcommand{\equationresize}[2]{
	\emptyvarerr{\equationresize}{#1}{Dimension no definida}
	\emptyvarerr{\equationresize}{#2}{Ecuacion a redimensionar no definida}
	\resizebox{#1\textwidth}{!}{$#2$}
}

% Insertar una ecuación
% 	#1	Label (opcional)
%	#2	Ecuación
\newcommand{\insertequation}[2][]{
	\emptyvarerr{\insertequation}{#2}{Ecuacion no definida}
	\ifthenelse{\equal{\numberedequation}{true}}{
		\vspace{-0.1cm}
		\begin{equation}
			\text{#1} #2
		\end{equation}
		\vspace{-0.26cm}
	}{
		\ifx\hfuzz#1\hfuzz
		\else
			\throwwarning{Label invalido en ecuacion sin numero}
		\fi
		\insertequationanum{#2}
	}
}

% Insertar una ecuación sin número
%	#1	Ecuación
\newcommand{\insertequationanum}[1]{
	\emptyvarerr{\insertequationanum}{#1}{Ecuacion no definida}
	\vspace{-0.1cm}
	\begin{equation*}
		\ensuremath{#1}
	\end{equation*}
	\vspace{-0.26cm}
}

% Insertar una ecuación con leyenda
% 	#1	Label (opcional)
%	#2	Ecuación
%	#3	Leyenda
\newcommand{\insertequationcaptioned}[3][]{
	\emptyvarerr{\insertequationcaptioned}{#2}{Ecuacion no definida}
	\ifx\hfuzz#3\hfuzz
		\insertequation[#1]{#2}
	\else
		\ifthenelse{\equal{\numberedequation}{true}}{
			\vspace{-0.10cm}
			\begin{equation}
				\text{#1} #2
			\end{equation}
			\vspace{-0.65cm}
			\begin{changemargin}{\captionlrmargin cm}{\captionlrmargin cm}
				\centering \textcolor{\captiontextcolor}{#3}
				\vspace{0.05cm}
			\end{changemargin}
			\vspace{0.05cm}
		}{
			\ifx\hfuzz#1\hfuzz
			\else
				\throwwarning{Label invalido en ecuacion sin numero}
			\fi
			\insertequationcaptionedanum{#2}{#3}
		}
	\fi
}

% Insertar una ecuación con leyenda sin número
%	#1	Ecuación
%	#2	Leyenda
\newcommand{\insertequationcaptionedanum}[2]{
	\emptyvarerr{\insertequationcaptionedanum}{#1}{Ecuacion no definida}
	\ifx\hfuzz#2\hfuzz
		\insertequationanum{#1}
	\else
		\vspace{-0.10cm}
		\begin{equation*}
			\ensuremath{#1}
		\end{equation*}
		\vspace{-0.65cm}
		\begin{changemargin}{\captionlrmargin cm}{\captionlrmargin cm}
			\centering \textcolor{\captiontextcolor}{#2}
			\vspace{0.05cm}
		\end{changemargin}
		\vspace{0.05cm}
	\fi
}

% Insertar una ecuación con el ambiente gather
%	#1	Ecuación
\newcommand{\insertgather}[1]{
	\emptyvarerr{\insertgather}{#1}{Ecuacion no definida}
	\ifthenelse{\equal{\numberedequation}{true}}{
		\vspace{-0.4cm}
		\begin{gather}
			\ensuremath{#1}
		\end{gather}
		\vspace{-0.4cm}
	}{
		\insertgatheranum{#1}
	}
}

% Insertar una ecuación con el ambiente gather sin número
%	#1	Ecuación
\newcommand{\insertgatheranum}[1]{
	\emptyvarerr{\insertgatheranum}{#1}{Ecuacion no definida}
	\vspace{-0.4cm}
	\begin{gather*}
		\ensuremath{#1}
	\end{gather*}
	\vspace{-0.4cm}
}

% Insertar una ecuación (gather) con leyenda
%	#1	Ecuación
%	#2	Leyenda
\newcommand{\insertgathercaptioned}[2]{
	\emptyvarerr{\insertgathercaptioned}{#1}{Ecuacion no definida}
	\ifx\hfuzz#2\hfuzz
		\insertgather{#1}
	\else
		\ifthenelse{\equal{\numberedequation}{true}}{
			\vspace{-0.45cm}
			\begin{gather}
				\ensuremath{#1}
			\end{gather}
			\vspace{-0.77cm}
			\begin{changemargin}{\captionlrmargin cm}{\captionlrmargin cm}
				\centering \textcolor{\captiontextcolor}{#2}
				\vspace{0.05cm}
			\end{changemargin}
			\vspace{0cm}
		}{
			\insertgathercaptionedanum{#1}{#2}
		}
	\fi
}

% Insertar una ecuación (gather) con leyenda sin número
%	#1	Ecuación
%	#2	Leyenda
\newcommand{\insertgathercaptionedanum}[2]{
	\emptyvarerr{\insertgathercaptionedanum}{#1}{Ecuacion no definida}
	\ifx\hfuzz#2\hfuzz
		\insertgatheranum{#1}
	\else
		\vspace{-0.45cm}
		\begin{gather*}
			\ensuremath{#1}
		\end{gather*}
		\vspace{-0.77cm}
		\begin{changemargin}{\captionlrmargin cm}{\captionlrmargin cm}
			\centering \textcolor{\captiontextcolor}{#2}
			\vspace{0.05cm}
		\end{changemargin}
		\vspace{0cm}
	\fi
}

% Insertar una ecuación con el ambiente gathered
% 	#1	Label (opcional)
%	#2	Ecuación
\newcommand{\insertgathered}[2][]{
	\emptyvarerr{\insertgathered}{#2}{Ecuacion no definida}
	\ifthenelse{\equal{\numberedequation}{true}}{
		\vspace{-0.1cm}
		\begin{equation}
			\begin{gathered}
				\text{#1} \ensuremath{#2}
			\end{gathered}
		\end{equation}
		\vspace{-0.05cm}
	}{
		\ifx\hfuzz#1\hfuzz
		\else
			\throwwarning{Label invalido en ecuacion (gathered) sin numero}
		\fi
		\insertgatheredanum{#2}
	}
}

% Insertar una ecuación con el ambiente gathered sin número
%	#1	Ecuación
\newcommand{\insertgatheredanum}[1]{
	\emptyvarerr{\insertgatheredanum}{#1}{Ecuacion no definida}
	\vspace{-0.4cm}
	\begin{gather*}
		\ensuremath{#1}
	\end{gather*}
	\vspace{-0.4cm}
}

% Insertar una ecuación (gathered) con leyenda
% 	#1	Label (opcional)
%	#2	Ecuación
%	#3	Leyenda
\newcommand{\insertgatheredcaptioned}[3][]{
	\emptyvarerr{\insertgatheredcaptioned}{#2}{Ecuacion no definida}
	\ifx\hfuzz#3\hfuzz
		\insertgathered[#1]{#2}
	\else
		\ifthenelse{\equal{\numberedequation}{true}}{
			\vspace{0cm}
			\begin{equation}
				\begin{gathered}
					\text{#1} \ensuremath{#2}
				\end{gathered}
			\end{equation}
			\vspace{-0.65cm}
			\begin{changemargin}{\captionlrmargin cm}{\captionlrmargin cm}
				\centering \textcolor{\captiontextcolor}{#3}
				\vspace{0.05cm}
			\end{changemargin}
			\vspace{0cm}
		}{
			\ifx\hfuzz#1\hfuzz
			\else
				\throwwarning{Label invalido en ecuacion (gathered) sin numero}
			\fi
			\insertgatheredcaptionedanum{#2}{#3}
		}
		\fi
}

% Insertar una ecuación (gathered) con leyenda sin número
%	#1	Ecuación
%	#2	Leyenda
\newcommand{\insertgatheredcaptionedanum}[2]{
	\emptyvarerr{\insertgatheredcaptionedanum}{#1}{Ecuacion no definida}
	\ifx\hfuzz#2\hfuzz
		\insertgatheredanum{#1}
	\else
		\vspace{-0.45cm}
		\begin{gather*}
			\ensuremath{#1}
		\end{gather*}
		\vspace{-0.7cm}
		\begin{changemargin}{\captionlrmargin cm}{\captionlrmargin cm}
			\centering \textcolor{\captiontextcolor}{#2}
			\vspace{0.05cm}
		\end{changemargin}
		\vspace{0cm}
	\fi
}

% Insertar una ecuación con el ambiente align
%	#1	Ecuación
\newcommand{\insertalign}[1]{
	\emptyvarerr{\insertalign}{#1}{Ecuacion no definida}
	\ifthenelse{\equal{\numberedequation}{true}}{
		\vspace{-0.45cm}
		\begin{align}
			\ensuremath{#1}
		\end{align}
		\vspace{-0.4cm}
	}{
		\insertalignanum{#1}
	}
}

% Insertar una ecuación con el ambiente align sin número
%	#1	Ecuación
\newcommand{\insertalignanum}[1]{
	\emptyvarerr{\insertalignanum}{#1}{Ecuacion no definida}
	\vspace{-0.45cm}
	\begin{align*}
		\ensuremath{#1}
	\end{align*}
	\vspace{-0.4cm}
}

% Insertar una ecuación (align) con leyenda
%	#1	Ecuación
%	#2	Leyenda
\newcommand{\insertaligncaptioned}[2]{
	\emptyvarerr{\insertaligncaptioned}{#1}{Ecuacion no definida}
	\ifx\hfuzz#2\hfuzz
		\insertalign{#1}
	\else
		\ifthenelse{\equal{\numberedequation}{true}}{
			\vspace{-0.45cm}
			\begin{align}
				\ensuremath{#1}
			\end{align}
			\vspace{-0.77cm}
			\begin{changemargin}{\captionlrmargin cm}{\captionlrmargin cm}
				\centering \textcolor{\captiontextcolor}{#2}
				\vspace{0.05cm}
			\end{changemargin}
			\vspace{0cm}
		}{
			\insertaligncaptionedanum{#1}{#2}
		}
	\fi
}

% Insertar una ecuación (align) con leyenda sin número
%	#1	Ecuación
%	#2	Leyenda
\newcommand{\insertaligncaptionedanum}[2]{
	\emptyvarerr{\insertaligncaptioned}{#1}{Ecuacion no definida}
	\ifx\hfuzz#2\hfuzz
		\insertalignanum{#1}
	\else
		\vspace{-0.45cm}
		\begin{align*}
			\ensuremath{#1}
		\end{align*}
		\vspace{-0.77cm}
		\begin{changemargin}{\captionlrmargin cm}{\captionlrmargin cm}
			\centering \textcolor{\captiontextcolor}{#2}
			\vspace{0.05cm}
		\end{changemargin}
		\vspace{0cm}
	\fi
}

% Insertar una ecuación con el ambiente aligned
% 	#1	Label (opcional)
%	#2	Ecuación
\newcommand{\insertaligned}[2][]{
	\emptyvarerr{\insertaligned}{#2}{Ecuacion no definida}
	\ifthenelse{\equal{\numberedequation}{true}}{
		\vspace{-0.1cm}
		\begin{equation}
			\begin{aligned}
				\text{#1} \ensuremath{#2}
			\end{aligned}
		\end{equation}
		\vspace{-0.05cm}
	}{
		\ifx\hfuzz#1\hfuzz
		\else
			\throwwarning{Label invalido en ecuacion (aligned) sin numero}
		\fi
		\insertalignedanum{#2}
	}
}

% Insertar una ecuación con el ambiente aligned sin número
%	#1	Ecuación
\newcommand{\insertalignedanum}[1]{
	\emptyvarerr{\insertalignedanum}{#1}{Ecuacion no definida}
	\vspace{-0.45cm}
	\begin{align*}
		\ensuremath{#1}
	\end{align*}
	\vspace{-0.4cm}
}

% Insertar una ecuación (aligned) con leyenda
% 	#1	Label (opcional)
%	#2	Ecuación
%	#3	Leyenda
\newcommand{\insertalignedcaptioned}[3][]{
	\emptyvarerr{\insertalignedcaptioned}{#2}{Ecuacion no definida}
	\ifx\hfuzz#3\hfuzz
		\insertaligned[#1]{#2}
	\else
		\ifthenelse{\equal{\numberedequation}{true}}{
			\vspace{0cm}
			\begin{equation}
				\begin{aligned}
					\text{#1} \ensuremath{#2}
				\end{aligned}
			\end{equation}
			\vspace{-0.65cm}
			\begin{changemargin}{\captionlrmargin cm}{\captionlrmargin cm}
				\centering \textcolor{\captiontextcolor}{#3}
				\vspace{0.05cm}
			\end{changemargin}
			\vspace{0cm}
		}{
			\ifx\hfuzz#1\hfuzz
			\else
				\throwwarning{Label invalido en ecuacion (aligned) sin numero}
			\fi
			\insertalignedcaptionedanum{#2}{#3}
		}
	\fi
}

% Insertar una ecuación (aligned) con leyenda sin número
%	#1	Ecuación
%	#2	Leyenda
\newcommand{\insertalignedcaptionedanum}[2]{
	\emptyvarerr{\insertalignedcaptioned}{#1}{Ecuacion no definida}
	\ifx\hfuzz#2\hfuzz
		\insertalignedanum{#1}
	\else
		\vspace{0cm}
		\begin{equation}
			\begin{aligned}
				\ensuremath{#1}
			\end{aligned}
		\end{equation}
		\vspace{-0.65cm}
		\begin{changemargin}{\captionlrmargin cm}{\captionlrmargin cm}
			\centering \textcolor{\captiontextcolor}{#2}
			\vspace{0.05cm}
		\end{changemargin}
		\vspace{0cm}
	\fi
}

% Template:     Informe/Reporte LaTeX
% Documento:    Funciones para insertar imágenes
% Versión:      4.7.6 (06/04/2018)
% Codificación: UTF-8
%
% Autor: Pablo Pizarro R.
%        Facultad de Ciencias Físicas y Matemáticas
%        Universidad de Chile
%        pablo.pizarro@ing.uchile.cl, ppizarror.com
%
% Manual template: [http://latex.ppizarror.com/Template-Informe/]
% Licencia MIT:    [https://opensource.org/licenses/MIT/]

\newcommand{\addimage}[3]{
	% Añade una imagen a un env 'image' con borde
	%	#1	Dirección de la imagen
	%	#2	Parámetros de la imagen
	%	#3	Leyenda de la imagen (opcional)
	\addimageboxed{#1}{#2}{0}{#3}
}

\newcommand{\addimageboxed}[4]{
	% Añade una imagen a un env 'image' con borde
	%	#1	Dirección de la imagen
	%	#2	Parámetros de la imagen
	%	#3	Ancho de la línea (en pt)
	%	#4	Leyenda de la imagen (opcional)
	\begingroup
		\setlength{\fboxsep}{0 pt}
		\setlength{\fboxrule}{#3 pt}
		\subfloat[#4]{
			\fbox{\includegraphics[#2]{#1}}
		}
		\hspace{0.1cm}
	\endgroup
}

\newcommand{\insertimage}[4][]{
	% Insertar una imagen
	% 	#1	Label (opcional)
	%	#2	Dirección de la imagen
	%	#3	Parámetros de la imagen
	%	#4	Leyenda de la imagen (opcional)
	\insertimageboxed[#1]{#2}{#3}{0}{#4}
}

\newcommand{\insertimageboxed}[5][]{
	% Insertar una imagen con recuadro
	% 	#1	Label (opcional)
	%	#2	Dirección de la imagen
	%	#3	Parámetros de la imagen
	%	#4	Ancho de la línea (en pt)
	%	#5	Leyenda de la imagen (opcional)
	\emptyvarerr{\insertimageboxed}{#2}{Direccion de la imagen no definida}
	\emptyvarerr{\insertimageboxed}{#3}{Parametros de la imagen no definidos}
	\emptyvarerr{\insertimageboxed}{#4}{Ancho de la linea no definido}
	\vspace{\margintopimages cm}
	\begin{figure}[H]
		\begingroup
			\setlength{\fboxsep}{0 pt}
			\setlength{\fboxrule}{#4 pt}
			\centering
			\fbox{\includegraphics[#3]{#2}}
		\endgroup
		\ifx\hfuzz#5\hfuzz
			\vspace{\captionlessmarginimage cm}
		\else
			\hspace{0cm}\caption{#5 #1}
		\fi
	\end{figure}
	\vspace{\marginbottomimages cm}
}

\newcommand{\insertdoubleimage}[8][]{
	% Insertar una imagen doble
	% 	#1	Label (opcional)
	%	#2	Dirección de la imagen 1
	%	#3	Parámetros de la imagen 1
	%	#4	Leyenda de la imagen 1
	%	#5	Dirección de la imagen 2
	%	#6	Parámetros de la imagen 2
	%	#7	Leyenda de la imagen 2
	%	#8	Leyenda de la imagen doble (opcional)
	\emptyvarerr{\insertdoubleimage}{#2}{Direccion de la imagen 1 no definida}
	\emptyvarerr{\insertdoubleimage}{#3}{Parametros de la imagen 1 no definidos}
	\emptyvarerr{\insertdoubleimage}{#5}{Direccion de la imagen 2 no definida}
	\emptyvarerr{\insertdoubleimage}{#6}{Parametros de la imagen 2 no definidos}
	\vspace{\margintopimages cm}
	\captionsetup{margin=0.45cm}
	\begin{figure}[H] \centering
		\subfloat[#4]{
			\includegraphics[#3]{#2}
		}
		\hspace{0.1cm}
		\subfloat[#7]{
			\includegraphics[#6]{#5}
		}
		\setcaptionmargincm{\captionlrmargin}
		\ifx\hfuzz#8\hfuzz
			\vspace{\captionlessmarginimage cm}
		\else
			\caption{#8 #1}
		\fi
	\end{figure}
	\setcaptionmargincm{\captionlrmargin}
	\vspace{\marginbottomimages cm}
}

\newcommand{\insertdoubleeqimage}[7][]{
	% Insertar una imagen doble, igual propiedades
	% 	#1	Label (opcional)
	%	#2	Dirección de la imagen 1
	%	#3	Leyenda de la imagen 1
	%	#4	Dirección de la imagen 2
	%	#5	Leyenda de la imagen 2
	%	#6	Propiedades de las imágenes
	%	#7 	Leyenda de la imagen doble (opcional)
	\insertdoubleimage[#1]{#2}{#6}{#3}{#4}{#6}{#5}{#7}
}

\newcommand{\inserttripleimage}[8][]{
	% Insertar una imagen triple
	% 	#1	Label (opcional)
	%	#2	Dirección de la imagen 1
	%	#3	Parámetros de la imagen 1
	%	#4	Dirección de la imagen 2
	%	#5	Parámetros de la imagen 2
	%	#6	Dirección de la imagen 3
	%	#7	Parámetros de la imagen 3
	%	#8	Leyenda de la imagen triple (opcional)
	\emptyvarerr{\inserttripleimage}{#2}{Direccion de la imagen 1 no definida}
	\emptyvarerr{\inserttripleimage}{#3}{Parametros de la imagen 1 no definidos}
	\emptyvarerr{\inserttripleimage}{#4}{Direccion de la imagen 2 no definida}
	\emptyvarerr{\inserttripleimage}{#5}{Parametros de la imagen 2 no definidos}
	\emptyvarerr{\inserttripleimage}{#6}{Direccion de la imagen 3 no definida}
	\emptyvarerr{\inserttripleimage}{#7}{Parametros de la imagen 3 no definidos}
	\vspace{\margintopimages cm}
	\captionsetup{margin=0.45cm}
	\begin{figure}[H] \centering
		\subfloat[]{
			\includegraphics[#3]{#2}
		}
		\hspace{0.1cm}
		\subfloat[]{
			\includegraphics[#5]{#4}
		}
		\hspace{0.1cm}
		\subfloat[]{
			\includegraphics[#7]{#6}
		}
		\setcaptionmargincm{\captionlrmargin}
		\ifx\hfuzz#8\hfuzz
			\vspace{\captionlessmarginimage cm}
		\else
			\caption{#8 #1}
		\fi
	\end{figure}
	\setcaptionmargincm{\captionlrmargin}
	\vspace{\marginbottomimages cm}
}

\newcommand{\inserttripleeqimage}[6][]{
	% Insertar una imagen triple, igual propiedades
	% 	#1	Label (opcional)
	%	#2	Dirección de la imagen 1
	%	#3	Dirección de la imagen 2
	%	#4	Dirección de la imagen 3
	%	#5	Propiedades de las imágenes
	%	#6	Leyenda de la imagen triple (opcional)
	\inserttripleimage[#1]{#2}{#5}{#3}{#5}{#4}{#5}{#6}
}

\newcommand{\insertquadimage}[7][]{
	% Insertar una imagen cuádruple, igual propiedades
	% 	#1	Label (opcional)
	%	#2	Dirección de la imagen 1
	%	#3	Dirección de la imagen 2
	%	#4	Dirección de la imagen 3
	%	#5	Dirección de la imagen 4
	%	#6	Propiedades de las imágenes
	%	#7	Leyenda de la imagen cuádruple (opcional)
	\emptyvarerr{\insertquadimage}{#2}{Direccion de la imagen 1 no definida}
	\emptyvarerr{\insertquadimage}{#3}{Direccion de la imagen 2 no definida}
	\emptyvarerr{\insertquadimage}{#4}{Direccion de la imagen 3 no definida}
	\emptyvarerr{\insertquadimage}{#5}{Direccion de la imagen 4 no definida}
	\emptyvarerr{\insertquadimage}{#6}{Propiedades de las imagenes no definidos}
	\vspace{\margintopimages cm}
	\captionsetup{margin=0.45cm}
	\begin{figure}[H] \centering
		\subfloat[]{
			\includegraphics[#6]{#2}
		}
		\hspace{0.1cm}
		\subfloat[]{
			\includegraphics[#6]{#3}
		}
		\hspace{0.1cm}
		\subfloat[]{
			\includegraphics[#6]{#4}
		}
		\hspace{0.1cm}
		\subfloat[]{
			\includegraphics[#6]{#5}
		}
		\setcaptionmargincm{\captionlrmargin}
		\ifx\hfuzz#7\hfuzz
			\vspace{\captionlessmarginimage cm}
		\else
			\caption{#7 #1}
		\fi
	\end{figure}
	\setcaptionmargincm{\captionlrmargin}
	\vspace{\marginbottomimages cm}
}

\newcommand{\insertpentaimage}[8][]{
	% Insertar una imagen quíntuple, igual propiedades
	% 	#1	Label (opcional)
	%	#2	Dirección de la imagen 1
	%	#3	Dirección de la imagen 2
	%	#4	Dirección de la imagen 3
	%	#5	Dirección de la imagen 4
	%	#6	Dirección de la imagen 5
	%	#7	Propiedades de las imágenes
	%	#8	Leyenda de la imagen quíntuple (opcional)
	\emptyvarerr{\insertpentaimage}{#2}{Direccion de la imagen 1 no definida}
	\emptyvarerr{\insertpentaimage}{#3}{Direccion de la imagen 2 no definida}
	\emptyvarerr{\insertpentaimage}{#4}{Direccion de la imagen 3 no definida}
	\emptyvarerr{\insertpentaimage}{#5}{Direccion de la imagen 4 no definida}
	\emptyvarerr{\insertpentaimage}{#6}{Direccion de la imagen 5 no definida}
	\emptyvarerr{\insertpentaimage}{#7}{Propiedades de las imagenes no definidas}
	\vspace{\margintopimages cm}
	\captionsetup{margin=0.45cm}
	\begin{figure}[H] \centering
		\subfloat[]{
			\includegraphics[#7]{#2}
		}
		\hspace{0.1cm}
		\subfloat[]{
			\includegraphics[#7]{#3}
		}
		\hspace{0.1cm}
		\subfloat[]{
			\includegraphics[#7]{#4}
		}
		\hspace{0.1cm}
		\subfloat[]{
			\includegraphics[#7]{#5}
		}
		\hspace{0.1cm}
		\subfloat[]{
			\includegraphics[#7]{#6}
		}
		\setcaptionmargincm{\captionlrmargin}
		\ifx\hfuzz#8\hfuzz
			\vspace{\captionlessmarginimage cm}
		\else
			\caption{#8 #1}
		\fi
	\end{figure}
	\setcaptionmargincm{\captionlrmargin}
	\vspace{\marginbottomimages cm}
}

\newcommand{\inserthexaimage}[9][]{
	% Insertar una imagen con 6 imágenes, igual propiedades
	% 	#1	Label (opcional)
	%	#2	Dirección de la imagen 1
	%	#3	Dirección de la imagen 2
	%	#4	Dirección de la imagen 3
	%	#5	Dirección de la imagen 4
	%	#6	Dirección de la imagen 5
	%	#7	Dirección de la imagen 6
	%	#8	Propiedades de las imágenes
	%	#9	Leyenda de la imagen global (opcional)
	\emptyvarerr{\inserthexaimage}{#2}{Direccion de la imagen 1 no definida}
	\emptyvarerr{\inserthexaimage}{#3}{Direccion de la imagen 2 no definida}
	\emptyvarerr{\inserthexaimage}{#4}{Direccion de la imagen 3 no definida}
	\emptyvarerr{\inserthexaimage}{#5}{Direccion de la imagen 4 no definida}
	\emptyvarerr{\inserthexaimage}{#6}{Direccion de la imagen 5 no definida}
	\emptyvarerr{\inserthexaimage}{#7}{Direccion de la imagen 6 no definida}
	\emptyvarerr{\inserthexaimage}{#8}{Propiedades de las imagenes no definidas}
	\vspace{\margintopimages cm}
	\captionsetup{margin=0.45cm}
	\begin{figure}[H] \centering
		\subfloat[]{
			\includegraphics[#8]{#2}
		}
		\hspace{0.1cm}
		\subfloat[]{
			\includegraphics[#8]{#3}
		}
		\hspace{0.1cm}
		\subfloat[]{
			\includegraphics[#8]{#4}
		}
		\hspace{0.1cm}
		\subfloat[]{
			\includegraphics[#8]{#5}
		}
		\hspace{0.1cm}
		\subfloat[]{
			\includegraphics[#8]{#6}
		}
		\hspace{0.1cm}
		\subfloat[]{
			\includegraphics[#8]{#7}
		}
		\setcaptionmargincm{\captionlrmargin}
		\ifx\hfuzz#9\hfuzz
			\vspace{\captionlessmarginimage cm}
		\else
			\caption{#9 #1}
		\fi
	\end{figure}
	\setcaptionmargincm{\captionlrmargin}
	\vspace{\marginbottomimages cm}
}

\newcommand{\insertimageleft}[4][]{
	% Insertar una imagen a la izquierda
	% 	#1	Label (opcional)
	%	#2	Dirección de la imagen
	%	#3	Ancho de la imagen (en textwidth)
	%	#4	Leyenda de la imagen (opcional)
	\insertimageleftboxed[#1]{#2}{#3}{0}{#4}
}

\newcommand{\insertimageleftboxed}[5][]{
	% Insertar una imagen a la izquierda
	% 	#1	Label (opcional)
	%	#2	Dirección de la imagen
	%	#3	Ancho de la imagen (en textwidth)
	%	#4	Ancho de la línea (en pt)
	%	#5	Leyenda de la imagen (opcional)
	\emptyvarerr{\insertimageleftboxed}{#2}{Direccion de la imagen no definida}
	\emptyvarerr{\insertimageleftboxed}{#3}{Ancho de la imagen no definido}
	\emptyvarerr{\insertimageleftboxed}{#4}{Ancho de la linea no definido}
	\begin{wrapfigure}{l}{#3\textwidth}
		\setcaptionmargincm{0}
		\ifthenelse{\equal{\figurecaptiontop}{true}}{}{
			\vspace{\marginfloatimages pt}}
		\begingroup
			\setlength{\fboxsep}{0 pt}
			\setlength{\fboxrule}{#4 pt}
			\centering
			\fbox{\includegraphics[width=\linewidth]{#2}}
		\endgroup
		\ifx\hfuzz#5\hfuzz
			\vspace{\captionlessmarginimage cm}
		\else
			\caption{#5 #1}
		\fi
	\end{wrapfigure}
	\setcaptionmargincm{\captionlrmargin}
}


\newcommand{\insertimageleftline}[5][]{
	% Insertar una imagen a la izquierda
	% 	#1	Label (opcional)
	%	#2	Dirección de la imagen
	%	#3	Ancho de la imagen (en textwidth)
	%	#4	Altura en líneas de la imagen
	%	#5	Leyenda de la imagen (opcional)
	\insertimageleftboxed[#1]{#2}{#3}{0}{#4}
}

\newcommand{\insertimageleftlineboxed}[6][]{
	% Insertar una imagen a la izquierda
	% 	#1	Label (opcional)
	%	#2	Dirección de la imagen
	%	#3	Ancho de la imagen (en textwidth)
	%	#4	Altura en líneas de la imagen
	%	#5	Ancho de la línea (en pt)
	%	#6	Leyenda de la imagen (opcional)
	\emptyvarerr{\insertimageleftlineboxed}{#2}{Direccion de la imagen no definida}
	\emptyvarerr{\insertimageleftlineboxed}{#3}{Ancho de la imagen no definido}
	\emptyvarerr{\insertimageleftlineboxed}{#4}{Altura en lineas de la imagen flotante izquierda no definida}
	\emptyvarerr{\insertimageleftlineboxed}{#5}{Ancho de la linea no definido}
	\begin{wrapfigure}[#4]{l}{#3\textwidth}
		\setcaptionmargincm{0}
		\ifthenelse{\equal{\figurecaptiontop}{true}}{}{
			\vspace{\marginfloatimages pt}}
		\begingroup
			\setlength{\fboxsep}{0 pt}
			\setlength{\fboxrule}{#5 pt}
			\centering
			\fbox{\includegraphics[width=\linewidth]{#2}}
		\endgroup
		\ifx\hfuzz#6\hfuzz
			\vspace{\captionlessmarginimage cm}
		\else
			\caption{#6 #1}
		\fi
	\end{wrapfigure}
	\setcaptionmargincm{\captionlrmargin}
}

\newcommand{\insertimageright}[4][]{
	% Insertar una imagen a la derecha
	% 	#1	Label (opcional)
	%	#2	Dirección de la imagen
	%	#3	Ancho de la imagen (en textwidth)
	%	#4	Leyenda de la imagen (opcional)
	\insertimagerightboxed[#1]{#2}{#3}{0}{#4}
}

\newcommand{\insertimagerightboxed}[5][]{
	% Insertar una imagen a la derecha
	% 	#1	Label (opcional)
	%	#2	Dirección de la imagen
	%	#3	Ancho de la imagen (en textwidth)
	%	#4	Ancho de la línea (en pt)
	%	#5	Leyenda de la imagen (opcional)
	\emptyvarerr{\insertimagerightboxed}{#2}{Direccion de la imagen no definida}
	\emptyvarerr{\insertimagerightboxed}{#3}{Ancho de la imagen no defindo}
	\emptyvarerr{\insertimagerightboxed}{#4}{Ancho de la linea no definido}
	\begin{wrapfigure}{r}{#3\textwidth}
		\setcaptionmargincm{0}
		\ifthenelse{\equal{\figurecaptiontop}{true}}{}{
			\vspace{\marginfloatimages pt}}
		\begingroup
			\setlength{\fboxsep}{0 pt}
			\setlength{\fboxrule}{#4 pt}
			\centering
			\fbox{\includegraphics[width=\linewidth]{#2}}
		\endgroup
		\ifx\hfuzz#5\hfuzz
			\vspace{\captionlessmarginimage cm}
		\else
			\caption{#5 #1}
		\fi
	\end{wrapfigure}
	\setcaptionmargincm{\captionlrmargin}
}

\newcommand{\insertimagerightline}[5][]{
	% Insertar una imagen a la derecha
	% 	#1	Label (opcional)
	%	#2	Dirección de la imagen
	%	#3	Ancho de la imagen (en textwidth)
	%	#4	Altura en líneas de la imagen
	%	#5	Leyenda de la imagen (opcional)
	\insertimagerightlineboxed[#1]{#2}{#3}{#4}{0}{#5}
}

\newcommand{\insertimagerightlineboxed}[6][]{
	% Insertar una imagen a la derecha
	% 	#1	Label (opcional)
	%	#2	Dirección de la imagen
	%	#3	Ancho de la imagen (en textwidth)
	%	#4	Altura en líneas de la imagen
	%	#5	Ancho de la línea (en pt)
	%	#6	Leyenda de la imagen (opcional)
	\emptyvarerr{\insertimagerightlineboxed}{#2}{Direccion de la imagen no definida}
	\emptyvarerr{\insertimagerightlineboxed}{#3}{Ancho de la imagen no defindo}
	\emptyvarerr{\insertimagerightlineboxed}{#4}{Altura en lineas de la imagen flotante derecha no definida}
	\emptyvarerr{\insertimagerightlineboxed}{#5}{Ancho de la linea no definido}
	\begin{wrapfigure}[#4]{r}{#3\textwidth}
		\setcaptionmargincm{0}
		\ifthenelse{\equal{\figurecaptiontop}{true}}{}{
			\vspace{\marginfloatimages pt}}
		\begingroup
			\setlength{\fboxsep}{0 pt}
			\setlength{\fboxrule}{#5 pt}
			\centering
			\fbox{\includegraphics[width=\linewidth]{#2}}
		\endgroup
		\ifx\hfuzz#6\hfuzz
			\vspace{\captionlessmarginimage cm}
		\else
			\caption{#6 #1}
		\fi
	\end{wrapfigure}
	\setcaptionmargincm{\captionlrmargin}
}

% Template:     Informe/Reporte LaTeX
% Documento:    Funciones para insertar títulos
% Versión:      5.3.0 (21/05/2018)
% Codificación: UTF-8
%
% Autor: Pablo Pizarro R. @ppizarror
%        Facultad de Ciencias Físicas y Matemáticas
%        Universidad de Chile
%        pablo.pizarro@ing.uchile.cl, ppizarror.com
%
% Manual template: [http://latex.ppizarror.com/Template-Informe/]
% Licencia MIT:    [https://opensource.org/licenses/MIT/]

% Insertar un título sin número
% 	#1	Título
\newcommand{\sectionanum}[1]{
	\emptyvarerr{\sectionanum}{#1}{Titulo no definido}
	\phantomsection
	\needspace{3\baselineskip}
	\section*{#1}
	\addcontentsline{toc}{section}{#1}
	\ifthenelse{\equal{\anumsecaddtocounter}{true}}{\stepcounter{section}}{}
	\changeheadertitle{#1}
}

% Insertar un título sin número y sin indexar
% 	#1	Título
\newcommand{\sectionanumnoi}[1]{
	\emptyvarerr{\sectionanumnoi}{#1}{Titulo no definido}
	\phantomsection
	\needspace{3\baselineskip}
	\section*{#1}
	\ifthenelse{\equal{\anumsecaddtocounter}{true}}{\stepcounter{section}}{}
	\changeheadertitle{#1}
}

% Insertar un título sin número sin cambiar el título del header
% 	#1	Título
\newcommand{\sectionanumheadless}[1]{
	\emptyvarerr{\sectionanumnoheadless}{#1}{Titulo no definido}
	\section*{#1}
	\addcontentsline{toc}{section}{#1}
	\ifthenelse{\equal{\anumsecaddtocounter}{true}}{\stepcounter{section}}{}
}

% Insertar un título sin número, sin indexar y sin cambiar el título del header
% 	#1	Título
\newcommand{\sectionanumnoiheadless}[1]{
	\emptyvarerr{\sectionanumnoi}{#1}{Titulo no definido}
	\section*{#1}
	\ifthenelse{\equal{\anumsecaddtocounter}{true}}{\stepcounter{section}}{}
}

% Insertar un subtítulo sin número
% 	#1	Subtítulo
\newcommand{\subsectionanum}[1]{
	\emptyvarerr{\subsectionanum}{#1}{Subtitulo no definido}
	\subsection*{#1}
	\addcontentsline{toc}{subsection}{#1}
	\ifthenelse{\equal{\anumsecaddtocounter}{true}}{\stepcounter{subsection}}{}
}

% Insertar un subtítulo sin número y sin indexar
% 	#1	Subtítulo
\newcommand{\subsectionanumnoi}[1]{
	\emptyvarerr{\subsectionanumnoi}{#1}{Subtitulo no definido}
	\subsection*{#1}
	\ifthenelse{\equal{\anumsecaddtocounter}{true}}{\stepcounter{subsection}}{}
}

% Insertar un sub-subtítulo sin número
% 	#1	Sub-subtítulo
\newcommand{\subsubsectionanum}[1]{
	\emptyvarerr{\subsubsectionanum}{#1}{Sub-subtitulo no definido}
	\subsubsection*{#1}
	\addcontentsline{toc}{subsubsection}{#1}
	\ifthenelse{\equal{\anumsecaddtocounter}{true}}{\stepcounter{subsubsection}}{}
}

% Insertar un sub-subtítulo sin número y sin indexar
% 	#1	Sub-subtítulo
\newcommand{\subsubsectionanumnoi}[1]{
	\emptyvarerr{\subsubsectionanumnoi}{#1}{Sub-subtitulo no definido}
	\subsubsection*{#1}
	\ifthenelse{\equal{\anumsecaddtocounter}{true}}{\stepcounter{subsubsection}}{}
}

% Cambia el título del header
%	#1	Título
\newcommand{\changeheadertitle}[1]{
	\emptyvarerr{\changeheadertitle}{#1}{Titulo no definido}
	\markboth{#1}{}
}

% Insertar un título en un índice
%	#1	Título
%	#2	Margen superior en pt. (opcional), 10pt por defecto
\newcommand{\insertindextitle}[2]{
	\emptyvarerr{\insertindextitle}{#1}{Titulo no definido}
	\ifx\hfuzz#2\hfuzz
		\addtocontents{toc}{\protect\addvspace{\indextitlemargin pt}}
	\else
		\addtocontents{toc}{\protect\addvspace{#2 pt}}
	\fi
	\addtocontents{toc}{\noindent\hyperref[swpn]{\textbf{#1}}}
}

% Crear un capítulo como una sección
%	#1	Título
\newcommand{\newchapter}[1]{
	\emptyvarerr{\newchapter}{#1}{Titulo no definido}
	\newpage
	\stepcounter{section}
	\phantomsection
	\needspace{3\baselineskip}
	\vspace* {3cm}
	\noindent {\huge{\textbf{\nomchapter\ \thesection}}} \\
	\vspace* {0.5cm} \\
	\noindent {\Huge{\textbf{#1}}} \\
	\vspace {0.5cm} \\
	\addcontentsline{toc}{section}{\protect\numberline{\thesection}#1}
	\markboth{#1}{}
}

% Template:     Informe/Reporte LaTeX
% Documento:    Funciones para insertar elementos
% Versión:      6.0.0 (13/10/2018)
% Codificación: UTF-8
%
% Autor: Pablo Pizarro R. @ppizarror
%        Facultad de Ciencias Físicas y Matemáticas
%        Universidad de Chile
%        pablo.pizarro@ing.uchile.cl, ppizarror.com
%
% Manual template: [http://latex.ppizarror.com/Template-Informe/]
% Licencia MIT:    [https://opensource.org/licenses/MIT/]

% Insertar párrafo
\newcommand{\newp}{
	\hbadness=10000 \vspace{\defaultnewlinesize pt} \par
}

% Insertar párrafo
% 	#1	Párrafo
\newcommand{\newpar}[1]{
	\hbadness=10000 #1 \newp
}

% Insertar párrafo sin nueva línea al final
% 	#1	Párrafo
\newcommand{\newparnl}[1]{
	#1 \par
}

% Redimensiona un ítem
% 	#1	Tamaño del nuevo objeto (En textwidth)
%	#2	Objeto a redimensionar
\newcommand{\itemresize}[2]{
	\emptyvarerr{\itemresize}{#1}{Tamano del nuevo objeto no definido}
	\emptyvarerr{\itemresize}{#2}{Objeto a redimensionar no definido}
	\resizebox{#1\textwidth}{!}{#2}
}

% Crea una página vacía
\newcommand{\insertemptypage}{
	\newpage
	\setcounter{templatepagecounter}{\thepage}
	\pagenumbering{gobble}
	\null
	\thispagestyle{empty}
	\newpage
	\pagenumbering{arabic}
	\setcounter{page}{\thetemplatepagecounter}
}

% Inserta un texto entre comillas
\newcommand{\quotes}[1]{\enquote*{#1}}

% Inserta un email con un link cliqueable
\newcommand{\insertemail}[1]{\href{mailto:#1}{\texttt{#1}}}

% END
% Template:     Informe/Reporte LaTeX
% Documento:    Definición de entornos
% Versión:      5.2.0 (21/05/2018)
% Codificación: UTF-8
%
% Autor: Pablo Pizarro R. @ppizarror
%        Facultad de Ciencias Físicas y Matemáticas
%        Universidad de Chile
%        pablo.pizarro@ing.uchile.cl, ppizarror.com
%
% Manual template: [http://latex.ppizarror.com/Template-Informe/]
% Licencia MIT:    [https://opensource.org/licenses/MIT/]

% Crea una sección de referencias solo para bibtex
\newenvironment{references}{
	\ifthenelse{\equal{\stylecitereferences}{bibtex}}{
	}{
		\throwerror{\references}{Solo se puede usar entorno references con estilo citas \noexpand\stylecitereferences=bibtex}
	}
	\begingroup
	% Se configura las referencias como una sección
	\ifthenelse{\equal{\donumrefsection}{true}}{
		\section{\namereferences}
	}{
		\sectionanum{\namereferences}
	}
	\renewcommand{\section}[2]{}
	\begin{thebibliography}{99}
	}
	{
	\end{thebibliography}
	\endgroup
}

% Crea una sección de anexos
\newenvironment{anexo}{
	\begingroup
	\clearpage
	\phantomsection
	\ifthenelse{\equal{\showappendixsectitle}{true}}{
		\appendixpage}{
	}
	\appendixtitleon
	\appendicestocpagenum
	\appendixtitletocon
	\bookmarksetup{
		numbered,
		openlevel=0
	}
	\begin{appendices}
		\bookmarksetupnext{level=part}
		\ifthenelse{\equal{\showappendixsecindex}{true}}{}{
			\belowpdfbookmark{\nameappendixsection}{contents}
		}
		\setcounter{secnumdepth}{3}
		\setcounter{tocdepth}{3}
		\ifthenelse{\equal{\appendixindepobjnum}{true}}{
			\counterwithin{equation}{section}
			\counterwithin{figure}{section}
			\counterwithin{lstlisting}{section}
			\counterwithin{table}{section}}{
		}
	}{
	\end{appendices}
	\bookmarksetupnext{level=0}
	\endgroup
}

% Inserta código fuente con parámetros
%	#1	Label (opcional)
%	#2	Estilo de código
%	#3	Parámetros
%	#4	Caption
\lstnewenvironment{sourcecodep}[4][]{
	\emptyvarerr{sourcecodep}{#2}{Estilo no definido}
	\checkvalidsourcecodestyle{#2}
	\ifthenelse{\equal{\showlinenumbers}{true}}{
		\rightlinenumbers}{
	}
	\ifthenelse{\equal{\codecaptiontop}{true}}{
		\ifx\hfuzz#4\hfuzz
			\ifx\hfuzz#3\hfuzz
				\lstset{
					style=#2
				}
			\else
				\lstset{
					style=#2,
					#3
				}
			\fi
		\else
			\ifx\hfuzz#3\hfuzz
				\lstset{
					caption={#4 #1},
					captionpos=t,
					style=#2
				}
			\else
				\lstset{
					caption={#4 #1},
					captionpos=t,
					style=#2,
					#3
				}
			\fi
		\fi
	}{
		\ifx\hfuzz#4\hfuzz
			\ifx\hfuzz#3\hfuzz
				\lstset{
					style=#2
				}
			\else
				\lstset{
					style=#2,
					#3
				}
			\fi
		\else
			\ifx\hfuzz#3\hfuzz
				\lstset{
					caption={#4 #1},
					captionpos=b,
					style=#2
				}
			\else
				\lstset{
					caption={#4 #1},
					captionpos=b,
					style=#2,
					#3
				}
			\fi
		\fi	
	}
}{
	\ifthenelse{\equal{\showlinenumbers}{true}}{
		\leftlinenumbers}{
	}
}

% Inserta código fuente sin parámetros
%	#1	Label (opcional)
%	#2	Estilo de código
%	#3	Caption
\lstnewenvironment{sourcecode}[3][]{
	\emptyvarerr{\equationresize}{#2}{Estilo no definido}
	\checkvalidsourcecodestyle{#2}
	\ifthenelse{\equal{\showlinenumbers}{true}}{
		\rightlinenumbers}{
	}
	\ifthenelse{\equal{\codecaptiontop}{true}}{
		\ifx\hfuzz#3\hfuzz
			\lstset{
				style=#2
			}
		\else
			\lstset{
				style=#2,
				caption={#3 #1},
				captionpos=t
			}
		\fi
	}{
		\ifx\hfuzz#3\hfuzz
			\lstset{
				style=#2
			}
		\else
			\lstset{
				style=#2,
				caption={#3 #1},
				captionpos=b
			}
		\fi
	}
}{
	\ifthenelse{\equal{\showlinenumbers}{true}}{
		\leftlinenumbers}{
	}
}

% Itemize en negrita
%	#1	Parámetros opcionales
\newenvironment{itemizebf}[1][]{
	\begin{itemize}[font=\bfseries,#1]
	}{
	\end{itemize}
}

% Enumerate en negrita
%	#1	Parámetros opcionales
\newenvironment{enumeratebf}[1][]{
	\begin{enumerate}[font=\bfseries,#1]
	}{
	\end{enumerate}
}

% Crea una sección de resumen
\newenvironment{resumen}{
	% Tipo de título para abstract
	\sectionfont{\color{\titlecolor} \fontsizetitle \styletitle \selectfont}
	
	% Inserta un título sin número, sin cabecera y sin aparecer en el índice, para que aparezca en el índice utilizar la función \sectionanumheadless
	\sectionanumnoiheadless{\nameabstract}}{
	
	% Salta de página si está imprimiendo por ambas caras
	\ifthenelse{\equal{\addemptypagetwosides}{true}}{
		\checkoddpage
		\ifoddpage
			\newpage
			\null
			\thispagestyle{empty}
			\newpage
			\addtocounter{page}{-1}
		\else
		\fi}{
	}
}

% Crea una sección de imágenes múltiples
%	#1	Label (opcional)
%	#2	Caption
\newenvironment{images}[2][]{
	\def\envimageslabelvar {#1}
	\def\envimagescaptionvar {#2}
	\def\envimagesinitialized {true}
	\vspace{\margintopimages cm}
	\captionsetup{margin=\captionmarginmultimg cm}
	\begin{figure}[H] \centering
		\vspace{-\marginrightmultimage cm}
		\vspace{-\marginrightmultimage cm}
		\vspace{-\marginrightmultimage cm}
		}{
		\setcaptionmargincm{\captionlrmargin}
		\ifx\hfuzz\envimagescaptionvar\hfuzz
			\vspace{\captionlessmarginimage cm}
		\else
			\caption{\envimagescaptionvar\envimageslabelvar}
		\fi
	\end{figure}
	\setcaptionmargincm{\captionlrmargin}
	\vspace{\marginbottomimages cm}
	\def\envimagesinitialized {false}
}
\def\envimagesinitialized {false}



% IMPORTACIÓN DE ESTILOS
% Template:     Informe/Reporte LaTeX
% Documento:    Estilos del template
% Versión:      5.2.1 (21/05/2018)
% Codificación: UTF-8
%
% Autor: Pablo Pizarro R. @ppizarror
%        Facultad de Ciencias Físicas y Matemáticas
%        Universidad de Chile
%        pablo.pizarro@ing.uchile.cl, ppizarror.com
%
% Manual template: [http://latex.ppizarror.com/Template-Informe/]
% Licencia MIT:    [https://opensource.org/licenses/MIT/]

% Template:     Informe/Reporte LaTeX
% Documento:    Definición de colores
% Versión:      4.7.6 (06/04/2018)
% Codificación: UTF-8
%
% Autor: Pablo Pizarro R.
%        Facultad de Ciencias Físicas y Matemáticas
%        Universidad de Chile
%        pablo.pizarro@ing.uchile.cl, ppizarror.com
%
% Manual template: [http://latex.ppizarror.com/Template-Informe/]
% Licencia MIT:    [https://opensource.org/licenses/MIT/]

\definecolor{backcolour}{rgb}{0.95, 0.95, 0.92}
\definecolor{codegray}{rgb}{0.5, 0.5, 0.5}
\definecolor{codegreen}{rgb}{0, 0.6, 0}
\definecolor{codepurple}{rgb}{0.58, 0, 0.82}
\definecolor{dgray}{rgb}{0.35, 0.35, 0.35}
\definecolor{dkblue}{rgb}{0, 0, 0.6}
\definecolor{dkgreen}{rgb}{0, 0.6, 0}
\definecolor{dkyellow}{cmyk}{0, 0, 0.8, 0.3}
\definecolor{editorGray}{rgb}{0.95, 0.95, 0.95}
\definecolor{editorGreen}{rgb}{0, 0.5, 0}
\definecolor{editorOcher}{rgb}{1, 0.5, 0}
\definecolor{gray}{rgb}{0.5, 0.5, 0.5}
\definecolor{lbrown}{RGB}{255, 252, 249}
\definecolor{lgray}{RGB}{180, 180, 180}
\definecolor{lyellow}{rgb}{1.0, 1.0, 0.88}
\definecolor{mauve}{rgb}{0.58, 0, 0.82}
\definecolor{mygray}{rgb}{0.5, 0.5, 0.5}
\definecolor{mygreen}{rgb}{0, 0.6, 0}
\definecolor{mylilas}{RGB}{170, 55, 241}
\definecolor{ocre}{RGB}{243, 102, 25}

% Template:     Informe/Reporte LaTeX
% Documento:    Estilos de código fuente
% Versión:      DEV
% Codificación: UTF-8
%
% Autor: Pablo Pizarro R. @ppizarror
%        Facultad de Ciencias Físicas y Matemáticas
%        Universidad de Chile
%        pablo.pizarro@ing.uchile.cl, ppizarror.com
%
% Manual template: [https://latex.ppizarror.com/Template-Informe/]
% Licencia MIT:    [https://opensource.org/licenses/MIT/]

% -----------------------------------------------------------------------------
% Bash
% -----------------------------------------------------------------------------
\lstdefinestyle{bash}{
	language=bash,
	breakatwhitespace=false,
	morecomment=[l]{rem},
	morecomment=[s]{::}{::},
	morekeywords={call,cp,dig,gcc,git,grep,ls,mv,python,rm,sudo,vim},
	sensitive=false
}

% -----------------------------------------------------------------------------
% C
% -----------------------------------------------------------------------------
\lstdefinestyle{c}{
	language=C,
	breakatwhitespace=false,
	keepspaces=true
}

% -----------------------------------------------------------------------------
% C++
% -----------------------------------------------------------------------------
\lstdefinestyle{cpp}{
	language=C++,
	breakatwhitespace=false,
	morecomment=[l][\color{magenta}]{\#}
}

% -----------------------------------------------------------------------------
% C#
% -----------------------------------------------------------------------------
\lstdefinestyle{csharp}{
	language=csh,
	morecomment=[l]{//},
	morecomment=[s]{/*}{*/},
	morekeywords={abstract,as,base,bool,break,byte,case,catch,char,checked,
		class,const,continue,decimal,default,delegate,do,double,else,enum,event,
		explicit,extern,false,finally,fixed,float,for,foreach,goto,if,implicit,
		in,int,interface,internal,is,lock,long,namespace,new,null,object,
		operator,out,override,params,private,protected,public,readonly,ref,
		return,sbyte,sealed,short,sizeof,stackalloc,static,string,struct,switch,
		this,throw,true,try,typeof,uint,ulong,unchecked,unsafe,ushort,using,
		virtual,void,volatile,while}
}

% -----------------------------------------------------------------------------
% CSS
% -----------------------------------------------------------------------------
\lstdefinelanguage{CSS}{
	morecomment=[s]{/*}{*/},
	morekeywords={-moz-binding,-moz-border-bottom-colors,
		-moz-border-left-colors,-moz-border-radius,
		-moz-border-radius-bottomleft,-moz-border-radius-bottomright,
		-moz-border-radius-topleft,-moz-border-radius-topright,
		-moz-border-right-colors,-moz-border-top-colors,-moz-opacity,
		-moz-outline,-moz-outline-color,-moz-outline-style,-moz-outline-width,
		-moz-user-focus,-moz-user-input,-moz-user-modify,-moz-user-select,
		-replace,-set-link-source,-use-link-source,accelerator,azimuth,
		background,background-attachment,background-color,background-image,
		background-position,background-position-x,background-position-y,
		background-repeat,behavior,border,border-bottom,border-bottom-color,
		border-bottom-style,border-bottom-width,border-collapse,border-color,
		border-left,border-left-color,border-left-style,border-left-width,
		border-right,border-right-color,border-right-style,border-right-width,
		border-spacing,border-style,border-top,border-top-color,
		border-top-style,border-top-width,border-width,bottom,caption-side,
		clear,clip,color,content,counter-increment,counter-reset,cue,cue-after,
		cue-before,cursor,direction,display,elevation,empty-cells,filter,float,
		font,font-family,font-size,font-size-adjust,font-stretch,font-style,
		font-variant,font-weight,height,ime-mode,include-source,
		layer-background-color,layer-background-image,layout-flow,layout-grid,
		layout-grid-char,layout-grid-char-spacing,layout-grid-line,
		layout-grid-mode,layout-grid-type,left,letter-spacing,line-break,
		line-height,list-style,list-style-image,list-style-position,
		list-style-type,margin,margin-bottom,margin-left,margin-right,
		margin-top,marker-offset,marks,max-height,max-width,min-height,
		min-width,orphans,outline,outline-color,outline-style,outline-width,
		overflow,overflow-X,overflow-Y,padding,padding-bottom,padding-left,
		padding-right,padding-top,page,page-break-after,page-break-before,
		page-break-inside,pause,pause-after,pause-before,pitch,pitch-range,
		play-during,position,quotes,richness,right,ruby-align,ruby-overhang,
		ruby-position,scrollbar-3d-light-color,scrollbar-arrow-color,
		scrollbar-base-color,scrollbar-dark-shadow-color,scrollbar-face-color,
		scrollbar-highlight-color,scrollbar-shadow-color,scrollbar-track-color,
		size,speak,speak-header,speak-numeral,speak-punctuation,speech-rate,
		stress,table-layout,text-align,text-align-last,text-autospace,
		text-decoration,text-indent,text-justify,text-kashida-space,
		text-overflow,text-shadow,text-transform,text-underline-position,top,
		unicode-bidi,vertical-align,visibility,voice-family,volume,white-space,
		widows,width,word-break,word-spacing,word-wrap,writing-mode,z-index,
		zoom},
	morestring=[s]{:}{;},
	sensitive=true
}

% -----------------------------------------------------------------------------
% CUDA
% -----------------------------------------------------------------------------
\lstdefinestyle{cuda}{
	language=C++,
	breakatwhitespace=false,
	emph={cudaFree,cudaMalloc,__device__,__global__,__host__,__shared__,
		__syncthreads},
	emphstyle=\color{dkcyan}\ttfamily,
	morecomment=[l][\color{magenta}]{\#},
	moredelim=[s][\ttfamily]{<<<}{>>>}
}

% -----------------------------------------------------------------------------
% DOCKER
% -----------------------------------------------------------------------------
\lstdefinelanguage{docker}{
	comment=[l]{\#},
	keywords={ADD,CMD,COPY,ENTRYPOINT,ENV,EXPOSE,FROM,LABEL,MAINTAINER,ONBUILD,
		RUN,STOPSIGNAL,USER,VOLUME,WORKDIR},
	morestring=[b]',
	morestring=[b]"
}
\lstdefinestyle{docker}{
	language=docker,
	breakatwhitespace=true
}

% -----------------------------------------------------------------------------
% GLSL - Shaders
% -----------------------------------------------------------------------------
\lstdefinelanguage{GLSL}{
	alsoletter={\#},
	morekeywords=[1]{
		attribute,bool,break,bvec2,bvec3,bvec4,case,centroid,const,continue,
		default,discard,do,else,false,flat,float,for,highp,if,in,inout,int,
		invariant,isampler1D,isampler1DArray,isampler2D,isampler2DArray,
		isampler2DMS,isampler2DMSArray,isampler2DRect,isampler3D,isamplerBuffer,
		isamplerCube,ivec2,ivec3,ivec4,layout,lowp,mat2,mat2x2,mat2x3,mat2x4,
		mat3,mat3x2,mat3x3,mat3x4,mat4,mat4x2,mat4x3,mat4x4,mediump,
		noperspective,out,precision,return,sampler1D,sampler1DArray,
		sampler1DArrayShadow,sampler1DShadow,sampler2D,sampler2DArray,
		sampler2DArrayShadow,sampler2DMS,sampler2DMSArray,sampler2DRect,
		sampler2DRectShadow,sampler2DShadow,sampler3D,samplerBuffer,
		samplerCube,samplerCubeShadow,smooth,struct,switch,true,uint,uniform,
		usampler1D,usampler1DArray,usampler2D,usampler2DArray,usampler2DMS,
		usampler2DMSArray,usampler2DRect,usampler3D,usamplerBuffer,usamplerCube,
		uvec2,uvec3,uvec4,varying,vec2,vec3,vec4,void,while
	},
	morekeywords=[2]{
		abs,acos,acosh,all,any,asin,asinh,atan,atan,atanh,ceil,clamp,cos,cosh,
		cross,degrees,determinant,dFdx,dFdy,distance,dot,EmitVertex,EndPrimitive,
		equal,exp,exp2,faceforward,floatBitsToInt,floatBitsToUint,floor,fract,
		fwidth,greaterThan,greaterThanEqual,intBitsToFloat,inverse,inversesqrt,
		isinf,isnan,length,lessThan,lessThanEqual,log,log2,matrixCompMult,max,
		min,mix,mod,modf,noise1,noise2,noise3,noise4,normalize,not,notEqual,
		outerProduct,pow,radians,reflect,refract,round,roundEven,shadow1D,
		shadow1DLod,shadow1DProj,shadow1DProjLod,shadow2D,shadow2DLod,
		shadow2DProj,shadow2DProjLod,sign,sin,sinh,smoothstep,sqrt,step,tan,tanh,
		texelFetch,texelFetchOffset,texture,texture1D,texture1DProj,
		texture1DProjLod,texture2D,texture2DLod,texture2DProj,texture2DProjLod,
		texture3D,texture3DLod,texture3DProj,texture3DProjLod,textureCube,
		textureCubeLod,textureGrad,textureGradOffset,textureLod,textureLodOffset,
		textureOffset,textureProj,textureProjGrad,textureProjGradOffset,
		textureProjLod,textureProjLodOffset,textureProjOffset,textureSize,
		transpose,trunc,uintBitsToFloat
	},
	morekeywords=[3]{
		\#version,core,gl_ClipDistance,gl_ClipDistance,gl_ClipVertex,gl_DepthRange,
		gl_FragColor,gl_FragCoord,gl_FragData,gl_FragDepth,gl_FrontFacing,
		gl_InstanceID,gl_Layer,gl_MaxClipDistances,gl_MaxCombinedTextureImageUnits,
		gl_MaxDrawBuffers,gl_MaxDrawBuffers,gl_MaxFragmentInputComponents,
		gl_MaxFragmentUniformComponents,gl_MaxGeometryInputComponents,
		gl_MaxGeometryOutputComponents,gl_MaxGeometryOutputVertices,
		gl_MaxGeometryOutputVertices,gl_MaxGeometryTextureImageUnits,
		gl_MaxGeometryTotalOutputComponents,gl_MaxGeometryUniformComponents,
		gl_MaxGeometryVaryingComponents,gl_MaxTextureImageUnits,
		gl_MaxVaryingComponents,gl_MaxVaryingFloats,gl_MaxVertexAttribs,
		gl_MaxVertexOutputComponents,gl_MaxVertexTextureImageUnits,
		gl_MaxVertexUniformComponents,gl_PerVertex,gl_PointCoord,gl_PointSize,
		gl_Position,gl_PrimitiveID,gl_VertexID
	},
	morecomment=[l]{//},
	morecomment=[s]{/*}{*/}
}
\lstdefinestyle{glsl}{
	language=GLSL,
	keywordstyle=[3]\color{dkcyan}\ttfamily,
	prebreak=\raisebox{0ex}[0ex][0ex]{\ensuremath{\hookleftarrow}},
	sensitive=true,
	upquote=true
}

% -----------------------------------------------------------------------------
% HTML5
% -----------------------------------------------------------------------------
\lstdefinelanguage{HTML5}{
	language=html,
	alsoletter={<>=-},
	morecomment=[s]{<!--}{-->},
	ndkeywords={-moz-transform:,<!,<body>,<canvas,</body>,</canvas>,</head>,
		</html>,</script>,</style>,</title>,<head>,<meta,<script>,<style>,
		<title>,=,><,border:,charset=,/>,height=,html>,id=,
		otherkeywords={<html>,transform:,transition-duration:,
			transition-property:,transition-timing-function:
		},width=},
	sensitive=true,
	tag=[s]
}
\lstdefinestyle{html5}{
	language=HTML5,
	alsodigit={.:;},
	alsolanguage=JavaScript,
	firstnumber=1,
	ndkeywordstyle=\color{dkgreen}\bfseries,
	numberfirstline=true
}

% -----------------------------------------------------------------------------
% Java
% -----------------------------------------------------------------------------
\lstdefinestyle{java}{
	language=Java,
	breakatwhitespace=true,
	keepspaces=true
}

% -----------------------------------------------------------------------------
% Javascript
% -----------------------------------------------------------------------------
\lstdefinelanguage{JavaScript}{
	comment=[l]{//},
	keepspaces=true,
	keywords={async,await,break,case,catch,catch,do,else,false,function,if,in,
		new,null,return,switch,true,typeof,var,while},
	morecomment=[s]{/*}{*/},
	morestring=[b]',
	morestring=[b]",
	ndkeywords={boolean,class,export,implements,import,this,throw},
	ndkeywordstyle=\color{darkgray}\bfseries,
	sensitive=false
}
\lstdefinestyle{js}{
	language=JavaScript
}

% -----------------------------------------------------------------------------
% JSON
% -----------------------------------------------------------------------------
\lstdefinestyle{json}{
	literate=*{0}{{{\color{numb}0}}}{1}{1}{{{\color{numb}1}}}{1}{2}
	{{{\color{numb}2}}}{1}{3}{{{\color{numb}3}}}{1}{4}{{{\color{numb}4}}}
	{1}{5}{{{\color{numb}5}}}{1}{6}{{{\color{numb}6}}}{1}{7}{{{\color{numb}7}}}
	{1}{8}{{{\color{numb}8}}}{1}{9}{{{\color{numb}9}}}{1}{:}
	{{{\color{punct}{:}}}}{1}{,}{{{\color{punct}{,}}}}{1}{\{}
	{{{\color{delim}{\{}}}}{1}{\}}{{{\color{delim}{\}}}}}
	{1}{[}{{{\color{delim}{[}}}}{1}{]}{{{\color{delim}{]}}}}{1},
	tabsize=2
}

% -----------------------------------------------------------------------------
% Kotlin
% -----------------------------------------------------------------------------
\lstdefinestyle{kotlin}{
	comment=[l]{//},
	emph={delegate,filter,first,firstOrNull,forEach,lazy,map,mapNotNull,println,
		return@},
	emphstyle={\color{OrangeRed}},
	keywords={abstract,actual,as,as?,break,by,class,companion,continue,data,do,
		dynamic,else,enum,expect,false,final,for,fun,get,if,import,in,interface,
		internal,is,null,object,override,package,private,public,return,set,
		super,suspend,this,throw,true,try,typealias,val,var,vararg,when,where,
		while},
	morecomment=[s]{/*}{*/},
	morestring=[b]",
	morestring=[s]{"""*}{*"""},
	ndkeywords={@Deprecated,@JvmField,@JvmName,@JvmOverloads,@JvmStatic,
		@JvmSynthetic,Array,Byte,Double,Float,Int,Integer,Iterable,Long,
		Runnable,Short,String},
	ndkeywordstyle=\color{BurntOrange}\bfseries,
	sensitive=true
}

% -----------------------------------------------------------------------------
% LaTeX
% -----------------------------------------------------------------------------
\lstdefinestyle{latex}{
	language=TeX,
	morekeywords={aacos,aasin,aatan,acos,addimage,addimageboxed,
		align,asin,atan,begin,bibitem,bibliography,bigstrut,boldmath,
		bookmarksetup,boxed,cancelto,caption,changeheadertitle,checkmark,
		checkvardefined,cite,dd,degree,eqref,equal,frac,fracnpartial,fullcite,
		hline,href,ifthenelse,imagesnewline,imageshspace,imagesvspace,
		includehfpdf,includefullhfpdf,insertalign,
		insertalignanum,insertaligncaptioned,insertaligncaptioned,
		insertaligncaptionedanum,insertaligned,insertalignedanum,
		insertalignedcaptioned,insertalignedcaptionedanum,insertemail,
		insertemptypage,inserteqimage,insertequation,insertequationanum,
		insertequationcaptioned,insertequationcaptionedanum,insertgather,
		insertgatheranum,insertgathercaptioned,insertgathercaptionedanum,
		insertgathered,insertgatheredanum,insertgatheredcaptioned,
		insertgatheredcaptionedanum,insertimage,insertimageleft,
		insertimageright,insertindextitle,insertindextitlepage,insertphone,
		isundefined,itemresize,label,LaTeX,lipsum,lpow,makeatletter,makeatother,
		newcommand,newcounter,newp,newpage,pow,quotes,ref,renewcommand,section,
		sectionanum,setcounter,setlength,shortcite,sourcecode,sourcecodep,
		subsection,subsectionanum,subsubsection,subsubsectionanum,
		subsubsubsection,subsubsubsection,subsubsubsectionanum,textbf,textit,
		textregistered,textsuperscript,texttt,throwbadconfig,unboldmath,url,
		xspace}
}

% -----------------------------------------------------------------------------
% Matlab
% -----------------------------------------------------------------------------
\lstdefinestyle{matlab}{
	language=Matlab,
	keepspaces=true,
	morekeywords={break,end,for,if,matlab2tikz,warning}
}

% -----------------------------------------------------------------------------
% OpenCL
% -----------------------------------------------------------------------------
\lstdefinestyle{opencl}{
	language=C++,
	breakatwhitespace=false,
	emph={bool3,bool4,bool8,bool16,char2,char3,char4,char8,char16,
		complex,constant,event_t,bool2,float2,float3,float4,float8,float16,
		global,half2,half3,half4,half8,half16,image2d_t,image3d_t,imaginary,
		int2,int3,int4,int8,int16,kernel,local,long2,long3,long4,long8,long16,
		private,quad,quad2,quad3,quad4,quad8,quad16,sampler_t,short2,short3,
		short4,short8,short16,uchar2,uchar3,uchar4,uchar8,uchar16,uint2,uint3,
		uint4,uint8,uint16,ulong2,ulong3,ulong4,ulong8,ulong16,ushort2,ushort3,
		ushort4,ushort8,ushort16,__constant,__global,__kernel,__local,__private},
	emphstyle=\color{dkcyan}\ttfamily,
	morecomment=[l][\color{magenta}]{\#}
}

% -----------------------------------------------------------------------------
% OpenSees
% -----------------------------------------------------------------------------
\lstdefinestyle{opensees}{
	language=tcl,
	breakatwhitespace=false,
	emph=[1]{-dof,-ele,-eleRange,-file,-iNode,-jNode,-kNode,-ndf,-ndm,-node,
		-nodeRange,-region,-time},
	emphstyle=[1]\color{black}\bfseries\em,
	keepspaces=true,
	morecomment=[l]{\#},
	ndkeywords={Aggregator,algorithm,analysis,analyze,BasicBuilder,constraints,
		deformation,disp,DisplayModel2D,DisplayModel3D,Elastic,element,equalDOF,fix,
		fixX,fixY,fixZ,geomTransf,integrator,Linear,loadConst,LoadControl,mass,
		model,Node,node,nonlinearBeamColumn,numberer,pattern,Plain,printA,
		rayleigh,reaction,recorder,region,section,Steel01,system,test,truss,
		uniaxialMaterial,wipe},
	ndkeywordstyle=\color{dkcyan}\ttfamily
}

% -----------------------------------------------------------------------------
% Perl
% -----------------------------------------------------------------------------
\lstdefinestyle{perl}{
	language=Perl,
	alsoletter={\%},
	breakatwhitespace=false,
	keepspaces=true
}

% -----------------------------------------------------------------------------
% PHP
% -----------------------------------------------------------------------------
\lstdefinestyle{php}{
	language=php,
	emph=[1]{php},
	emph=[2]{if,and,or,else},
	emph=[3]{abstract,as,const,else,elseif,endfor,endforeach,endif,extends,
		final,for,foreach,global,if,implements,private,protected,public,static,
		var},
	emphstyle=[1]\color{black},
	emphstyle=[2]\color{dkyellow},
	keywords={abstract,and,array,as,break,callable,case,catch,class,clone,const,
		continue,declare,default,die,do,echo,else,elseif,empty,enddeclare,
		endfor,endforeach,endif,endswitch,endwhile,eval,exit,extends,final,
		finally,for,foreach,function,global,goto,if,implements,include,
		include_once,instanceof,insteadof,interface,isset,list,namespace,new,or,
		print,private,protected,public,require,require_once,return,static,
		switch,throw,trait,try,unset,use,var,while,xor,yield,__halt_compiler},
	showlines=true,
	upquote=true
}

% -----------------------------------------------------------------------------
% Texto plano
% -----------------------------------------------------------------------------
\lstdefinestyle{plaintext}{
	language={},
	keepspaces=true,
	literate={\ \ }{{\ }}1,
	tabsize=4
}

% -----------------------------------------------------------------------------
% Pseudocódigo
% -----------------------------------------------------------------------------
\lstdefinestyle{pseudocode}{
	language=Matlab,
	backgroundcolor=\color{white},
	breakatwhitespace=false,
	commentstyle=\color{gray}\upshape,
	frame=tb,
	keepspaces=true,
	keywords={and,be,begin,break,datatype,do,elif,else,end,for,foreach,fun,
		function,if,in,input,let,null,or,output,pop,procedure,push,repeat,
		return,swap,until,while,xor},
	keywordstyle=\color{black}\bfseries\em,
	mathescape=true,
	morecomment=[l]{//},
	morecomment=[l]{\#},
	morecomment=[s]{/*}{*/},
	morecomment=[s]{/**}{*/},
	sensitive=false,
	stringstyle=\color{dkgray}\bfseries\em
}

% -----------------------------------------------------------------------------
% Python
% -----------------------------------------------------------------------------
\lstdefinestyle{python}{
	language=Python,
	breakatwhitespace=false,
	keepspaces=true
}

% -----------------------------------------------------------------------------
% Ruby
% -----------------------------------------------------------------------------
\lstdefinestyle{ruby} {
	language=Ruby,
	breakatwhitespace=true,
	morestring=[s][]{\#\{}{\}},
	morestring=*[d]{"},
	sensitive=true
}

% -----------------------------------------------------------------------------
% Scala
% -----------------------------------------------------------------------------
\lstdefinestyle{scala}{
	language=scala,
	breakatwhitespace=true,
	morecomment=[l]{//},
	morecomment=[n]{/*}{*/},
	morekeywords={abstract,case,catch,class,def,do,else,extends,false,final,
		finally,for,if,implicit,import,match,mixin,new,null,object,override,
		package,private,protected,requires,return,sealed,super,this,throw,trait,
		true,try,type,val,var,while,with,yield},
	morestring=[b]',
	morestring=[b]",
	morestring=[b]""",
	otherkeywords={=>,<-,<\%,<:,>:,\#,@}
}

% -----------------------------------------------------------------------------
% SQL
% -----------------------------------------------------------------------------
\lstdefinestyle{sql}{
	language=SQL,
	breakatwhitespace=true
}

% -----------------------------------------------------------------------------
% TCL
% -----------------------------------------------------------------------------
\lstdefinestyle{tcl}{
	language=tcl,
	breakatwhitespace=false,
	keepspaces=true,
	morecomment=[l]{\#}
}

% -----------------------------------------------------------------------------
% XML
% -----------------------------------------------------------------------------
\lstdefinelanguage{XML}{
	morecomment=[s]{<?}{?>},
	morekeywords={encoding,type,version,xmlns},
	morestring=[b]",
	morestring=[s]{>}{<}
}
\lstdefinestyle{xml}{
	language=XML,
	tabsize=2
}

% -----------------------------------------------------------------------------
% Configuración de códigos fuente
% -----------------------------------------------------------------------------
\lstset{
	% \mbox{\textcolor{black}{$\hookrightarrow$}\space}
	aboveskip=3mm,
	basicstyle={\small\ttfamily},
	belowskip=3mm,
	breaklines=true,
	columns=fullflexible,
	commentstyle=\color{dkgreen}\ttfamily\upshape,
	extendedchars=true,
	identifierstyle=\color{black},
	keepspaces=true,
	keywordstyle=\color{blue},
	literate={á}{{\'a}}1 {é}{{\'e}}1 {í}{{\'i}}1 {ó}{{\'o}}1 {ú}{{\'u}}1
		{Á}{{\'A}}1 {É}{{\'E}}1 {Í}{{\'I}}1 {Ó}{{\'O}}1 {Ú}{{\'U}}1 {à}{{\`a}}1
		{è}{{\`e}}1 {ì}{{\`i}}1 {ò}{{\`o}}1 {ù}{{\`u}}1 {À}{{\`A}}1 {È}{{\'E}}1
		{Ì}{{\`I}}1 {Ò}{{\`O}}1 {Ù}{{\`U}}1 {ä}{{\"a}}1 {ë}{{\"e}}1 {ï}{{\"i}}1
		{ö}{{\"o}}1 {ü}{{\"u}}1 {Ä}{{\"A}}1 {Ë}{{\"E}}1 {Ï}{{\"I}}1 {Ö}{{\"O}}1
		{Ü}{{\"U}}1 {â}{{\^a}}1 {ê}{{\^e}}1 {î}{{\^i}}1 {ô}{{\^o}}1 {û}{{\^u}}1
		{Â}{{\^A}}1 {Ê}{{\^E}}1 {Î}{{\^I}}1 {Ô}{{\^O}}1 {Û}{{\^U}}1 {œ}{{\oe}}1
		{Œ}{{\OE}}1 {æ}{{\ae}}1 {Æ}{{\AE}}1 {ß}{{\ss}}1 {ű}{{\H{u}}}1
		{Ű}{{\H{U}}}1 {ő}{{\H{o}}}1 {Ő}{{\H{O}}}1 {ç}{{\c c}}1 {Ç}{{\c C}}1
		{ø}{{\o}}1 {å}{{\r a}}1 {Å}{{\r A}}1 {€}{{\EUR}}1 {£}{{\pounds}}1
		{ñ}{{\~n}}1 {Ñ}{{\~N}}1 {¿}{{?``}}1 {¡}{{!``}}1,
	numbers=left,
	numbersep=5pt,
	numberstyle=\tiny\color{dkgray},
	postbreak=\mbox{$\hookrightarrow$\space},
	showspaces=false,
	showstringspaces=false,
	showtabs=false,
	stepnumber=1,
	stringstyle=\color{mauve},
	tabsize=3
}

% -----------------------------------------------------------------------------
% Chequeo de estilos, cualquier nuevo estilo añadirlo a esta lista
% -----------------------------------------------------------------------------
\newcommand{\checkvalidsourcecodestyle}[1]{
	\ifthenelse{\equal{#1}{bash}}{}{
	\ifthenelse{\equal{#1}{c}}{}{
	\ifthenelse{\equal{#1}{cpp}}{}{
	\ifthenelse{\equal{#1}{csharp}}{}{
	\ifthenelse{\equal{#1}{cuda}}{}{
	\ifthenelse{\equal{#1}{docker}}{}{
	\ifthenelse{\equal{#1}{glsl}}{}{
	\ifthenelse{\equal{#1}{html5}}{}{
	\ifthenelse{\equal{#1}{java}}{}{
	\ifthenelse{\equal{#1}{js}}{}{
	\ifthenelse{\equal{#1}{json}}{}{
	\ifthenelse{\equal{#1}{kotlin}}{}{
	\ifthenelse{\equal{#1}{latex}}{}{
	\ifthenelse{\equal{#1}{matlab}}{}{
	\ifthenelse{\equal{#1}{opencl}}{}{
	\ifthenelse{\equal{#1}{opensees}}{}{
	\ifthenelse{\equal{#1}{perl}}{}{
	\ifthenelse{\equal{#1}{php}}{}{
	\ifthenelse{\equal{#1}{plaintext}}{}{
	\ifthenelse{\equal{#1}{pseudocode}}{}{
	\ifthenelse{\equal{#1}{python}}{}{
	\ifthenelse{\equal{#1}{ruby}}{}{
	\ifthenelse{\equal{#1}{scala}}{}{
	\ifthenelse{\equal{#1}{sql}}{}{
	\ifthenelse{\equal{#1}{tcl}}{}{
	\ifthenelse{\equal{#1}{xml}}{}{
		\def\sourcecodestyle {#1}
		\throwbadconfig{Estilo de codigo desconocido}{\sourcecodestyle}{bash,c,cpp,
		csharp,cuda,docker,glsl,html5,java,js,json,kotlin,latex,matlab,opencl,
		opensees,perl,php,plaintext,pseudocode,python,ruby,scala,sql,tcl,xml}
	}}}}}}}}}}}}}}}}}}}}}}}}}}
}

% END
% Template:     Informe/Reporte LaTeX
% Documento:    Otros estilos
% Versión:      6.1.0 (03/11/2018)
% Codificación: UTF-8
%
% Autor: Pablo Pizarro R. @ppizarror
%        Facultad de Ciencias Físicas y Matemáticas
%        Universidad de Chile
%        pablo.pizarro@ing.uchile.cl, ppizarror.com
%
% Manual template: [https://latex.ppizarror.com/Template-Informe/]
% Licencia MIT:    [https://opensource.org/licenses/MIT/]

% -----------------------------------------------------------------------------
% Estilo de enumeración en griego
% -----------------------------------------------------------------------------
\RequirePackage{enumitem}
\makeatletter
\def\greek#1{\expandafter\@greek\csname c@#1\endcsname}
\def\Greek#1{\expandafter\@Greek\csname c@#1\endcsname}
\def\@greek#1{
	\ifcase#1
		\or $\alpha$
		\or $\beta$
		\or $\gamma$
		\or $\delta$
		\or $\epsilon$
		\or $\zeta$
		\or $\eta$
		\or $\theta$
		\or $\iota$
		\or $\kappa$
		\or $\lambda$
		\or $\mu$
		\or $\nu$
		\or $\xi$
		\or $o$
		\or $\pi$
		\or $\rho$
		\or $\sigma$
		\or $\tau$
		\or $\upsilon$
		\or $\phi$
		\or $\chi$
		\or $\psi$
		\or $\omega$
	\fi
}
\def\@Greek#1{
	\ifcase#1
		\or $\mathrm{A}$
		\or $\mathrm{B}$
		\or $\Gamma$
		\or $\Delta$
		\or $\mathrm{E}$
		\or $\mathrm{Z}$
		\or $\mathrm{H}$
		\or $\Theta$
		\or $\mathrm{I}$
		\or $\mathrm{K}$
		\or $\Lambda$
		\or $\mathrm{M}$
		\or $\mathrm{N}$
		\or $\Xi$
		\or $\mathrm{O}$
		\or $\Pi$
		\or $\mathrm{P}$
		\or $\Sigma$
		\or $\mathrm{T}$
		\or $\mathrm{Y}$
		\or $\Phi$
		\or $\mathrm{X}$
		\or $\Psi$
		\or $\Omega$
	\fi
}
\makeatother
\AddEnumerateCounter{\greek}{\@greek}{24}
\AddEnumerateCounter{\Greek}{\@Greek}{12}

% -----------------------------------------------------------------------------
% Columna centrada en tablas
% -----------------------------------------------------------------------------
\newcolumntype{P}[1]{
	>{\centering\arraybackslash}p{#1}
}

% END


% CONFIGURACIÓN INICIAL DEL DOCUMENTO
% Template:     Informe/Reporte LaTeX
% Documento:    Configuración inicial del template
% Versión:      5.5.1 (04/09/2018)
% Codificación: UTF-8
%
% Autor: Pablo Pizarro R. @ppizarror
%        Facultad de Ciencias Físicas y Matemáticas
%        Universidad de Chile
%        pablo.pizarro@ing.uchile.cl, ppizarror.com
%
% Manual template: [http://latex.ppizarror.com/Template-Informe/]
% Licencia MIT:    [https://opensource.org/licenses/MIT/]

% Se revisa si las variables no han sido borradas
\checkvardefined{\titulodelinforme}
\checkvardefined{\temaatratar}
\checkvardefined{\autordeldocumento}
\checkvardefined{\nombredelcurso}
\checkvardefined{\codigodelcurso}
\checkvardefined{\nombreuniversidad}
\checkvardefined{\nombrefacultad}
\checkvardefined{\departamentouniversidad}
\checkvardefined{\localizacionuniversidad}

% Se añade \xspace a las variables
\makeatletter
	\g@addto@macro\titulodelinforme\xspace
	\g@addto@macro\temaatratar\xspace
	\g@addto@macro\autordeldocumento\xspace
	\g@addto@macro\nombredelcurso\xspace
	\g@addto@macro\codigodelcurso\xspace
	\g@addto@macro\nombreuniversidad\xspace
	\g@addto@macro\nombrefacultad\xspace
	\g@addto@macro\departamentouniversidad\xspace
	\g@addto@macro\localizacionuniversidad\xspace
\makeatother

% Se crean variables si se borraron
\ifthenelse{\isundefined{\tablaintegrantes}}{
	\errmessage{LaTeX Warning: Se borro la variable \noexpand\tablaintegrantes, creando una vacia}
	\def\tablaintegrantes {}}{
}

% Se activan números en menú marcadores del pdf
\ifthenelse{\equal{\cfgpdfsecnumbookmarks}{true}}{
	\bookmarksetup{numbered}}{
}

% Se define metadata del pdf
\ifthenelse{\equal{\cfgshowbookmarkmenu}{true}}{
	\def\cdfpagemodepdf {UseOutlines}
	}{
	\def\cdfpagemodepdf {UseNone}
}
\hypersetup{
	bookmarksopen={\cfgpdfbookmarkopen},
	bookmarksopenlevel={\cfgbookmarksopenlevel},
	bookmarkstype={toc},
	pdfauthor={\autordeldocumento},
	pdfcenterwindow={\cfgpdfcenterwindow},
	pdfcopyright={\cfgpdfcopyright},
	pdfcreator={LaTeX},
	pdfdisplaydoctitle={\cfgpdfdisplaydoctitle},
	pdffitwindow={\cfgpdffitwindow},
	pdfinfo={
		Curso.Codigo={\codigodelcurso},
		Curso.Nombre={\nombredelcurso},
		Documento.Autor={\autordeldocumento},
		Documento.Tema={\temaatratar},
		Documento.Titulo={\titulodelinforme},
		Template.Autor.Alias={ppizarror},
		Template.Autor.Email={pablo.pizarro@ing.uchile.cl},
		Template.Autor.Nombre={Pablo Pizarro R.},
		Template.Autor.Web={http://ppizarror.com/},
		Template.Codificacion={UTF-8},
		Template.Fecha={04/09/2018},
		Template.Latex.Compilador={pdflatex},
		Template.Licencia.Tipo={MIT},
		Template.Licencia.Web={https://opensource.org/licenses/MIT/},
		Template.Nombre={Template-Informe},
		Template.Tipo={Normal},
		Template.Version.Dev={5.5.1-N},
		Template.Version.Hash={2BCC8DCE7F3F001260ED33CAF94B9672},
		Template.Version.Release={5.5.1},
		Template.Web.Dev={https://github.com/Template-Latex/Template-Informe/},
		Template.Web.Manual={http://latex.ppizarror.com/Template-Informe/},
		Universidad.Departamento={\departamentouniversidad},
		Universidad.Nombre={\nombreuniversidad},
		Universidad.Ubicacion={\localizacionuniversidad}
	},
	pdfkeywords={\codigodelcurso, \nombredelcurso, \nombreuniversidad, \localizacionuniversidad},
	pdflang={\documentlang},
	pdfmenubar={\cfgpdfmenubar},
	pdfpagelayout={\cfgpdfpagemode},
	pdfpagemode={\cdfpagemodepdf},
	pdfproducer={Template-Informe v5.5.1 | (Pablo Pizarro R.) ppizarror.com},
	pdfremotestartview={Fit},
	pdfstartpage={1},
	pdfstartview={\cfgpdfpageview},
	pdfsubject={\temaatratar},
	pdftitle={\titulodelinforme},
	pdftoolbar={\cfgpdftoolbar},
	pdfview={\cfgpdfpageview}
}

% Establece la carpeta de imágenes por defecto
\graphicspath{{./\defaultimagefolder}}

% Definición de valores e dimensiones
\newcounter{templatepagecounter} % Administra números de páginas
\renewcommand{\baselinestretch}{\defaultinterline} % Ajuste del entrelineado
\setlength{\headheight}{64 pt} % Tamaño de la cabecera sin fancyhdr
\setcounter{MaxMatrixCols}{20} % Número máximo de columnas en matrices
\setlength{\footnotemargin}{4 mm} % Margen del footnote
\setlength{\columnsep}{\columnsepwidth em} % Separación entre columnas
\ifthenelse{\equal{\showlinenumbers}{true}}{
	\setlength{\linenumbersep}{0.22cm}
	\renewcommand\linenumberfont{\normalfont\tiny\color{\linenumbercolor}}
	}{
}

% Posición inicial de los objetos
\floatplacement{figure}{\imagedefaultplacement}
\floatplacement{table}{\tabledefaultplacement}
\floatplacement{tikz}{\tikzdefaultplacement}

% Configuración de los colores
\color{\maintextcolor} % Color principal
\arrayrulecolor{\tablelinecolor} % Color de las líneas de las tablas
\sethlcolor{\highlightcolor} % Color del subrayado por defecto
\ifthenelse{\equal{\showborderonlinks}{true}}{
	% Color de links con borde
	\hypersetup{
		citebordercolor=\numcitecolor,
		linkbordercolor=\linkcolor,
		urlbordercolor=\urlcolor
	}
}{
	% Color de links sin borde
	\hypersetup{
		hidelinks,
		colorlinks=true,
		citecolor=\numcitecolor,
		linkcolor=\linkcolor,
		urlcolor=\urlcolor
	}
}
\ifthenelse{\equal{\colorpage}{white}}{}{
	\pagecolor{\colorpage}
}

% Configuración de las leyendas
\setcaptionmargincm{\captionlrmargin} % Márgenes de las leyendas por defecto
\ifthenelse{\equal{\captiontextbold}{true}}{ % Texto en negrita en etiquetas
	\renewcommand{\captiontextbold}{bf}}{
	\renewcommand{\captiontextbold}{}
}
\captionsetup{ % Se configura el texto de los caption
	labelfont={color=\captioncolor, \captiontextbold},
	textfont={color=\captiontextcolor},
	singlelinecheck=on
}
\floatsetup[figure]{ % Configuración de márgenes en las figuras
	captionskip=\captiontbmarginfigure pt
}
\floatsetup[table]{ % Configuración de márgenes en las tablas
	captionskip=\captiontbmargintable pt
}
\ifthenelse{\equal{\figurecaptiontop}{true}}{ % Caption superior en figuras
	\floatsetup[figure]{position=above}}{
}
\ifthenelse{\equal{\tablecaptiontop}{true}}{ % Caption superior en tablas
	\floatsetup[table]{position=top}
	}{
	\floatsetup[table]{position=bottom}
}
\ifthenelse{\equal{\captionalignment}{justified}}{ % Leyenda justificada
	\captionsetup{
		justification=justified,
		format=plain
	}
}{
\ifthenelse{\equal{\captionalignment}{centered}}{ % Leyenda centrada
	\captionsetup{
		justification=centering
	}
}{
\ifthenelse{\equal{\captionalignment}{left}}{ % Leyenda alineada a la izquierda
	\captionsetup{
		justification=raggedright,
		singlelinecheck=false
	}
}{
\ifthenelse{\equal{\captionalignment}{right}}{ % Leyenda alineada a la derecha
	\captionsetup{
		justification=raggedleft,
		singlelinecheck=false
	}
}{
	\throwbadconfig{Posicion de leyendas desconocida}{\captionalignment}{justified,centered,left,right}}}}
}

% Configuración de referencias y citas
\ifthenelse{\equal{\stylecitereferences}{natbib}}{
	\bibliographystyle{apa}
	\setlength{\bibsep}{\natbibrefsep pt}
}{
\ifthenelse{\equal{\stylecitereferences}{apacite}}{
	\bibliographystyle{apacite}
	\setlength{\bibitemsep}{\apaciterefsep pt}
}{
\ifthenelse{\equal{\stylecitereferences}{bibtex}}{
	\bibliographystyle{apa}
	\newlength{\bibitemsep}
	\setlength{\bibitemsep}{.2\baselineskip plus .05\baselineskip minus .05\baselineskip}
	\newlength{\bibparskip}\setlength{\bibparskip}{0pt}
	\let\oldthebibliography\thebibliography
	\renewcommand\thebibliography[1]{
		\oldthebibliography{#1}
		\setlength{\parskip}{\bibitemsep}
		\setlength{\itemsep}{\bibparskip}
	}
	\setlength{\bibitemsep}{\bibtexrefsep pt}
}{
	\throwbadconfig{Estilo citas desconocido}{\stylecitereferences}{bibtex,apacite,natbib}}}
}
\makeatletter
\ifthenelse{\equal{\twocolumnreferences}{true}}{ % Referencias en 2 columnas
	\renewenvironment{thebibliography}[1]
	{\begin{multicols}{2}[\section*{\refname}]
		\@mkboth{\MakeUppercase\refname}{\MakeUppercase\refname}
		\list{\@biblabel{\@arabic\c@enumiv}}
		{\settowidth\labelwidth{\@biblabel{#1}}
			\leftmargin\labelwidth
			\advance\leftmargin\labelsep
			\@openbib@code
			\usecounter{enumiv}
			\let\p@enumiv\@empty
			\renewcommand\theenumiv{\@arabic\c@enumiv}}
		\sloppy
		\clubpenalty 4000
		\@clubpenalty \clubpenalty
		\widowpenalty 4000
		\sfcode`\.\@m}
		{\def\@noitemerr
		{\@latex@warning{Ambiente `thebibliography' no definido}}
		\endlist\end{multicols}}}{}
\makeatother

% Configuración anexo
\patchcmd{\appendices}{\quad}{\sectionappendixlastchar\quad}{}{}

% Se añade listings (código fuente) a tocloft
\begingroup
	\makeatletter
	\let\newcounter\@gobble\let\setcounter\@gobbletwo
	\globaldefs\@ne\let\c@loldepth\@ne
	\newlistof{listings}{lol}{\lstlistlistingname}
	\newlistentry{lstlisting}{lol}{0}
	\makeatother
\endgroup

% Reconfiguración de tamaño de páginas
\makeatletter
	\def\ifGm@preamble#1{\@firstofone}
	\appto\restoregeometry{
		\pdfpagewidth=\paperwidth
		\pdfpageheight=\paperheight}
	\apptocmd\newgeometry{
		\pdfpagewidth=\paperwidth
		\pdfpageheight=\paperheight}{}{}
\makeatother

% Configuración de hbox y vbox
\hfuzz=200pt
\vfuzz=200pt
\hbadness=\maxdimen
\vbadness=\maxdimen

% Se activa el modo estricto de revisión de números de página
\strictpagecheck

% Actualización márgen títulos
\titlespacing{\section}{0pt}{20pt}{10pt}
\titlespacing{\subsection}{0pt}{15pt}{10pt}

% Se activa el word-wrap para textos con \texttt{}
\ttfamily \hyphenchar\the\font=`\-

% Se define el tipo de texto de los url
\urlstyle{tt}

% Se revisa si se importa tikz
\ifthenelse{\equal{\portraitstyle}{style16}}{\coreimporttikz}{}
\ifthenelse{\equal{\portraitstyle}{\bgtemplatetestcode}}{\coreimporttikz}{}

% Se define la versión menor a compilar
\pdfminorversion=\pdfcompileversion

% Profundidad de las secciones
\setcounter{secnumdepth}{4}

% Crea las subsubsubsecciones
\titleclass{\subsubsubsection}{straight}[\subsection]

% Crea contador subsubsubsección
\newcounter{subsubsubsection}[subsubsection]

% Establece el número de las subsubsubsecciones
\ifthenelse{\equal{\showdotontitles}{true}}{ % Agrega punto tras el número
	\renewcommand\thesubsubsubsection{\thesubsubsection.\arabic{subsubsubsection}.}
	\renewcommand\theparagraph{\thesubsubsubsection.\arabic{paragraph}.}
}{
	\renewcommand\thesubsubsubsection{\thesubsubsection.\arabic{subsubsubsection}}
	\renewcommand\theparagraph{\thesubsubsubsection.\arabic{paragraph}}
}

% Agrega compatibilidad de subsubsubsecciones al TOC
\makeatletter
	\renewcommand\paragraph{\@startsection{paragraph}{5}{\z@}
		{3.25ex \@plus1ex \@minus.2ex}
		{-1em}
		{\normalfont\normalsize\bfseries}}
	\renewcommand\subparagraph{\@startsection{subparagraph}{6}{\parindent}
		{3.25ex \@plus1ex \@minus .2ex}
		{-1em}
		{\normalfont\normalsize\bfseries}}
	\def\toclevel@subsubsubsection{4}
	\def\toclevel@paragraph{5}
	\def\toclevel@paragraph{6}
	\def\l@subsubsubsection{\@dottedtocline{4}{7em}{4em}}
	\def\l@paragraph{\@dottedtocline{5}{10em}{5em}}
	\def\l@subparagraph{\@dottedtocline{6}{14em}{6em}}
\makeatother

% Profundidad del índice y bookmarks pdf
\setcounter{tocdepth}{\indexdepth}


% INICIO DE LAS PÁGINAS
\begin{document}
	
% PORTADA
% Template:     Informe/Reporte LaTeX
% Documento:    Definición de portadas
% Versión:      6.1.2 (05/12/2018)
% Codificación: UTF-8
%
% Autor: Pablo Pizarro R. @ppizarror
%        Facultad de Ciencias Físicas y Matemáticas
%        Universidad de Chile
%        pablo.pizarro@ing.uchile.cl, ppizarror.com
%
% Manual template: [https://latex.ppizarror.com/Template-Informe/]
% Licencia MIT:    [https://opensource.org/licenses/MIT/]

% -----------------------------------------------------------------------------
% Importa las configuraciones
% -----------------------------------------------------------------------------
% Template:     Informe/Reporte LaTeX
% Documento:    Configuraciones adicionales de portadas
% Versión:      5.2.7 (21/05/2018)
% Codificación: UTF-8
%
% Autor: Pablo Pizarro R. @ppizarror
%        Facultad de Ciencias Físicas y Matemáticas
%        Universidad de Chile
%        pablo.pizarro@ing.uchile.cl, ppizarror.com
%
% Manual template: [http://latex.ppizarror.com/Template-Informe/]
% Licencia MIT:    [https://opensource.org/licenses/MIT/]

% style15 [A]
\def\headerimageA {departamentos/uchile}            % Imagen en el header
\def\headerimagescaleA {0.4}                        % Escala de la imagen

% style16 [B]
\def\portraitbackgroundimageB {ejemplos/portada-1}  % Imagen de fondo
\def\portraitbackgroundcolorB {ocre}                % Color principal

% style17 [C]
\def\portraitimageC {img/ejemplos/test-image}       % Imagen de la portada
\def\portraitimageboxedC {true}                     % Imagen recuadrada
\def\portraitimageboxedwidthC {0.5}                 % Grosor línea recuadro
\def\portraitimagewidthC {8cm}                      % Ancho de la imagen en cm

% style18 [D]
\def\portraitimageD {img/ejemplos/test-image-wrap}  % Imagen de la portada
\def\portraitimageboxedD {false}                    % Imagen recuadrada
\def\portraitimageboxedwidthD {0.5}                 % Grosor línea recuadro
\def\portraitimagewidthD {4cm}                      % Ancho de la imagen en cm
 % !FILE <NL,NODIST>

% -----------------------------------------------------------------------------
% Se escribe el header de la portada
% -----------------------------------------------------------------------------
\newpage
\renewcommand{\thepage}{\nameportraitpage} % !NL

% -----------------------------------------------------------------------------
% Estilos de portada
% -----------------------------------------------------------------------------
\ifthenelse{\equal{\portraitstyle}{style1}}{
	\setpagemargincm{\pagemarginleft}{\firstpagemargintop}{\pagemarginright}{\pagemarginbottom}
	\pagestyle{fancy}

	% Escribe el header
	\fancyhf{}
	\fancyhead[L]{
		\nombreuniversidad ~ \\
		\nombrefacultad ~ \\
		\departamentouniversidad ~ \\
		\vspace{-0.43cm}
	}
	\fancyhead[R]{
		\includegraphics[scale=\imagendepartamentoescala]{\imagendepartamento}
		\hspace{-0.255cm}
	}

	% Título y tema a tratar
	~ \\
	\vfill
	\begin{center}
		\textcolor{\portraittitlecolor}{
			{\noindent \Huge{\titulodelinforme} \vspace{0.5cm}} ~ \\
			{\noindent \large{\temaatratar}}
		}
	\end{center}

	% Tabla de integrantes
	\vfill
	\noindent
	\begin{minipage}{1.0\textwidth}
		\begin{flushright}
			\tablaintegrantes
		\end{flushright}
	\end{minipage}
}{
\ifthenelse{\equal{\portraitstyle}{style2}}{
	\setpagemargincm{\pagemarginleft}{\firstpagemargintop}{\pagemarginright}{\pagemarginbottom}
	\pagestyle{fancy}

	% Escribe el header
	\fancyhf{}
	\fancyhead[L]{
		\nombreuniversidad ~ \\
		\nombrefacultad ~ \\
		\departamentouniversidad ~ \\
		\vspace{-0.43cm}
	}
	\fancyhead[R]{
		\includegraphics[scale=\imagendepartamentoescala]{\imagendepartamento}
		\hspace{-0.255cm}
	}

	% Nombre de curso y título
	~ \\
	\vfill
	\begin{center}
		{\noindent \LARGE{\nombredelcurso} \vspace{0.3cm}} ~ \\
		\vspace*{1.5cm}
		\textcolor{\portraittitlecolor}{
			{\centering \noindent \Huge{\titulodelinforme} \vspace{0.3cm}} ~ \\
			{\noindent \large{\temaatratar}}
		}
	\end{center}

	% Tabla de integrantes
	\vfill
	\noindent
	\begin{minipage}{1.0\textwidth}
		\begin{flushright}
			\tablaintegrantes
		\end{flushright}
	\end{minipage}
}{
\ifthenelse{\equal{\portraitstyle}{style3}}{
	\setpagemargincm{\pagemarginleft}{\firstpagemargintop}{\pagemarginright}{\pagemarginbottom}
	\pagestyle{fancy}

	% Escribe el header
	\fancyhf{}
	\fancyhead[L]{
		\nombreuniversidad ~ \\
		\nombrefacultad ~ \\
		\departamentouniversidad ~ \\
		\vspace{-0.43cm}
	}
	\fancyhead[R]{
		\includegraphics[scale=\imagendepartamentoescala]{\imagendepartamento}
		\hspace{-0.255cm}
	}

	% Curso y código - título informe
	~ \\
	\vfill
	\begin{center}
		\vspace*{-1.0cm}
		{\noindent \huge{\nombredelcurso} \vspace{0.3cm}} ~ \\
		{\noindent \large{Código del curso: \codigodelcurso}} ~ \\
		\vspace*{1.8cm}
		\textcolor{\portraittitlecolor}{
			{\noindent \Huge{\titulodelinforme} \vspace{0.3cm}} ~ \\
			{\noindent \large{\temaatratar}}
		}
	\end{center}

	% Tabla de integrantes
	\vfill
	\noindent
	\begin{minipage}{1.0\textwidth}
		\begin{flushright}
			\tablaintegrantes
		\end{flushright}
	\end{minipage}
}{
\ifthenelse{\equal{\portraitstyle}{style4}}{
	\setpagemargincm{\pagemarginleft}{\pagemargintop}{\pagemarginright}{\pagemarginbottom}
	\thispagestyle{empty}

	% Escribe el header
	\vspace*{-1.5cm}
	\noindent \includegraphics[width=1.75cm]{departamentos/uchile2}
	\hspace*{-0.15cm}
	\begin{tabular}{l}
		\small \scshape{\MakeUppercase{\nombreuniversidad}} ~ \\
		\small \scshape{\MakeUppercase{\nombrefacultad}} ~ \\
		\small \scshape{\MakeUppercase{\departamentouniversidad}} ~ \\
		\vspace*{1.25cm}\mbox{}
	\end{tabular}

	% Título portada - curso
	\vfill
	\begin{center}
		{\fontsize{22pt}{10pt} \selectfont
			\noindent \textcolor{\portraittitlecolor}{\titulodelinforme} \vspace*{0.35cm}} ~ \\
		{\noindent \fontsize{10pt}{5pt} \selectfont \textcolor{\portraittitlecolor}{\codigodelcurso\ - \nombredelcurso}} ~ \\
		\vspace*{3cm}
	\end{center}

	% Tabla de integrantes
	\vfill
	\noindent
	\begin{minipage}{1.0\textwidth}
		\begin{flushright}
			\tablaintegrantes
		\end{flushright}
	\end{minipage}
}{
\ifthenelse{\equal{\portraitstyle}{style5}}{
	\setpagemargincm{\pagemarginleft}{\pagemargintop}{\pagemarginright}{\pagemarginbottom}
	\thispagestyle{empty}

	% Escribe el header
	\includegraphics[width=1.5cm]{departamentos/uchile3}
	\hspace{-0.2cm}
	\begin{tabular}{l}
		\small \scshape{\MakeUppercase{\nombreuniversidad}} ~ \\
		\small \scshape{\MakeUppercase{\nombrefacultad}} ~ \\
		\small \scshape{\MakeUppercase{\departamentouniversidad}} ~ \\
		\vspace*{1cm}\mbox{}
	\end{tabular}

	% Título informe - tema
	\vfill
	\begin{center}
		\fontsize{8mm}{9mm}\selectfont
		\textcolor{\portraittitlecolor}{
			\noindent \titulodelinforme ~ \\
			\noindent \temaatratar ~ \\
		}
		\vspace*{1cm}
		\footnotesize{\codigodelcurso\ - \nombredelcurso} ~ \\
		\vspace*{1.4cm}
	\end{center}

	% Tabla de integrantes
	\vfill
	\begin{center}
		\noindent \normalsize{\tablaintegrantes}
	\end{center}
}{
\ifthenelse{\equal{\portraitstyle}{style6}}{
	\setpagemargincm{\pagemarginleft}{\pagemargintop}{\pagemarginright}{\pagemarginbottom}
	\thispagestyle{empty}

	% Escribe el header
	\begin{wrapfigure}{l}{0.3\textwidth}
		\vspace{-0.69cm}
		\noindent \hspace{-1.10cm} \includegraphics[scale=1.35]{departamentos/fcfm2}
	\end{wrapfigure}
	\hspace*{0.3cm}
	\noindent \textsc{\color{red} \hspace{-2.2cm} \departamentouniversidad} ~ \\
	\hspace*{0.3cm}
	\noindent \textsc{\color{dkgray} \hspace{-1.6cm} \nombrefacultad} ~ \\
	\hspace*{0.3cm}
	\noindent \textsc{\color{dkgray} \hspace{-1.6cm} \nombreuniversidad} ~ \\
	\hspace*{0.3cm}
	\noindent \textsc{\color{dkgray} \hspace{-1.6cm} \codigodelcurso \nombredelcurso} ~ \\

	% Título informe - tema
	\vfill
	\begin{center}
		\vspace*{0.5cm}
		{\color{dkgray} \Large \textbf{\MakeUppercase{\temaatratar}}} ~ \\
		\noindent \rule{\linewidth}{0.3mm} ~ \\
		\Huge \textup \bfseries \textsc{\textcolor{\portraittitlecolor}{\titulodelinforme}} ~ \\
		\noindent \rule{\linewidth}{0.3mm} ~ \\
	\end{center}
	\begin{minipage}{.5\textwidth}
		~
	\end{minipage}

	% Tabla de integrantes
	\vfill
	\begin{minipage}{1.0\textwidth}
		\begin{flushright}
			\noindent \tablaintegrantes
		\end{flushright}
	\end{minipage}
}{
\ifthenelse{\equal{\portraitstyle}{style7}}{
	\setpagemargincm{\pagemarginleft}{\pagemargintop}{\pagemarginright}{\pagemarginbottom}
	\thispagestyle{empty}

	% Escribe el header
	\begin{center}
		\vspace*{-1.5cm}
		\includegraphics[scale=\imagendepartamentoescala]{\imagendepartamento}
		\hspace*{-0.15cm}
		\begin{tabular}{l}
			\vspace*{0.26cm}\mbox{} ~ \\
			\small \textsc{\MakeUppercase{\nombreuniversidad}} ~ \\
			\small \textsc{\MakeUppercase{\nombrefacultad}} ~ \\
			\small \textsc{\MakeUppercase{\departamentouniversidad}} ~ \\
			\vspace*{1.25cm}\mbox{}
		\end{tabular}
	\end{center}

	% Título informe - tema
	\vfill
	\begin{center}
		\noindent \rule{\textwidth}{0.4mm} \\ \vspace{0.3cm}
		{\huge \textcolor{\portraittitlecolor}{\titulodelinforme} \vspace{0.2cm} ~ \\}
		\noindent \rule{\textwidth}{0.4mm} ~ \\ \vspace{0.40cm}
		{\large \textcolor{\portraittitlecolor}{\temaatratar} ~ \\}
	\end{center}

	% Tabla de integrantes
	\vfill
	\noindent
	\begin{minipage}{1.0\textwidth}
		\begin{flushright}
			\scshape{\tablaintegrantes}
		\end{flushright}
	\end{minipage}
}{
\ifthenelse{\equal{\portraitstyle}{style8}}{
	\setpagemargincm{\pagemarginleft}{\pagemargintop}{\pagemarginright}{\pagemarginbottom}
	\thispagestyle{empty}

	% Escribe el header
	\begin{center}
		\vspace*{-1.0cm}
		\begin{tabular}{c}
			\includegraphics[scale=\imagendepartamentoescala]{\imagendepartamento} \vspace{0.5cm} ~ \\
			\small \scshape{\MakeUppercase{\nombreuniversidad}} ~ \\
			\small \scshape{\MakeUppercase{\nombrefacultad}} ~ \\
			\small \scshape{\MakeUppercase{\departamentouniversidad}}
		\end{tabular}
	\end{center}

	% Título informe - tema
	\vfill
	\begin{center}
		\noindent \rule{\textwidth}{0.4mm} \\ \vspace{0.3cm}
		{\huge \textcolor{\portraittitlecolor}{\titulodelinforme} \vspace{0.2cm} ~ \\}
		\noindent \rule{\textwidth}{0.4mm} ~ \\ \vspace{0.40cm}
		{\large \textcolor{\portraittitlecolor}{\temaatratar} ~ \\}
	\end{center}

	% Tabla de integrantes
	\vfill
	\noindent
	\begin{minipage}{1.0\textwidth}
		\begin{flushright}
			\scshape{\tablaintegrantes}
		\end{flushright}
	\end{minipage}
}{
\ifthenelse{\equal{\portraitstyle}{style9}}{
	\setpagemargincm{\pagemarginleft}{\pagemargintop}{\pagemarginright}{\pagemarginbottom}
	\thispagestyle{empty}

	% Título del informe
	\noindent \includegraphics[scale=\imagendepartamentoescala]{\imagendepartamento}
	\vfill
	\begin{center}
		\noindent \rule{\textwidth}{0.4mm} \\ \vspace{0.3cm}
		{\huge \textcolor{\portraittitlecolor}{\titulodelinforme} \vspace{0.2cm} \\}
		\noindent \rule{\textwidth}{0.4mm} \\ \vspace{0.35cm}
		{\large \textcolor{\portraittitlecolor}{\temaatratar} \\}
	\end{center}

	% Nombre universidad - departamento
	\vfill
	\begin{center}
		\begin{tabular}{c}
			\small \scshape{\MakeUppercase{\nombreuniversidad}} ~ \\
			\small \scshape{\MakeUppercase{\nombrefacultad}} ~ \\
			\small \scshape{\MakeUppercase{\departamentouniversidad}}
		\end{tabular}
	\end{center}

	% Tabla de integrantes
	\vfill
	\begin{center}
		\indent \scshape{\tablaintegrantes}
	\end{center}
}{
\ifthenelse{\equal{\portraitstyle}{style10}}{
	\setpagemargincm{\pagemarginleft}{\pagemargintop}{\pagemarginright}{\pagemarginbottom}
	\thispagestyle{empty}

	% Header
	~ \\
	\vfill
	\begin{center}
		\ifthenelse{\equal{\nombreuniversidad}{\xspace}}{
			\noindent {\large \textsc{\departamentouniversidad}}
		}{
			\noindent {\large \textsc{\nombreuniversidad, \departamentouniversidad}}
		}
		\vspace{1.0cm}
	\end{center}

	% Título informe
	\vfill
	\begin{center}
		\noindent {\large \scshape{\nombredelcurso}} \vspace{0.5cm} ~ \\
		\noindent {\large \scshape{\codigodelcurso}} \vspace{0.5cm} ~ \\
		\noindent \rule{\textwidth}{0.4mm} \\ \vspace{0.3cm}
		{\huge \bfseries \textcolor{\portraittitlecolor}{\titulodelinforme} \vspace{0.2cm} \\}
		\noindent \rule{\textwidth}{0.4mm} \\ \vspace{2.5cm}
	\end{center}

	% Tabla de integrantes
	\vfill
	\begin{center}
		\indent \tablaintegrantes
	\end{center}
	\vfill
	~ \\
}{
\ifthenelse{\equal{\portraitstyle}{style11}}{
	\setpagemargincm{\pagemarginleft}{\pagemargintop}{\pagemarginright}{\pagemarginbottom}
	\thispagestyle{empty}

	% Header
	\begin{center}
		\vspace*{-1.0cm}
		\scshape{\nombreuniversidad} ~ \\
		\scshape{\nombrefacultad} ~ \\
		\scshape{\departamentouniversidad}
	\end{center}

	% Título informe
	\vfill
	\begin{center}
		{\setstretch{1.2} \fontsize{21pt}{22pt} \selectfont \textcolor{\portraittitlecolor}{\scshape{\titulodelinforme}} \vspace{0.5cm}} ~ \\
		{\fontsize{13pt}{10pt} \selectfont \textcolor{\portraittitlecolor}{\scshape{\temaatratar}}}
	\end{center}

	% Tabla de integrantes
	\vfill
	\begin{center}
		\indent \tablaintegrantes
	\end{center}
}{
\ifthenelse{\equal{\portraitstyle}{style12}}{
	\setpagemargincm{\pagemarginleft}{\pagemargintop}{\pagemarginright}{\pagemarginbottom}
	\thispagestyle{empty}

	% Imagen departamento
	\begin{center}
		\vspace*{-1.0cm}
		\includegraphics[scale=\imagendepartamentoescala]{\imagendepartamento}
	\end{center}

	% Título informe
	\vfill
	\begin{center}
		{\bf \Huge \scshape{\textcolor{\portraittitlecolor}{\titulodelinforme}} \vspace{0.3cm}} \\
		{\bf \Large \textcolor{\portraittitlecolor}{\temaatratar}}
	\end{center}

	% Tabla de integrantes
	\vfill
	\begin{flushright}
		\noindent \tablaintegrantes
	\end{flushright}

	% Footer
	\vspace{0.5cm}
	\noindent \rule{\textwidth}{0.4mm}
	\begin{center}
		\ifthenelse{\equal{\nombreuniversidad}{\xspace}}{
			\scshape{\nombrefacultad} \\
		}{
			\scshape{\nombreuniversidad, \nombrefacultad} \\
		}
		\scshape{\departamentouniversidad}
	\end{center}
}{
\ifthenelse{\equal{\portraitstyle}{style13}}{
	\setpagemargincm{\pagemarginleft}{\pagemargintop}{\pagemarginright}{\pagemarginbottom}
	\thispagestyle{empty}

	% Header
	\noindent
	\vspace*{-1.5cm}
	\begin{flushleft}
		\begin{minipage}{0.65\textwidth}
			\ifthenelse{\equal{\nombreuniversidad}{\xspace}}{
				{\fontsize{3.5mm}{0.5mm} \selectfont \noindent \textsf{\nombrefacultad}} ~ \\
			}{
				{\fontsize{3.5mm}{0.5mm} \selectfont \noindent \textsf{\nombreuniversidad, \nombrefacultad}} ~ \\
			}
			\noindent {\fontsize{3.0mm}{0.5mm} \selectfont \textsf{\departamentouniversidad} \vspace{-0.2cm}} ~ \\
			\noindent \textcolor{gray}{\rule{\textwidth}{0.3mm}}
		\end{minipage}
	\end{flushleft}
	\vspace*{-2.15cm}
	\begin{flushright}
		\begin{minipage}{0.3\textwidth}
			\noindent \includegraphics[width=1.0\textwidth]{\imagendepartamento}
		\end{minipage}
	\end{flushright}

	% Título informe
	\vfill
	\begin{center}
		\begin{minipage}{0.9\textwidth}
			\begin{framed}
				\LARGE
				\vspace{1cm}
				\centering \textcolor{\portraittitlecolor}{\textbf{\titulodelinforme}}
				\vspace{1cm}
			\end{framed}
		\end{minipage}
	\end{center}

	% Tabla de integrantes
	\vfill
	\begin{flushright}
		\noindent \textsf{\tablaintegrantes}
	\end{flushright}
}{
\ifthenelse{\equal{\portraitstyle}{style14}}{
	\setpagemargincm{\pagemarginleft}{\pagemargintop}{\pagemarginright}{\pagemarginbottom}
	\thispagestyle{empty}

	% Header
	\noindent
	\begin{flushleft}
		\vspace*{-1.0cm}
		\noindent \includegraphics[scale=\imagendepartamentoescala]{\imagendepartamento} \\
	\end{flushleft}

	% Título del informe
	\vfill
	{\bf \huge \noindent \textcolor{\portraittitlecolor}{\textsf{\MakeUppercase{\titulodelinforme}} \vspace*{0.05cm}}} \\
	{\bf \large \noindent \textcolor{\portraittitlecolor}{\textsf{\MakeUppercase{\temaatratar}}}} \\

	% Tabla de integrantes
	\vfill
	\begin{flushright}
		\noindent \textsf{\tablaintegrantes}
	\end{flushright}
}{
\ifthenelse{\equal{\portraitstyle}{style15}}{
	\setpagemargincm{\pagemarginleft}{\pagemargintop}{\pagemarginright}{\pagemarginbottom}
	\thispagestyle{empty}
	\checkextravarexist{\headerimageA}{Defina la imagen extra de la portada en el archivo lib/page/portrait-config.tex (VERSION NORMAL) o bien en el bloque PORTADA (VERSION COMPACTA)}
	\checkextravarexist{\headerimagescaleA}{Defina la escala de la imagen extra de la portada en el archivo lib/page/portrait-config.tex (VERSION NORMAL) o bien en el bloque PORTADA (VERSION COMPACTA)}

	% Header
	\vspace*{-1.5cm}
	\noindent \begin{minipage}{0.8\textwidth}
		\noindent \begin{minipage}{0.22\textwidth}
			\includegraphics[scale=1.0]{departamentos/fcfm2} \\
		\end{minipage}
		\begin{minipage}{0.6\textwidth}
			\begin{flushleft}
				\textsc{
				\begin{tabular}{l}
					{\small \nombreuniversidad} ~ \\
					{\small \nombrefacultad} ~ \\
					{\small \departamentouniversidad}
				\end{tabular}
				}
			\end{flushleft}
		\end{minipage}
	\end{minipage}
	\noindent \begin{minipage}{0.2\textwidth}
		\begin{flushright}
			\ifthenelse{\isundefined{\headerimageA}}{}{
				\ifthenelse{\isundefined{\headerimagescaleA}}{}{
					\noindent \includegraphics[scale=\headerimagescaleA]{\headerimageA} \\
				}
			}
		\end{flushright}
	\end{minipage}

	% Título informe
	\vfill
	\begin{center}
		{\fontsize{25pt}{15pt} \selectfont \textcolor{\portraittitlecolor}{\textbf{\titulodelinforme}} \vspace{0.7cm}} \\
		{\Large \textcolor{\portraittitlecolor}{\temaatratar}}
	\end{center}

	% Tabla de integrantes
	\vfill
	\begin{center}
		\noindent \tablaintegrantes
	\end{center}
}{
\ifthenelse{\equal{\portraitstyle}{style16}}{
	\setpagemargincm{\pagemarginleft}{\pagemargintop}{\pagemarginright}{\pagemarginbottom}
	\checkextravarexist{\portraitbackgroundimageB}{[portrait-style16] Defina el fondo de la portada en el archivo lib/page/portrait-config.tex (VERSION NORMAL) o bien en el bloque PORTADA (VERSION COMPACTA)}
	\checkextravarexist{\portraitbackgroundcolorB}{[portrait-style16] Defina el color del bloque del titulo de la portada en el archivo lib/page/portrait-config.tex (VERSION NORMAL) o bien en el bloque PORTADA (VERSION COMPACTA)}

	% Imagen fondo, título informe
	\begingroup
		\thispagestyle{empty}
		\begin{tikzpicture}[remember picture,overlay]
			\node[inner sep=0pt] (background) at (current page.center) {\includegraphics[width=\paperwidth]{\portraitbackgroundimageB}};
			\draw (current page.center) node [fill=\portraitbackgroundcolorB!30!white,fill opacity=0.6,text opacity=1,inner sep=1cm]{\Huge\centering\bfseries\sffamily\parbox[c][][t]{\paperwidth}{
					\centering \textcolor{\portraittitlecolor}{\titulodelinforme} \\ [10pt]
					{\Large \temaatratar} \\ [25pt]
					{\huge \autordeldocumento}}};
		\end{tikzpicture}
		\vfill
	\endgroup
}{
\ifthenelse{\equal{\portraitstyle}{style17}}{
	\setpagemargincm{\pagemarginleft}{\firstpagemargintop}{\pagemarginright}{\pagemarginbottom}
	\pagestyle{fancy}
	\checkextravarexist{\portraitimageC}{[portrait-style17] Defina la imagen de la portada en el archivo lib/page/portrait-config.tex (VERSION NORMAL) o bien en el bloque PORTADA (VERSION COMPACTA)}
	\checkextravarexist{\portraitimageboxedC}{[portrait-style17] Defina si la imagen de la portada se encierra en un recuadro en el archivo lib/page/portrait-config.tex (VERSION NORMAL) o bien en el bloque PORTADA (VERSION COMPACTA)}
	\checkextravarexist{\portraitimageboxedwidthC}{[portrait-style17] Defina el grosor del recuadro de la imagen de la portada en el archivo lib/page/portrait-config.tex (VERSION NORMAL) o bien en el bloque PORTADA (VERSION COMPACTA)}
	\checkextravarexist{\portraitimagewidthC}{[portrait-style17] Defina los parametros de la imagen de la portada en el archivo lib/page/portrait-config.tex (VERSION NORMAL) o bien en el bloque PORTADA (VERSION COMPACTA)}

	% Header
	\fancyhf{}
	\fancyhead[L]{
		\nombreuniversidad ~ \\
		\nombrefacultad ~ \\
		\departamentouniversidad ~ \\
		\vspace{-0.43cm}
	}
	\fancyhead[R]{
		\includegraphics[scale=\imagendepartamentoescala]{\imagendepartamento}
		\hspace{-0.255cm}
	}

	% Título del informe
	~ \\
	\vfill
	\begin{center}
		\textcolor{\portraittitlecolor}{
			{\noindent \Huge{\titulodelinforme} \vspace{0.5cm}} ~ \\
			{\noindent \large{\temaatratar}}
		}
	\end{center}

	% Imagen de la portada
	~ \\
	\ifthenelse{\equal{\portraitimageboxedC}{true}}{
		\insertimageboxed{\portraitimageC}{width=\portraitimagewidthC}{\portraitimageboxedwidthC}{}
	}{
		\insertimage{\portraitimageC}{width=\portraitimagewidthC}{}
	}
	~ \\
	\vfill

	% Tabla de integrantes
	\noindent
	\begin{minipage}{1.0\textwidth}
		\begin{flushright}
			\tablaintegrantes
		\end{flushright}
	\end{minipage}
}{
\ifthenelse{\equal{\portraitstyle}{style18}}{
	\setpagemargincm{\pagemarginleft}{\firstpagemargintop}{\pagemarginright}{\pagemarginbottom}
	\pagestyle{fancy}
	\checkextravarexist{\portraitimageD}{[portrait-style18] Defina la imagen de la portada en el archivo lib/page/portrait-config.tex (VERSION NORMAL) o bien en el bloque PORTADA (VERSION COMPACTA)}
	\checkextravarexist{\portraitimageboxedD}{[portrait-style18] Defina si la imagen de la portada se encierra en un recuadro en el archivo lib/page/portrait-config.tex (VERSION NORMAL) o bien en el bloque PORTADA (VERSION COMPACTA)}
	\checkextravarexist{\portraitimageboxedwidthD}{[portrait-style18] Defina el grosor del recuadro de la imagen de la portada en el archivo lib/page/portrait-config.tex (VERSION NORMAL) o bien en el bloque PORTADA (VERSION COMPACTA)}
	\checkextravarexist{\portraitimagewidthD}{[portrait-style18] Defina los parametros de la imagen de la portada en el archivo lib/page/portrait-config.tex (VERSION NORMAL) o bien en el bloque PORTADA (VERSION COMPACTA)}

	% Header
	\fancyhf{}
	\fancyhead[L]{
		\nombreuniversidad ~ \\
		\nombrefacultad ~ \\
		\departamentouniversidad ~ \\
		\vspace{-0.43cm}
	}
	\fancyhead[R]{
		\includegraphics[scale=\imagendepartamentoescala]{\imagendepartamento}
		\hspace{-0.255cm}
	}

	% Imagen de la portada
	~ \\
	\ifthenelse{\equal{\portraitimageboxedD}{true}}{
		\insertimageboxed{\portraitimageD}{width=\portraitimagewidthD}{\portraitimageboxedwidthD}{}
	}{
		\insertimage{\portraitimageD}{width=\portraitimagewidthD}{}
	}

	% Título del informe
	\vfill
	\begin{center}
		\textcolor{\portraittitlecolor}{
			{\noindent \Huge{\titulodelinforme} \vspace{0.5cm}} ~ \\
			{\noindent \large{\temaatratar}}
		}
	\end{center}
	\vfill

	% Tabla de integrantes
	\noindent
	\begin{minipage}{1.0\textwidth}
		\begin{flushright}
			\tablaintegrantes
		\end{flushright}
	\end{minipage}
}{
\ifthenelse{\equal{\portraitstyle}{\bgtemplatetestcode}}{
	\setpagemargincm{\pagemarginleft}{\pagemargintop}{\pagemarginright}{\pagemarginbottom}
	\pagestyle{empty}
	\pagecolor{lbrown}
	\begin{center}
		\vspace*{-1.0cm}
		\scshape{\nombreuniversidad} ~ \\
		\scshape{\nombrefacultad} ~ \\
		\scshape{\departamentouniversidad}
	\end{center}
	~ \\
	\begin{center}
		\bgtemplatetestimg
	\end{center}
	\begin{center}
		\vspace*{-6cm}
		{\setstretch{1.2} \fontsize{25pt}{22pt} \selectfont \textcolor{\portraittitlecolor}{\scshape{\titulodelinforme}} \vspace{0.5cm}} \\
		{\fontsize{15pt}{10pt} \selectfont \textcolor{\portraittitlecolor}{\scshape{\temaatratar}}}
	\end{center}
	\vfill
	\begin{flushright}
		\noindent \tablaintegrantes
	\end{flushright}
	\newpage
	\pagecolor{white}
}{
	\throwbadconfigondoc{Estilo de portada incorrecto}{\portraitstyle}{style1 .. style18}}}}}}}}}}}}}}}}}}}
}

% -----------------------------------------------------------------------------
% Añade una página en blanco al imprimir por las dos caras
% -----------------------------------------------------------------------------
\ifthenelse{\equal{\addemptypagetwosides}{true}}{
	\newpage
	\null
	\thispagestyle{empty}
	\renewcommand{\thepage}{}
	\newpage}{
}

% END % Se puede borrar

% CONFIGURACIÓN DE PÁGINA Y ENCABEZADOS
% Template:     Informe/Reporte LaTeX
% Documento:    Configuración de página
% Versión:      5.3.6 (21/05/2018)
% Codificación: UTF-8
%
% Autor: Pablo Pizarro R. @ppizarror
%        Facultad de Ciencias Físicas y Matemáticas
%        Universidad de Chile
%        pablo.pizarro@ing.uchile.cl, ppizarror.com
%
% Manual template: [http://latex.ppizarror.com/Template-Informe/]
% Licencia MIT:    [https://opensource.org/licenses/MIT/]

% Numeración de páginas
\newpage
\ifthenelse{\equal{\romanpageuppercase}{true}}{
	\pagenumbering{Roman}
}{
	\pagenumbering{roman}
}
\setcounter{page}{1}
\setcounter{footnote}{1}

% Márgenes de páginas y tablas
\setpagemargincm{\pagemarginleft}{\pagemargintop}{\pagemarginright}{\pagemarginbottom}
\def\arraystretch {\tablepadding} % Se ajusta el padding de las tablas
\def\pagewidth {0.995\linewidth} % Define ancho de la página en cm

% Se define el punto decimal
\ifthenelse{\equal{\pointdecimal}{true}}{
	\decimalpoint}{
}

% Definición de nombres de objetos
\renewcommand{\appendixname}{\nomltappendixsection} % Nombre del anexo (título)
\renewcommand{\appendixpagename}{\nameappendixsection} % Nombre del anexo en índice
\renewcommand{\appendixtocname}{\nameappendixsection} % Nombre del anexo en índice
\renewcommand{\contentsname}{\nomltcont}  % Nombre del índice
\renewcommand{\figurename}{\nomltwfigure} % Nombre de la leyenda de las fig.
\renewcommand{\listfigurename}{\nomltfigure} % Nombre del índice de figuras
\renewcommand{\listtablename}{\nomlttable} % Nombre del índice de tablas
\renewcommand{\lstlistingname}{\nomltwsrc} % Nombre leyenda del código fuente
\renewcommand{\lstlistlistingname}{\nomltsrc} % Nombre índice código fuente
\renewcommand{\refname}{\namereferences} % Nombre de las referencias
\renewcommand{\tablename}{\nomltwtable} % Nombre de la leyenda de tablas

% Estilo de títulos
\sectionfont{\color{\titlecolor} \fontsizetitle \styletitle \selectfont}
\subsectionfont{\color{\subtitlecolor} \fontsizesubtitle \stylesubtitle \selectfont}
\subsubsectionfont{\color{\subsubtitlecolor} \fontsizesubsubtitle \stylesubsubtitle \selectfont}

% Se crean los header-footer
\ifthenelse{\equal{\hfstyle}{style1}}{
	\pagestyle{fancy} \fancyhf{}
	\fancyhead[L]{\nouppercase{\rightmark}}
	\fancyhead[R]{\small \rm \thepage}
	\fancyfoot[L]{\small \rm \textit{\titulodelinforme}}
	\fancyfoot[R]{\small \rm \textit{\codigodelcurso \nombredelcurso}}
	\renewcommand{\headrulewidth}{0.5pt}
	\renewcommand{\footrulewidth}{0.5pt}
	\renewcommand{\sectionmark}[1]{\markboth{#1}{}}
}{
\ifthenelse{\equal{\hfstyle}{style2}}{
	\pagestyle{fancy} \fancyhf{}
	\fancyhead[L]{\nouppercase{\rightmark}}
	\fancyhead[R]{\small \rm \thepage}
	\fancyfoot[L]{\small \rm \textit{\titulodelinforme}}
	\fancyfoot[R]{\small \rm \textit{\codigodelcurso \nombredelcurso}}
	\renewcommand{\headrulewidth}{0.5pt}
	\renewcommand{\footrulewidth}{0pt}
	\renewcommand{\sectionmark}[1]{\markboth{#1}{}}
}{
\ifthenelse{\equal{\hfstyle}{style3}}{
	\pagestyle{fancy} \fancyhf{}
	\fancyhead[L]{
		\small \rm \textit{\codigodelcurso \nombredelcurso}
		\vspace{0.04cm}
	}
	\fancyhead[R]{
		\includegraphics[width=1.2cm]{departamentos/fcfm2}
		\vspace{-0.10cm}
	}
	\fancyfoot[C]{\thepage}
	\renewcommand{\headrulewidth}{0.5pt}
	\renewcommand{\footrulewidth}{0pt}
}{
\ifthenelse{\equal{\hfstyle}{style4}}{
	\pagestyle{fancy} \fancyhf{}
	\fancyhead[L]{\nouppercase{\rightmark}}
	\fancyhead[R]{}
	\fancyfoot[C]{\small \rm \thepage}
	\renewcommand{\headrulewidth}{0.5pt}
	\renewcommand{\footrulewidth}{0pt}
	\renewcommand{\sectionmark}[1]{\markboth{#1}{}}
}{
\ifthenelse{\equal{\hfstyle}{style5}}{
	\pagestyle{fancy} \fancyhf{}
	\fancyhead[L]{\codigodelcurso \nombredelcurso}
	\fancyhead[R]{\nouppercase{\rightmark}}
	\fancyfoot[L]{\departamentouniversidad, \nombreuniversidad}
	\fancyfoot[R]{\small \rm \thepage}
	\renewcommand{\headrulewidth}{0pt}
	\renewcommand{\footrulewidth}{0pt}
	\renewcommand{\sectionmark}[1]{\markboth{#1}{}}
}{
\ifthenelse{\equal{\hfstyle}{style6}}{
	\pagestyle{fancy} \fancyhf{}
	\fancyfoot[L]{\departamentouniversidad}
	\fancyfoot[C]{\thepage}
	\fancyfoot[R]{\nombreuniversidad}
	\renewcommand{\headrulewidth}{0pt}
	\renewcommand{\footrulewidth}{0pt}
	\setlength{\headheight}{49pt}
}{
\ifthenelse{\equal{\hfstyle}{style7}}{
	\pagestyle{fancy} \fancyhf{}
	\fancyfoot[C]{\thepage}
	\renewcommand{\headrulewidth}{0pt}
	\renewcommand{\footrulewidth}{0pt}
	\setlength{\headheight}{49pt}
}{
\ifthenelse{\equal{\hfstyle}{style8}}{
	\pagestyle{fancy} \fancyhf{}
	\fancyfoot[R]{\thepage}
	\renewcommand{\headrulewidth}{0pt}
	\renewcommand{\footrulewidth}{0pt}
	\setlength{\headheight}{49pt}
}{
\ifthenelse{\equal{\hfstyle}{style9}}{
	\pagestyle{fancy} \fancyhf{}
	\fancyhead[L]{\nouppercase{\rightmark}}
	\fancyhead[R]{}
	\fancyfoot[L]{\small \rm \textit{\titulodelinforme}}
	\fancyfoot[R]{\small \rm \thepage}
	\renewcommand{\headrulewidth}{0.5pt}
	\renewcommand{\footrulewidth}{0.5pt}
	\renewcommand{\sectionmark}[1]{\markboth{#1}{}}
}{
\ifthenelse{\equal{\hfstyle}{style10}}{
	\pagestyle{fancy} \fancyhf{}
	\fancyhead[L]{\nouppercase{\rightmark}}
	\fancyhead[R]{\small \rm \textit{\titulodelinforme}}
	\fancyfoot[L]{}
	\fancyfoot[R]{\small \rm \thepage}
	\renewcommand{\headrulewidth}{0.5pt}
	\renewcommand{\footrulewidth}{0.5pt}
	\renewcommand{\sectionmark}[1]{\markboth{#1}{}}
}{
	\throwbadconfigondoc{Estilo de header-footer incorrecto}{\hfstyle}{style1 .. style10}}}}}}}}}}
}

% Muestra los números de línea
\ifthenelse{\equal{\showlinenumbers}{true}}{
	\linenumbers}{
}


% RESUMEN O ABSTRACT
\begin{resumen}
	\lipsum[1] % Párrafo ejemplo, se puede borrar
\end{resumen}

% TABLA DE CONTENIDOS - ÍNDICE
% Template:     Informe LaTeX
% Documento:    Índice
% Versión:      DEV
% Codificación: UTF-8
%
% Autor: Pablo Pizarro R.
%        Facultad de Ciencias Físicas y Matemáticas
%        Universidad de Chile
%        pablo@ppizarror.com
%
% Manual template: [https://latex.ppizarror.com/informe]
% Licencia MIT:    [https://opensource.org/licenses/MIT]

\ifthenelse{\equal{\showindex}{true}}{
	
	% -------------------------------------------------------------------------
	% Crea nueva página y establece estilo de títulos
	% -------------------------------------------------------------------------
	\newpage
	\begingroup
	\sectionfont{\color{\indextitlecolor} \fontsizetitlei \styletitlei \selectfont}
	
	% -------------------------------------------------------------------------
	% Salta de página si está imprimiendo por ambas caras
	% -------------------------------------------------------------------------
	\ifthenelse{\equal{\addemptypagetwosides}{true}}{
		\checkoddpage
		\ifoddpage
		\else
			\newpage
			\null
			\thispagestyle{empty}
			\newpage
			\addtocounter{page}{-1}
		\fi}{
	}
	
	% -------------------------------------------------------------------------
	% Añade la entrada del índice a los marcadores del pdf
	% -------------------------------------------------------------------------
	\ifthenelse{\equal{\addindextobookmarks}{true}}{
		\belowpdfbookmark{\nomltcont}{contents}}{
	}
	\tocloftpagestyle{fancy}
	
	% -------------------------------------------------------------------------
	% Configuración del punto en índice
	% -------------------------------------------------------------------------
	\ifthenelse{\equal{\showdotaftersnum}{true}}{
		% Agrega los puntos
		\def\cftchapaftersnum {.}
		\def\cftsecaftersnum {.}
		\def\cftsubsecaftersnum {.}
		\def\cftsubsubsecaftersnum {.}
		\def\cftsubsubsubsecaftersnum {.}
		% Modifica los márgenes
		\def\cftsecnumwidth {1.9em}
		\def\cftsubsecnumwidth {2.57em} % Incremento 0.67
		\renewcommand\cftsubsubsecnumwidth{3.35em} % Incremento 0.78
		\setlength{\cftsubsecindent}{1.91em}
		\setlength{\cftsubsubsecindent}{4.48em} % Incremento 2.57
		}{
	}
	
	% -------------------------------------------------------------------------
	% Configuración carácter número de página
	% -------------------------------------------------------------------------
	\renewcommand{\cftdot}{\charnumpageindex}
	
	% -------------------------------------------------------------------------
	% Configuración del punto en número de objetos
	% -------------------------------------------------------------------------
	\def\cftfigaftersnum {\charafterobjectindex\enspace} % Figuras
	\def\cftsubfigaftersnum {\charafterobjectindex\enspace} % Subfiguras
	\def\cfttabaftersnum {\charafterobjectindex\enspace} % Tablas
	\def\cftlstlistingaftersnum {\charafterobjectindex\enspace} % Códigos fuente
	\def\cftmyindexequationsaftersnum {\charafterobjectindex\enspace} % Ecuaciones
	
	% -------------------------------------------------------------------------
	% Desactiva los números de línea
	% -------------------------------------------------------------------------
	\ifthenelse{\equal{\showlinenumbers}{true}}{
		\nolinenumbers}{
	}

	% -------------------------------------------------------------------------
	% Cambia tabulación índice de objetos para alinear con contenidos
	% -------------------------------------------------------------------------
	\ifthenelse{\equal{\objectindexindent}{true}}{
		\setlength{\cfttabindent}{1.9em} % Tablas
		\setlength{\cftfigindent}{1.9em} % Figuras
		\setlength{\cftsubfigindent}{1.9em} % Subfiguras
		\setlength{\cftmyindexequationsindent}{1.9em} % Ecuaciones
		\def\cftlstlistingindent {1.9em} % Códigos fuente
	}{
		\setlength{\cfttabindent}{0em} % Tablas
		\setlength{\cftfigindent}{0em} % Figuras
		\setlength{\cftsubfigindent}{0em} % Subfiguras
		\setlength{\cftmyindexequationsindent}{0em} % Ecuaciones
		\def\cftlstlistingindent {0em} % Códigos fuente
	}
	
	% -------------------------------------------------------------------------
	% Calcula tamaño del margen de los números en objetos del índice
	% -------------------------------------------------------------------------
	
	% Código fuente
	\ifthenelse{\equal{\showsectioncaptioncode}{none}}{
		\def\cftdefautnumwidthcode {3.0em} % Añade +0.7em
		\def\cftdefaultnumwidthromancode {5.25em} % Añade +0.5em para no overflow
	}{
	\ifthenelse{\equal{\showsectioncaptioncode}{sec}}{
		\def\cftdefautnumwidthcode {3.7em}
		\def\cftdefaultnumwidthromancode {5.75em}
	}{
	\ifthenelse{\equal{\showsectioncaptioncode}{ssec}}{
		\def\cftdefautnumwidthcode {4.4em}
		\def\cftdefaultnumwidthromancode {6.25em}
	}{
	\ifthenelse{\equal{\showsectioncaptioncode}{sssec}}{
		\def\cftdefautnumwidthcode {5.1em}
		\def\cftdefaultnumwidthromancode {6.75em}
	}{
	\ifthenelse{\equal{\showsectioncaptioncode}{ssssec}}{
		\def\cftdefautnumwidthcode {5.8em}
		\def\cftdefaultnumwidthromancode {7.25em}
	}{
	\ifthenelse{\equal{\showsectioncaptioncode}{chap}}{
		\def\cftdefautnumwidthcode {3.0em}
		\def\cftdefaultnumwidthromancode {5.25em}
	}{
		\throwbadconfig{Valor configuracion incorrecto}{\showsectioncaptioncode}{none,chap,sec,ssec,sssec,ssssec}}}}}}
	}

	% Código fuente
	\ifthenelse{\equal{\showsectioncaptioneqn}{none}}{
		\def\cftdefautnumwidtheqn {3.0em} % Añade +0.7em
		\def\cftdefaultnumwidthromaneqn {5.25em} % Añade +0.5em para no overflow
	}{
	\ifthenelse{\equal{\showsectioncaptioneqn}{sec}}{
		\def\cftdefautnumwidtheqn {3.7em}
		\def\cftdefaultnumwidthromaneqn {5.75em}
	}{
	\ifthenelse{\equal{\showsectioncaptioneqn}{ssec}}{
		\def\cftdefautnumwidtheqn {4.4em}
		\def\cftdefaultnumwidthromaneqn {6.25em}
	}{
	\ifthenelse{\equal{\showsectioncaptioneqn}{sssec}}{
		\def\cftdefautnumwidtheqn {5.1em}
		\def\cftdefaultnumwidthromaneqn {6.75em}
	}{
	\ifthenelse{\equal{\showsectioncaptioneqn}{ssssec}}{
		\def\cftdefautnumwidtheqn {5.8em}
		\def\cftdefaultnumwidthromaneqn {7.25em}
	}{
	\ifthenelse{\equal{\showsectioncaptioneqn}{chap}}{
		\def\cftdefautnumwidtheqn {3.0em}
		\def\cftdefaultnumwidthromaneqn {5.25em}
	}{
		\throwbadconfig{Valor configuracion incorrecto}{\showsectioncaptioneqn}{none,chap,sec,ssec,sssec,ssssec}}}}}}
	}
	
	% Figuras
	\ifthenelse{\equal{\showsectioncaptionfig}{none}}{
		\def\cftdefautnumwidthfig {3.0em} % Añade +0.7em
		\def\cftdefaultnumwidthromanfig {5.25em} % Añade +0.5em
	}{
	\ifthenelse{\equal{\showsectioncaptionfig}{sec}}{
		\def\cftdefautnumwidthfig {3.7em}
		\def\cftdefaultnumwidthromanfig {5.75em}
	}{
	\ifthenelse{\equal{\showsectioncaptionfig}{ssec}}{
		\def\cftdefautnumwidthfig {4.4em}
		\def\cftdefaultnumwidthromanfig {6.25em}
	}{
	\ifthenelse{\equal{\showsectioncaptionfig}{sssec}}{
		\def\cftdefautnumwidthfig {5.1em}
		\def\cftdefaultnumwidthromanfig {6.75em}
	}{
	\ifthenelse{\equal{\showsectioncaptionfig}{ssssec}}{
		\def\cftdefautnumwidthfig {5.8em}
		\def\cftdefaultnumwidthromanfig {7.25em}
	}{
	\ifthenelse{\equal{\showsectioncaptionfig}{chap}}{
		\def\cftdefautnumwidthfig {3.0em}
		\def\cftdefaultnumwidthromanfig {5.25em}
	}{
		\throwbadconfig{Valor configuracion incorrecto}{\showsectioncaptionfig}{none,chap,sec,ssec,sssec,ssssec}}}}}}
	}
	
	% Tablas
	\ifthenelse{\equal{\showsectioncaptiontab}{none}}{
		\def\cftdefautnumwidthtab {3.0em} % Añade +0.7em
		\def\cftdefaultnumwidthromantab {5.25em} % Añade +0.5em
	}{
	\ifthenelse{\equal{\showsectioncaptiontab}{sec}}{
		\def\cftdefautnumwidthtab {3.7em}
		\def\cftdefaultnumwidthromantab {5.75em}
	}{
	\ifthenelse{\equal{\showsectioncaptiontab}{ssec}}{
		\def\cftdefautnumwidthtab {4.4em}
		\def\cftdefaultnumwidthromantab {6.25em}
	}{
	\ifthenelse{\equal{\showsectioncaptiontab}{sssec}}{
		\def\cftdefautnumwidthtab {5.1em}
		\def\cftdefaultnumwidthromantab {6.75em}
	}{
	\ifthenelse{\equal{\showsectioncaptiontab}{ssssec}}{
		\def\cftdefautnumwidthtab {5.8em}
		\def\cftdefaultnumwidthromantab {7.25em}
	}{
	\ifthenelse{\equal{\showsectioncaptiontab}{chap}}{
		\def\cftdefautnumwidthtab {3.0em}
		\def\cftdefaultnumwidthromantab {5.25em}
	}{
		\throwbadconfig{Valor configuracion incorrecto}{\showsectioncaptiontab}{none,chap,sec,ssec,sssec,ssssec}}}}}}
	}
	
	% Configuración identado de títulos de objetos después del número
	\def\cftfignumwidth {\cftdefautnumwidth}
	% \def\cftsubfignumwidth {\cftdefautnumwidth}
	\def\cfttabnumwidth {\cftdefautnumwidth}
	\def\cftlstlistingnumwidth {\cftdefautnumwidth}
	
	% Código fuente
	\ifthenelse{\equal{\captionnumcode}{arabic}}{ % No hace nada (default)
		\def\cftlstlistingnumwidth {\cftdefautnumwidthcode}
	}{
		\ifthenelse{\equal{\captionnumcode}{roman}}{
			\def\cftlstlistingnumwidth {\cftdefaultnumwidthromancode}
		}{
		\ifthenelse{\equal{\captionnumcode}{Roman}}{
			\def\cftlstlistingnumwidth {\cftdefaultnumwidthromancode}
		}{
			\def\cftlstlistingnumwidth {\cftdefautnumwidthcode}
		}}
	}

	% Ecuaciones
	\ifthenelse{\equal{\captionnumequation}{arabic}}{ % No hace nada (default)
		\def\cftmyindexequationsnumwidth {\cftdefautnumwidtheqn}
	}{
		\ifthenelse{\equal{\captionnumequation}{roman}}{
			\def\cftmyindexequationsnumwidth {\cftdefaultnumwidthromaneqn}
		}{
		\ifthenelse{\equal{\captionnumequation}{Roman}}{
			\def\cftmyindexequationsnumwidth {\cftdefaultnumwidthromaneqn}
		}{
			\def\cftmyindexequationsnumwidth {\cftdefautnumwidtheqn}
		}}
	}
	
	% Figuras
	\ifthenelse{\equal{\captionnumfigure}{arabic}}{ % No hace nada (default)
		\def\cftfignumwidth {\cftdefautnumwidthfig}
	}{
		\ifthenelse{\equal{\captionnumfigure}{roman}}{
			\def\cftfignumwidth {\cftdefaultnumwidthromanfig}
		}{
			\ifthenelse{\equal{\captionnumfigure}{Roman}}{
				\def\cftfignumwidth {\cftdefaultnumwidthromanfig}
			}{
				\def\cftfignumwidth {\cftdefautnumwidthfig}
			}}
	}

	% Tablas
	\ifthenelse{\equal{\captionnumtable}{arabic}}{ % No hace nada (default)
		\def\cfttabnumwidth {\cftdefautnumwidthtab}
	}{
		\ifthenelse{\equal{\captionnumtable}{roman}}{
			\def\cfttabnumwidth {\cftdefaultnumwidthromantab}
		}{
			\ifthenelse{\equal{\captionnumtable}{Roman}}{
				\def\cfttabnumwidth {\cftdefaultnumwidthromantab}
			}{
				\def\cfttabnumwidth {\cftdefautnumwidthtab}
			}}
	}

	% -------------------------------------------------------------------------
	% Genera las funciones para los índices
	% -------------------------------------------------------------------------
	\newcommand{\LoIf}{ % Lista de figuras
		\iftotalfigures
			\ifthenelse{\equal{\indexnewpagef}{true}}{\newpage}{}
			\listoffigures
		\fi
	}
	\newcommand{\LoIt}{ % Tablas
		\iftotaltables
			\ifthenelse{\equal{\indexnewpaget}{true}}{\newpage}{}
			\listoftables
		\fi
	}
	\newcommand{\LoIc}{ % Códigos fuente (listings)
		\iftotallstlistings
			\ifthenelse{\equal{\indexnewpagec}{true}}{\newpage}{}
			\lstlistoflistings
		\fi
	}
	\newcommand{\LoIe}{ % Ecuaciones
		\iftotaltemplateIndexEquationss
			\ifthenelse{\equal{\indexnewpagee}{true}}{\newpage}{}
			\listofmyindexequations
		\fi
	}

	% -------------------------------------------------------------------------
	% Índice de contenidos
	% -------------------------------------------------------------------------
	\ifthenelse{\equal{\showindexofcontents}{true}}{
		\tableofcontents
	}{}
	
	% -------------------------------------------------------------------------
	% Índice de objetos
	% -------------------------------------------------------------------------
	\ifthenelse{\equal{\indexstyle}{tfce}}{\LoIt\LoIf\LoIc\LoIe}{\ifthenelse{\equal{\indexstyle}{}}{}{\ifthenelse{\equal{\indexstyle}{e}}{\LoIe}{\ifthenelse{\equal{\indexstyle}{c}}{\LoIc}{\ifthenelse{\equal{\indexstyle}{f}}{\LoIf}{\ifthenelse{\equal{\indexstyle}{t}}{\LoIt}{\ifthenelse{\equal{\indexstyle}{ec}}{\LoIe\LoIc}{\ifthenelse{\equal{\indexstyle}{ce}}{\LoIc\LoIe}{\ifthenelse{\equal{\indexstyle}{ef}}{\LoIe\LoIf}{\ifthenelse{\equal{\indexstyle}{fe}}{\LoIf\LoIe}{\ifthenelse{\equal{\indexstyle}{et}}{\LoIe\LoIt}{\ifthenelse{\equal{\indexstyle}{te}}{\LoIt\LoIe}{\ifthenelse{\equal{\indexstyle}{cf}}{\LoIc\LoIf}{\ifthenelse{\equal{\indexstyle}{fc}}{\LoIf\LoIc}{\ifthenelse{\equal{\indexstyle}{ct}}{\LoIc\LoIt}{\ifthenelse{\equal{\indexstyle}{tc}}{\LoIt\LoIc}{\ifthenelse{\equal{\indexstyle}{ft}}{\LoIf\LoIt}{\ifthenelse{\equal{\indexstyle}{tf}}{\LoIt\LoIf}{\ifthenelse{\equal{\indexstyle}{ecf}}{\LoIe\LoIc\LoIf}{\ifthenelse{\equal{\indexstyle}{efc}}{\LoIe\LoIf\LoIc}{\ifthenelse{\equal{\indexstyle}{cef}}{\LoIc\LoIe\LoIf}{\ifthenelse{\equal{\indexstyle}{cfe}}{\LoIc\LoIf\LoIe}{\ifthenelse{\equal{\indexstyle}{fec}}{\LoIf\LoIe\LoIc}{\ifthenelse{\equal{\indexstyle}{fce}}{\LoIf\LoIc\LoIe}{\ifthenelse{\equal{\indexstyle}{ect}}{\LoIe\LoIc\LoIt}{\ifthenelse{\equal{\indexstyle}{etc}}{\LoIe\LoIt\LoIc}{\ifthenelse{\equal{\indexstyle}{cet}}{\LoIc\LoIe\LoIt}{\ifthenelse{\equal{\indexstyle}{cte}}{\LoIc\LoIt\LoIe}{\ifthenelse{\equal{\indexstyle}{tec}}{\LoIt\LoIe\LoIc}{\ifthenelse{\equal{\indexstyle}{tce}}{\LoIt\LoIc\LoIe}{\ifthenelse{\equal{\indexstyle}{eft}}{\LoIe\LoIf\LoIt}{\ifthenelse{\equal{\indexstyle}{etf}}{\LoIe\LoIt\LoIf}{\ifthenelse{\equal{\indexstyle}{fet}}{\LoIf\LoIe\LoIt}{\ifthenelse{\equal{\indexstyle}{fte}}{\LoIf\LoIt\LoIe}{\ifthenelse{\equal{\indexstyle}{tef}}{\LoIt\LoIe\LoIf}{\ifthenelse{\equal{\indexstyle}{tfe}}{\LoIt\LoIf\LoIe}{\ifthenelse{\equal{\indexstyle}{cft}}{\LoIc\LoIf\LoIt}{\ifthenelse{\equal{\indexstyle}{ctf}}{\LoIc\LoIt\LoIf}{\ifthenelse{\equal{\indexstyle}{fct}}{\LoIf\LoIc\LoIt}{\ifthenelse{\equal{\indexstyle}{ftc}}{\LoIf\LoIt\LoIc}{\ifthenelse{\equal{\indexstyle}{tcf}}{\LoIt\LoIc\LoIf}{\ifthenelse{\equal{\indexstyle}{tfc}}{\LoIt\LoIf\LoIc}{\ifthenelse{\equal{\indexstyle}{ecft}}{\LoIe\LoIc\LoIf\LoIt}{\ifthenelse{\equal{\indexstyle}{ectf}}{\LoIe\LoIc\LoIt\LoIf}{\ifthenelse{\equal{\indexstyle}{efct}}{\LoIe\LoIf\LoIc\LoIt}{\ifthenelse{\equal{\indexstyle}{eftc}}{\LoIe\LoIf\LoIt\LoIc}{\ifthenelse{\equal{\indexstyle}{etcf}}{\LoIe\LoIt\LoIc\LoIf}{\ifthenelse{\equal{\indexstyle}{etfc}}{\LoIe\LoIt\LoIf\LoIc}{\ifthenelse{\equal{\indexstyle}{ceft}}{\LoIc\LoIe\LoIf\LoIt}{\ifthenelse{\equal{\indexstyle}{cetf}}{\LoIc\LoIe\LoIt\LoIf}{\ifthenelse{\equal{\indexstyle}{cfet}}{\LoIc\LoIf\LoIe\LoIt}{\ifthenelse{\equal{\indexstyle}{cfte}}{\LoIc\LoIf\LoIt\LoIe}{\ifthenelse{\equal{\indexstyle}{ctef}}{\LoIc\LoIt\LoIe\LoIf}{\ifthenelse{\equal{\indexstyle}{ctfe}}{\LoIc\LoIt\LoIf\LoIe}{\ifthenelse{\equal{\indexstyle}{fect}}{\LoIf\LoIe\LoIc\LoIt}{\ifthenelse{\equal{\indexstyle}{fetc}}{\LoIf\LoIe\LoIt\LoIc}{\ifthenelse{\equal{\indexstyle}{fcet}}{\LoIf\LoIc\LoIe\LoIt}{\ifthenelse{\equal{\indexstyle}{fcte}}{\LoIf\LoIc\LoIt\LoIe}{\ifthenelse{\equal{\indexstyle}{ftec}}{\LoIf\LoIt\LoIe\LoIc}{\ifthenelse{\equal{\indexstyle}{ftce}}{\LoIf\LoIt\LoIc\LoIe}{\ifthenelse{\equal{\indexstyle}{tecf}}{\LoIt\LoIe\LoIc\LoIf}{\ifthenelse{\equal{\indexstyle}{tefc}}{\LoIt\LoIe\LoIf\LoIc}{\ifthenelse{\equal{\indexstyle}{tcef}}{\LoIt\LoIc\LoIe\LoIf}{\ifthenelse{\equal{\indexstyle}{tcfe}}{\LoIt\LoIc\LoIf\LoIe}{\ifthenelse{\equal{\indexstyle}{tfec}}{\LoIt\LoIf\LoIe\LoIc}{\throwbadconfig{Estilo desconocido del indice}{\indexstyle}{tfce,,e,c,f,t,ec,ce,ef,fe,et,te,cf,fc,ct,tc,ft,tf,ecf,efc,cef,cfe,fec,fce,ect,etc,cet,cte,tec,tce,eft,etf,fet,fte,tef,tfe,cft,ctf,fct,ftc,tcf,tfc,ecft,ectf,efct,eftc,etcf,etfc,ceft,cetf,cfet,cfte,ctef,ctfe,fect,fetc,fcet,fcte,ftec,ftce,tecf,tefc,tcef,tcfe,tfec}}}}}}}}}}}}}}}}}}}}}}}}}}}}}}}}}}}}}}}}}}}}}}}}}}}}}}}}}}}}}}}}}}

	% -------------------------------------------------------------------------
	% Termina el bloque de índice
	% -------------------------------------------------------------------------
	\endgroup
	
	% -------------------------------------------------------------------------
	% Se añade una página en blanco
	% -------------------------------------------------------------------------
	\newpage
	\ifthenelse{\equal{\addemptypagetwosides}{true}}{
		\vfill
		\checkoddpage
		\ifoddpage
			\newpage
			\null
			\thispagestyle{empty}
			\newpage
			\addtocounter{page}{-1}
		\else
		\fi}{
	}

}{}

% END % Se puede borrar

% CONFIGURACIONES FINALES
% Template:     Informe/Reporte LaTeX
% Documento:    Configuraciones finales
% Versión:      6.0.0 (13/10/2018)
% Codificación: UTF-8
%
% Autor: Pablo Pizarro R. @ppizarror
%        Facultad de Ciencias Físicas y Matemáticas
%        Universidad de Chile
%        pablo.pizarro@ing.uchile.cl, ppizarror.com
%
% Manual template: [http://latex.ppizarror.com/Template-Informe/]
% Licencia MIT:    [https://opensource.org/licenses/MIT/]

% -----------------------------------------------------------------------------
% Se reestablecen headers y footers
% -----------------------------------------------------------------------------
\markboth{}{}
\newpage
\ifthenelse{\equal{\disablehfrightmark}{false}}{
	\ifthenelse{\equal{\hfstyle}{style1}}{
		\fancyhead[L]{\nouppercase{\leftmark}}}{
	}
	\ifthenelse{\equal{\hfstyle}{style2}}{
		\fancyhead[L]{\nouppercase{\leftmark}}}{
	}
	\ifthenelse{\equal{\hfstyle}{style4}}{
		\fancyhead[L]{\nouppercase{\leftmark}}}{
	}
	\ifthenelse{\equal{\hfstyle}{style5}}{
		\fancyhead[R]{\nouppercase{\leftmark}}}{
	}
	\ifthenelse{\equal{\hfstyle}{style9}}{
		\fancyhead[L]{\nouppercase{\leftmark}}}{
	}
	\ifthenelse{\equal{\hfstyle}{style10}}{
		\fancyhead[L]{\nouppercase{\leftmark}}}{
	}
	\ifthenelse{\equal{\hfstyle}{style11}}{ % Similar a 1
		\fancyhead[L]{\nouppercase{\leftmark}}}{
	}
	\ifthenelse{\equal{\hfstyle}{style14}}{ % Similar a 4
		\fancyhead[L]{\nouppercase{\leftmark}}}{
	}}{
}

% -----------------------------------------------------------------------------
% Estilo de títulos - reestablece estilos por el índice
% -----------------------------------------------------------------------------
\sectionfont{\color{\titlecolor} \fontsizetitle \styletitle \selectfont}
\subsectionfont{\color{\subtitlecolor} \fontsizesubtitle \stylesubtitle \selectfont}
\subsubsectionfont{\color{\subsubtitlecolor} \fontsizesubsubtitle \stylesubsubtitle \selectfont}
\titleformat{\subsubsubsection}{\color{\ssstitlecolor} \normalfont \fontsizessstitle \stylessstitle}{\thesubsubsubsection}{1em}{}
\titlespacing*{\subsubsubsection}{0pt}{3.25ex plus 1ex minus .2ex}{1.5ex plus .2ex}

% -----------------------------------------------------------------------------
% Numeración de objetos
% -----------------------------------------------------------------------------
\ifthenelse{\equal{\showsectioncaption}{none}}{
}{
\ifthenelse{\equal{\showsectioncaption}{sec}}{
	\counterwithin{equation}{section}
	\counterwithin{figure}{section}
	\counterwithin{lstlisting}{section}
	\counterwithin{table}{section}
}{
\ifthenelse{\equal{\showsectioncaption}{ssec}}{
	\counterwithin{equation}{subsection}
	\counterwithin{figure}{subsection}
	\counterwithin{lstlisting}{subsection}
	\counterwithin{table}{subsection}
}{
\ifthenelse{\equal{\showsectioncaption}{sssec}}{
	\counterwithin{equation}{subsubsection}
	\counterwithin{figure}{subsubsection}
	\counterwithin{lstlisting}{subsubsection}
	\counterwithin{table}{subsubsection}
}{
	\throwbadconfig{Valor configuracion incorrecto}{\showsectioncaption}{none,sec,ssec,sssec}
}}}}

% -----------------------------------------------------------------------------
% Se reestablecen números de página y secciones
% -----------------------------------------------------------------------------
\renewcommand{\thepage}{\arabic{page}}
\setcounter{page}{1}
\setcounter{section}{0}
\setcounter{footnote}{0}

% -----------------------------------------------------------------------------
% Muestra los números de línea
% -----------------------------------------------------------------------------
\ifthenelse{\equal{\showlinenumbers}{true}}{
	\linenumbers}{
}

% END

% ======================= INICIO DEL DOCUMENTO =======================

% Template:     Informe/Reporte LaTeX
% Documento:    Archivo de ejemplo
% Versión:      6.1.1 (05/12/2018)
% Codificación: UTF-8
%
% Autor: Pablo Pizarro R. @ppizarror
%        Facultad de Ciencias Físicas y Matemáticas
%        Universidad de Chile
%        pablo.pizarro@ing.uchile.cl, ppizarror.com
%
% Manual template: [https://latex.ppizarror.com/Template-Informe/]
% Licencia MIT:    [https://opensource.org/licenses/MIT/]

% ------------------------------------------------------------------------------
% NUEVA SECCIÓN
% ------------------------------------------------------------------------------
% Las secciones se inician con \section, si se quiere una sección sin número se
% pueden usar las funciones \sectionanum (sección sin número) o la función
% \sectionanumnoi para crear el mismo título sin numerar y sin aparecer en el índice.
\section{Informes con \LaTeX}

	% SUB-SECCIÓN
	% Las sub-secciones se inician con \subsection, si se quiere una sub-sección sin
	% número se pueden usar las funciones \subsectionanum (nuevo subtítulo sin numeración)
	% o la función \subsectionanumnoi para crear el mismo subtítulo sin numerar y sin
	% aparecer en el índice
	\subsection{Una breve introducción}
		
		Este es un párrafo, puede contener múltiples \quotes{Expresiones} así como fórmulas o referencias\footnote{Las referencias se hacen utilizando la expresión \texttt{\textbackslash label}\{etiqueta\}.} a fórmulas como \eqref{eqn:identidad-imposible}. A continuación se muestra un ejemplo de inserción de imágenes o figuras (como la Figura \ref{img:testimage}) con el comando \href{https://latex.ppizarror.com/informe.html#hlp-imagen}{\texttt{\textbackslash insertimage}}:

		% Para insertar una imagen se puede usar la función \insertimage la cual
		% toma un primer parámetro opcional para definir una etiqueta (dentro de
		% los corchetes), luego toma la dirección de la imagen, sus parámetros
		% (en este caso se definió la escala de 0.15) y una leyenda opcional.
		\insertimage[\label{img:testimage}]{ejemplos/test-image.png}{scale=0.15}{Where are you? de \quotes{Internet}.}

		A continuación\footnote{Como se puede observar las funciones \texttt{\textbackslash insert...} añaden un párrafo automáticamente.} se muestra un ejemplo de inserción de ecuaciones simples con el comando \href{https://latex.ppizarror.com/informe.html#hlp-formulae}{\texttt{\textbackslash insertequation}}:

		% Se inserta una ecuación, el primer parámetro entre [] es opcional
		% (permite identificar con una etiqueta para poder referenciarlo después
		% con \ref), seguido de aquello se escribe la ecuación en modo bruto sin signos $.
		\insertequation[\label{eqn:identidad-imposible}]{\pow{a}{k}=\pow{b}{k}+\pow{c}{k} \quad \forall k>2}

		% Se añade párrafo de prueba. Notar que no se requiere añadir un salto
		% de línea después de insertar una ecuación.
		\lipsum[75]

		% Los párrafos se pueden añadir con \newp, esta función se hizo para
		% evitar errores y warnings por parte del compilador de LaTeX.
		\newp Este es un nuevo párrafo insertado con el comando \href{https://latex.ppizarror.com/informe.html#hlp-parrafo}{\texttt{\textbackslash newp}}. Si no te gustan los comandos \texttt{\textbackslash newp}, \texttt{\textbackslash newpar} o \texttt{\textbackslash newparnl} simplemente puedes usar los salto de línea convencionales acompañado de \texttt{\textbackslash par}.

	% SUB-SECCIÓN
	\subsection{Añadiendo tablas}

		También puedes usar tablas, ¡Crearlas es muy fácil!. Puedes usar el plugin \href{https://www.ctan.org/tex-archive/support/excel2latex/}{Excel2Latex} \cite{ref2} de Excel para convertir las tablas a \LaTeX\xspace o bien utilizar el \quotes{creador de tablas online} \cite{ref3}.

		% Tabla generada con el plugin Excel2Latex
		\begin{table}[htbp]
			\centering
			\caption{Ejemplo de tablas.}
			\begin{tabular}{ccc}
				\hline
				\textbf{Columna 1} & \textbf{Columna 2} & \textbf{Columna 3} \bigstrut\\
				\hline
				$\omega$ & $\nu$ & $\delta$ \bigstrut[t]\\
				$\beta$ & $\gamma$ & $\epsilon$ \\
				$\Phi$ & $\Theta$ & $\varSigma$ \bigstrut[b]\\
				\hline
			\end{tabular}
			\label{tab:tabla-1}
		\end{table}


% ------------------------------------------------------------------------------
% NUEVA SECCIÓN
% ------------------------------------------------------------------------------
\newpage
\section{Aquí un nuevo tema}

	% SUB-SECCIÓN
	\subsection{Haciendo informes como un profesional}

		% Se inserta una imagen flotante en la izquierda del documento con
		% \insertimageleft, al igual que las demás funciones, el primer parámetro
		% es opcional, luego viene la ubicación de la imagen, seguido de la escala
		% (un 30% del ancho de página) y por último su leyenda. Para insertar una
		% imagen flotante en la derecha se utiliza \insertimageright usando los
		% mismos parámetros.
		\insertimageleft[\label{img:imagen-izquierda}]{ejemplos/test-image-wrap}{0.3}{Apolo flotando a la izquierda.}

		\lipsum[1]

		% Párrafos de ejemplo
		\newp \lipsum[115]
		\newp \lipsum[2]

		% Agrega una ecuación con leyenda
		\insertequationcaptioned[\label{eqn:formulasinsentido}]{\int_{a}^{b} f(x) \dd{x} = \fracnpartial{f(x)}{x}{\eta} \cdotp \textstyle \sum_{x=a}^{b} f(x)\cancelto{1+\frac{\epsilon}{k}}{(1+\Delta x)}}{Ecuación sin sentido.}

		% Aquí no es necesario usar \newp dado que todas las funciones \insert... añaden un párrafo nuevo por defecto
		\lipsum[115]

		\newp \lipsum[4]

	% Inserta un subtítulo sin número
	\subsection{Otros párrafos más normales}

		% Párrafos
		\lipsum[7]
		\newp \lipsum[2]

		% Se inserta una ecuación larga con el entorno gathered (1 solo número de ecuación)
		\insertgathered[\label{eqn:eqn-larga}]{
			\lpow{\Lambda}{f} = \frac{L\cdot f}{W} \cdot \frac{\pow{\lpow{Q}{e}}{2}}{8 \pow{\pi}{2} \pow{W}{4} g} + \sum_{i=1}^{l} \frac{f \cdot \big( M - d\big)}{l \cdot W} \cdot \frac{\pow{\big(\lpow{Q}{e}- i\cdot Q\big)}{2}}{8 \pow{\pi}{2} \pow{W}{4} g}\\
			Q_e = 2.5Q \cdot \int_{0}^{e} V(x) \dd{x} + \aasin{ \bigg(1+\frac{1}{1-e}\bigg) }
		}

		% Nuevo párrafo
		\lipsum[4]

		% Se inserta un multicols, con esto se pueden escribir párrafos en varias columnas
		\begin{multicols}{2}

			% Párrafo 1
			\lipsum[4]

			% Ecuación encerrada en una caja
			\insertequation[]{ \boxed{f(x) = \fracdpartial{u}{t}} }

			% Párrafo 2 del multicols
			\lipsum[1]

		\end{multicols}

	% SUB-SECCIÓN
	\subsection{Ejemplos de inserción de código fuente}

		% A continuación se crea una función auxiliar, esta es una herramienta
		% extremadamente importante y muy útil. Esta función de ejemplo toma dos
		% parámetros, uno es el lenguaje del código fuente, el segundo el
		% identificador en el manual.
		\newcommand{\insertsrcmanual}[2]{\href{https://latex.ppizarror.com/informe.html\#hlp-srccode\&srctype=#1}{#2}}

		El template permite la inserción de los siguientes lenguajes de programación de forma nativa: \insertsrcmanual{bash}{bash}, \insertsrcmanual{c}{C}, \insertsrcmanual{csharp}{C\#}, \insertsrcmanual{cpp}{C++}, \insertsrcmanual{cuda}{cuda}, \insertsrcmanual{docker}{DOCKER}, \insertsrcmanual{html5}{HTML5}, \insertsrcmanual{java}{Java}, \insertsrcmanual{js}{Javascript}, \insertsrcmanual{json}{JSON}, \insertsrcmanual{kotlin}{Kotlin}, \insertsrcmanual{latex}{LaTeX}, \insertsrcmanual{matlab}{Matlab}, \insertsrcmanual{perl}{Perl}, \insertsrcmanual{php}{PHP}, \insertsrcmanual{plaintext}{Texto plano}, \insertsrcmanual{pseudocode}{Pseudocódigo}, \insertsrcmanual{python}{Python}, \insertsrcmanual{ruby}{Ruby}, \insertsrcmanual{scala}{Scala}, \insertsrcmanual{sql}{SQL} y \insertsrcmanual{xml}{XML}. Para insertar un código fuente se debe usar el entorno \texttt{sourcecode}, o el entorno \texttt{sourcecodep} si es que se quiere utilizar parámetros adicionales. \newp

		A continuación se presenta un ejemplo de inserción de código fuente en Python (Código \ref{codigo-python}), Java (Código \ref{codigo-java}) y Matlab (Código \ref{codigo-matlab}):

% Se define el lenguaje del código. Cuidado: Los códigos en LaTeX son sensibles a las tabulaciones y espacios en blanco
\begin{sourcecode}[\label{codigo-python}]{python}{Ejemplo en Python.}
import numpy as np

def incmatrix(genl1, genl2):
	m = len(genl1)
	n = len(genl2)
	M = None # Comentario 1
	VT = np.zeros((n*m, 1), int) # Comentario 2
\end{sourcecode}

\begin{sourcecode}[\label{codigo-java}]{java}{Ejemplo en Java.}
import java.io.IOException;
import javax.servlet.*;

// Hola mundo
public class Hola extends GenericServlet {
	public void service(ServletRequest request, ServletResponse response)
	throws ServletException, IOException{
		response.setContentType("text/html");
		PrintWriter pw = response.getWriter();
		pw.println("Hola, mundo!");
		pw.close();
	}
}
\end{sourcecode}

\begin{sourcecode}[\label{codigo-matlab}]{matlab}{Ejemplo en Matlab.}
% Se crea gráfico
f = figure(1); hold on;
movegui(f, 'center');
xlabel('td/Tn'); ylabel('FAD=Umax/Uf0');
title('Espectro de pulso de desplazamiento');

for j = 1:length(BETA)
	fad = ones(1, NDATOS); % Arreglo para el FAD
	for i = 1:NDATOS
		[t, u_t, ~, ~] = main(BETA(j), r(i), M, K, F0, 0);
		fad(i) = max(abs(u_t)) / uf0;
	end
end
\end{sourcecode}

	% SUB-SECCIÓN
	\subsection{Añadir múltiples imágenes}

	El template ofrece el entorno \href{https://latex.ppizarror.com/informe.html#hlp-images}{\texttt{images}} que permite insertar múltiples imágenes de una manera muy sencilla \footnote{Desde la versión \texttt{6.0.0} no se da soporte a las funciones de inserción de imágenes múltiples \texttt{\textbackslash insertdoubleimage}, \texttt{\textbackslash insertdoubleeqimage},\texttt{\textbackslash inserttripleimage}, \texttt{\textbackslash inserttripleeqimage}, \texttt{\textbackslash insertquadimage}, \texttt{\textbackslash insertpentaimage} y \texttt{\textbackslash inserthexaimage}.}. Para crear imágenes múltiples se deben usar las siguientes instrucciones:

\begin{sourcecode}{latex}{}
\begin{images}[\label{imagenmultiple}]{Ejemplo de imagen múltiple.}
	\addimage{ejemplos/test-image}{width=6.5cm}{Ciudad.}
	\addimage{ejemplos/test-image-wrap}{width=5cm}{Apolo.}
	\addimagenewline
	\addimage{ejemplos/test-image}{width=12cm}{Ciudad más grande.}
\end{images}
\end{sourcecode}

	Obteniendo así:

	\begin{images}{Ejemplo de imagen múltiple.}
		\addimage{ejemplos/test-image}{width=6.5cm}{Ciudad.}
		\addimage{ejemplos/test-image-wrap}{width=5cm}{Apolo.}
		\addimagenewline
		\addimage{ejemplos/test-image}{width=12cm}{Ciudad más grande.}
	\end{images}


% ------------------------------------------------------------------------------
% NUEVA SECCIÓN
% ------------------------------------------------------------------------------
% Inserta una sección sin número
\sectionanum{Más ejemplos}

	% Inserta un subtítulo sin número
	\subsectionanum{Listas y Enumeraciones}

		Hacer listas enumeradas con \LaTeX\ es muy fácil con el template\footnote{También puedes revisar el manual de las enumeraciones en \url{http://www.texnia.com/archive/enumitem.pdf}.}, para ello debes usar el comando \texttt{\textbackslash begin\{enumerate\}}, cada elemento comienza por \texttt{\textbackslash item}, resultando así:

		\begin{enumerate}
			\item Grecia
			\item Abracadabra
			\item Manzanas
		\end{enumerate}

		También se puede cambiar el tipo de enumeración, se pueden usar letras, números romanos, entre otros. Esto se logra cambiando el \textbf{label} del objeto \texttt{enumerate}. A continuación se muestra un ejemplo usando letras con el estilo \texttt{\textbackslash alph} \footnote{Con \texttt{\textbackslash Alph} las letras aparecen en mayúscula.}, números romanos con \texttt{\textbackslash roman} \footnote{Con \texttt{\textbackslash Roman} los números romanos salen en mayúscula.} o números griegos con \texttt{\textbackslash greek}\footnote{Una característica propia del template, con \texttt{\textbackslash Greek} las letras griegas están escritas en mayúscula.}:

		\begin{multicols}{3}
			\begin{enumeratebf}[label=\alph*) ] % Fuente en negrita
				\item Peras
				\item Manzanas
				\item Naranjas
			\end{enumeratebf}

			\begin{enumerate}[label=\greek*) ]
				\item Matemáticas
				\item Lenguaje
				\item Filosofía
			\end{enumerate}

			\begin{enumerate}[label=\roman*) ]
				\item Rojo
				\item Café
				\item Morado
			\end{enumerate}
		\end{multicols}

		Para hacer listas sin numerar con \LaTeX\ hay que usar el comando \texttt{\textbackslash begin\{itemize\}}, cada elemento empieza por \texttt{\textbackslash item}, resultando:

		\begin{multicols}{3}
			\begin{itemize}[label={--}]
				\item Peras
				\item Manzanas
				\item Naranjas
			\end{itemize}

			\begin{enumerate}[label={*}]
				\item Rojo
				\item Café
				\item Morado
			\end{enumerate}

			\begin{itemize}
				\item Árboles
				\item Pasto
				\item Flores
			\end{itemize}
		\end{multicols}

	% Inserta un subtítulo sin número
	\subsectionanum{Otros}

		Recuerda revisar el manual de todas las funciones y configuraciones de este template visitando el siguiente link: \url{https://latex.ppizarror.com/Template-Informe/}. Si necesitas una ayuda muy específica sobre el template, o si tienes alguna sugerencia, me puedes enviar un correo a \insertemail{pablo.pizarro@ing.uchile.cl}.


% ------------------------------------------------------------------------------
% REFERENCIAS (ESTILO BIBTEX), revisar configuración \stylecitereferences
% ------------------------------------------------------------------------------
\newpage % Salto de página
\begin{references}
	\bibitem{ref1}
	Template Informe en \LaTeX.
	\textit{¡Revisa el manual online de este template!} \\
	\url{https://latex.ppizarror.com/Template-Informe/}

	\bibitem{ref2}
	Excel2Latex.
	\textit{Importa de forma sencilla tus tablas de Excel a \LaTeX.} \\
	\url{https://www.ctan.org/tex-archive/support/excel2latex/}

	\bibitem{ref3}
	Overleaf.
	\textit{Uno de los mejores editores online para \LaTeX, renovado con su versión 2.0.} \\
	\href{https://es.overleaf.com}{\texttt{https://es.overleaf.com/}}
\end{references}


% ------------------------------------------------------------------------------
% ANEXO
% ------------------------------------------------------------------------------
\newpage
\begin{anexo}
	\section{Cálculos realizados}

		\subsection{Metodología}
			\lipsum[1]

			% Imagen, se numerará automáticamente con la letra del anexo según
			% la configuración \appendixindepobjnum
			\insertimage[\label{img:anexo-2}]{ejemplos/test-image.png}{scale=0.2}{Imagen en anexo.}

		\subsection{Resultados}
			\lipsum[10]

			% Tablas
			\begin{table}[htbp]
				\centering
				\caption{Tabla de cálculo.}
				\begin{tabular}{ccc}
					\hline
					\textbf{Elemento} & $\epsilon_i$ & \boldmath{}\textbf{Valor}\unboldmath{} \bigstrut\\
					\hline
					A     & 10    & 3,14$\pi$ \bigstrut[t]\\
					B     & 20    & 6 \\
					C     & 30    & 7 \\
					\end{tabular}
				\label{tab:anexo-1}
			\end{table}

	\newpage
	\section{Más cálculos}

		% Párrafo
		\lipsum[1]

		% Nuevo párrafo con identación
		\newp \lipsum[4]

		% Tabla de encuestas
		\begin{table}[htbp]
			\centering
			\caption{Resultados encuesta.}
			\begin{tabular}{ccc}
				\hline
				\textbf{Herramienta} & \textbf{Nota} & \textbf{Recomendado} \bigstrut\\
				\hline
				\LaTeX & 100\% & Si $\checkmark$ \\
				Microsoft Word \textsuperscript{\textregistered} & 0\% & No $\frownie$\\
			\end{tabular}
			\label{tab:anexo-2}
		\end{table}

\end{anexo}
 % Ejemplo, se puede borrar

% FIN DEL DOCUMENTO
\end{document}